\subsection{Cosection-twisted reduced d-critical orientation}
\label{subsec:cosection-twisted-orientation}

The full reduced d-critical orientation entry of the compact Hall sector
on $K3\times E$ is the second entry of the pro-object.  Its completed
form consists of a Joyce d-critical structure, a vanishing-cycle
coefficient system, a Joyce--Upmeier--type orientation line, and a
quotient-descent datum for the reduced $E$-quotient.  The rows below
separate these inputs.  Before any Pfaffian wall or Weyl transport can
be discussed, the orientation line must descend through the
cosection-reduced moduli.

The compact Hall heart is determined by the chosen stability datum
$(\sigma,S)$ of Chapter~\ref{ch:k3xe-setup}; its K3-surface component
lies in the Bridgeland--Bayer--Macri lineage
\cite{BridgelandK3Stability,BayerMacriProjectivity}.  The first finite
orientation input is the PTVV \((-1)\)-shifted symplectic form on the
retained derived substacks.  The Joyce d-critical truncation, the K3
semiregularity cosection, cosection surjectivity, the reduced
self-Ext complex, the determinant line, and its square root are
separate rows.

\begin{proposition}[PTVV form on retained finite substacks]
\label{prop:ptvv-finite-substack-symplectic}
\textnormal{[proved for retained derived enhancements; d-critical truncation not constructed]}
Let \(R\) be a finite HN height and let
\(\mathfrak M^{ss}_{R,c}\) be one of the retained finite-type
semistable substacks carrying the quasi-smooth derived enhancement
\[
\mathfrak M^{\mathrm{der}}_{R,c}
\]
of Proposition~\ref{prop:retained-quasi-smooth-derived-enhancements}.
Then \(\mathfrak M^{\mathrm{der}}_{R,c}\) carries the restricted PTVV
\((-1)\)-shifted symplectic form
\[
\omega^{\mathrm{PTVV}}_{R,c}.
\]
Its cotangent amplitude is \([-1,0]\), and its symplectic restriction
defect is zero on every retained row.

This is the unreduced derived symplectic input.  It does not construct
the Joyce d-critical truncation, the K3 semiregularity cosection,
cosection surjectivity, \(\operatorname{RHom}_{\mathrm{red}}\), the
determinant line, a square-root orientation, quotient orientation, an
O2 wall atlas, or protected integration.
\end{proposition}

\begin{proof}
The retained derived-enhancement datum of
Definition~\ref{def:retained-quasi-smooth-derived-enhancement-datum}
assigns to each retained finite-type semistable substack a derived
Artin enhancement with cotangent amplitude \([-1,0]\).  Since
\[
K_{K3\times E}\simeq\mathcal O_{K3\times E},
\]
the ambient perfect-complex moduli problem is Calabi--Yau of dimension
three.  Pantev--To\"en--Vaqui\'e--Vezzosi
\cite[Theorem~0.3]{PTVV2013} therefore supplies a
\((-1)\)-shifted symplectic form on the derived moduli of perfect
complexes.  Proposition~\ref{prop:retained-quasi-smooth-derived-enhancements}
is exactly the retained finite-row restriction of this form to
\(\mathfrak M^{\mathrm{der}}_{R,c}\), with zero Tor-amplitude and
symplectic restriction defects.  The final paragraph lists the later
orientation rows not supplied by this PTVV input.
\end{proof}

The packet
\[
\texttt{certificates/orientation/ptvv\_finite\_substack\_symplectic}
\]
verifies
\(\texttt{PTVV\_FINITE\_SUBSTACK\_SYMPLECTIC\_VERIFIED}\): it imports
the retained derived-enhancement packet, checks the three retained
finite substacks, verifies shifted degree \(-1\), cotangent amplitude
\([-1,0]\), and zero symplectic defect, and records that d-critical
truncation, K3 cosection, reduced self-Ext, determinant, orientation
square root, quotient orientation, O2, and protected integration rows
remain open.

\begin{proposition}[Joyce d-critical truncation of retained finite substacks]
\label{prop:joyce-dcritical-truncation}
\textnormal{[proved for classical truncations; vanishing cycles not constructed]}
Let \(\mathfrak M^{\mathrm{der}}_{R,c}\) be a retained finite
quasi-smooth derived enhancement carrying the PTVV form
\(\omega^{\mathrm{PTVV}}_{R,c}\) of
Proposition~\ref{prop:ptvv-finite-substack-symplectic}.  Then its
classical truncation
\[
\mathfrak M^{\mathrm{cl}}_{R,c}
:=t_0\bigl(\mathfrak M^{\mathrm{der}}_{R,c}\bigr)
\]
carries the Joyce d-critical structure
\[
s^{\mathrm{Joyce}}_{R,c}\in
H^0\bigl(\mathfrak M^{\mathrm{cl}}_{R,c},\mathcal S^0_{\mathfrak M^{\mathrm{cl}}_{R,c}}\bigr)
\]
canonically associated to \(\omega^{\mathrm{PTVV}}_{R,c}\).  The
d-critical truncation defect is zero on every retained finite row.

This row constructs only the d-critical structure on the classical
truncation.  It does not construct the BBDJS vanishing-cycle complex,
the K3 semiregularity cosection, cosection surjectivity,
\(\operatorname{RHom}_{\mathrm{red}}\), determinant lines, square-root
orientations, quotient orientation, transition compatibility, or
protected integration.
\end{proposition}

\begin{proof}
The input of Proposition~\ref{prop:ptvv-finite-substack-symplectic} is
a retained quasi-smooth derived enhancement equipped with a
\((-1)\)-shifted symplectic form.  Brav--Bussi--Dupont--Joyce--Szendr{\H{o}}i
\cite[Theorem~6.9]{BBDJS2015} associate to the classical truncation
of a \((-1)\)-shifted symplectic derived scheme or Artin chart a
canonical Joyce d-critical structure, functorial on overlaps.  Applying
this theorem to each retained finite derived row gives
\(s^{\mathrm{Joyce}}_{R,c}\) on \(t_0(\mathfrak M^{\mathrm{der}}_{R,c})\).
The retained packet records zero defect for these three truncation rows.
The final paragraph separates the d-critical truncation from the later
coefficient, cosection, reduced-obstruction, orientation, quotient, and
integration rows.
\end{proof}

The packet
\[
\texttt{certificates/orientation/joyce\_dcritical\_truncation}
\]
verifies
\(\texttt{JOYCE\_DCRITICAL\_TRUNCATION\_VERIFIED}\): it imports the
PTVV finite-substack packet, applies the BBDJS d-critical truncation
theorem to the three retained finite derived substacks, and records zero
d-critical defect.  Its firewall rows record that vanishing cycles, K3
cosections, reduced self-Ext complexes, determinant lines, square-root
orientations, quotient orientation, transition compatibility, and
protected integration remain open.

The holomorphic symplectic two-form on the K3 fibre makes the
unreduced Donaldson--Thomas theory of \(K3\times E\) vanish
identically: it generates trivial first-order deformations that
obstruct the virtual class.  The standard remedy, as in reduced K3
enumerative geometry \cite{MaulikPandharipande}, is to quotient these
out by a semiregularity cosection.

\begin{proposition}[K3 semiregularity cosection on retained finite substacks]
\label{prop:k3-semiregularity-cosection}
\textnormal{[proved as an unreduced cosection; surjectivity not proved]}
Let \(X=S\times E\), where \(S\) is a K3 surface with holomorphic
two-form \(\sigma_S\), and put \(\sigma=p_S^*\sigma_S\).  For each
retained finite classical truncation
\(\mathfrak M^{\mathrm{cl}}_{R,c}\) of
Proposition~\ref{prop:joyce-dcritical-truncation}, the obstruction
sheaf \(\operatorname{Ob}_{R,c}\) carries an
\(\mathcal O_{\mathfrak M^{\mathrm{cl}}_{R,c}}\)-linear cosection
\[
\operatorname{CS}_{R,c}:
\operatorname{Ob}_{R,c}\longrightarrow
\mathcal O_{\mathfrak M^{\mathrm{cl}}_{R,c}} .
\]
At a point represented by a perfect complex \(F\) on \(X\), and for
\(\xi\in\operatorname{Ext}^2_X(F,F)\), the fibre map is
\[
\operatorname{CS}_{R,c}|_F(\xi)
=
\int_X
\sigma\wedge
\operatorname{Tr}\bigl((\xi\otimes 1_{\Omega_X^1})\circ
\operatorname{At}(F)\bigr),
\]
where \(\operatorname{At}(F)\in
\operatorname{Ext}^1_X(F,F\otimes\Omega_X^1)\) is the Atiyah class.
The construction defect is zero on every retained finite row.

This row constructs only the unreduced K3 semiregularity cosection.  It
does not prove cosection surjectivity on retained strata, filtered Hall
additivity, the reduced self-Ext complex, perfectness of the reduced
obstruction theory, determinant lines, square-root orientations,
quotient orientation, transition compatibility, vanishing cycles, or
protected integration.
\end{proposition}

\begin{proof}
The obstruction sheaf of the retained quasi-smooth derived enhancement
has fibre \(\operatorname{Ext}^2_X(F,F)\) on the chart represented by
the universal perfect complex \(F\).  The Atiyah class and trace are
functorial under restriction to etale charts.  Hence the assignment
displayed above glues from charts to an
\(\mathcal O\)-linear morphism of obstruction sheaves.  The wedge
product lies in
\[
H^3(X,\Omega_X^3)=H^3(X,K_X)\simeq \C
\]
because \(\sigma\in H^0(X,\Omega_X^2)\).  This is the
Buchweitz--Flenner--Bloch semiregularity map in the \(K3\times E\)
form used by Maulik--Toda
\cite[Lemma~2.6, Theorem~2.7]{MaulikTodaGV} and by the Kiem--Li
cosection formalism \cite{KiemLi2013}.  The packet checks the same
formula on the three retained truncation rows.  No nonvanishing or
surjectivity assertion enters this construction.
\end{proof}

The packet
\[
\texttt{certificates/orientation/k3\_semiregularity\_cosection}
\]
verifies
\(\texttt{K3\_SEMIREGULARITY\_COSECTION\_VERIFIED}\): it imports the
Joyce d-critical truncation packet, writes the semiregularity formula
above on the three retained finite rows, and records zero construction
defect.  Its firewall rows record that surjectivity, filtered Hall
additivity, reduced self-Ext complexes, determinant lines, square-root
orientations, quotient orientation, transition compatibility, vanishing
cycles, and protected integration remain open.

\begin{proposition}[Cosection surjectivity criterion on retained K3-class strata]
\label{prop:k3-cosection-surjectivity-criterion}
\textnormal{[criterion proved; retained nonzero K3-class witnesses not supplied]}
Let \(\mathfrak M^{\mathrm{cl}}_{R,c}\) be a retained finite classical
truncation equipped with the cosection
\(\operatorname{CS}_{R,c}\) of
Proposition~\ref{prop:k3-semiregularity-cosection}.  Suppose the
retained row carries a K3-class witness
\[
\beta_{R,c}\in H_2(S,\mathbb Z),\qquad \beta_{R,c}\ne 0,
\]
recorded by the support or charge map for every geometric point of
\(\mathfrak M^{\mathrm{cl}}_{R,c}\), and suppose the row lies on the
reduced DT/PT branch on which the Maulik--Toda semiregularity theorem
applies.  Then
\[
\operatorname{CS}_{R,c}:
\operatorname{Ob}_{R,c}\twoheadrightarrow
\mathcal O_{\mathfrak M^{\mathrm{cl}}_{R,c}}
\]
is surjective, equivalently its cokernel has rank zero.

The current retained finite packets do not supply the witness
\(\beta_{R,c}\ne0\), a branch-identification row, or a cokernel-rank
row.  Therefore row \(274\) is not discharged here.
\end{proposition}

\begin{proof}
Surjectivity of a cosection of the obstruction sheaf is local on the
classical moduli stack.  At a geometric point \(F\), the fibre map is
the semiregularity map displayed in
Proposition~\ref{prop:k3-semiregularity-cosection}.  On the reduced
nonzero K3-class DT/PT branch, Maulik--Toda
\cite[Lemma~2.6, Theorem~2.7]{MaulikTodaGV} identify this map with the
standard K3 semiregularity cosection and prove surjectivity; the same
surjective branch is used in the reduced stable-pair theory of
Oberdieck \cite[Theorems~1 and~3]{OberdieckReducedPT}.  If the retained
row records a nonzero K3 class \(\beta_{R,c}\) at every geometric point
and identifies the row with that branch, these pointwise surjections
glue to a surjection of sheaves.  The last assertion is a table-level
audit: the retained class, charge, and D0-HN semiregularity tables have
no row recording \(\beta_{R,c}\ne0\) and no row recording cokernel rank
zero.
\end{proof}

The packet
\[
\texttt{certificates/orientation/k3\_cosection\_surjectivity}
\]
verifies
\(\texttt{K3\_COSECTION\_SURJECTIVITY\_OBSTRUCTION\_VERIFIED}\): it
records the criterion above, checks that the current retained finite
class tables lack a K3-class witness column or row, checks that the D0
semiregularity table is empty, and keeps the surjectivity and cokernel
tables empty.  This status is not a proof of surjectivity; it is the
fail-closed ledger for row \(274\).

After a future retained-stratum packet supplies the nonzero K3-class
witness and the cokernel-rank-zero row, the later additivity row checks
the cosection on extension strata.  On
the extension stack $\mathfrak E_{A,B}$ of short exact sequences
$0\to A\to C\to B\to 0$ — with projections $p_1,p_2$ to the substacks
$\mathfrak M_A,\mathfrak M_B$ of sub- and quotient-object and $q$ to
$\mathfrak M_C$, the correspondence set up in
\S\ref{subsec:hall-correspondences} — the later additivity row is the
filtered Hall identity
\[
q^*\operatorname{CS}_C
=p_1^*\operatorname{CS}_A+p_2^*\operatorname{CS}_B,
\]
the compatibility hypothesis in Kiem--Li cosection localization
\cite{JoyceSong,JoyceUpmeier}.

\begin{definition}[Cosection-reduced self-Ext cone]
\label{def:rhom-red-self-ext-cone}
Let \(F\) be the universal object over a retained finite classical
truncation \(\mathfrak M^{\mathrm{cl}}_{R,c}\).  Assume the row carries
the cosection \(\operatorname{CS}_{R,c}\) of
Proposition~\ref{prop:k3-semiregularity-cosection}.  Define the
cosection-reduced self-Ext cone by
\[
\operatorname{RHom}_{\mathrm{red}}(F,F)
:=
\operatorname{fib}\!\left(
\Theta_F:
\operatorname{RHom}_X(F,F)\longrightarrow
R\Gamma(X,\mathcal O_X)\oplus H^2(S,\mathcal O_S)[-1]
\right),
\]
where
\[
\Theta_F=(\operatorname{tr},\operatorname{CS}_F),
\]
and equivalently
\[
\operatorname{RHom}_{\mathrm{red}}(F,F)
=
\operatorname{Cone}(\Theta_F)[-1].
\]
The first component removes the scalar trace direction; the second is
the K3 semiregularity cosection direction.  If the surjectivity witness
of Proposition~\ref{prop:k3-cosection-surjectivity-criterion} is also
supplied, this cone is the complex entering the Kiem--Li reduced
obstruction theory on the retained branch.  This definition does not
prove surjectivity, does not prove that the cone is perfect, does not
identify its determinant line, does not construct a square root, and does
not supply quotient orientation, transition compatibility, vanishing
cycles, or protected integration.
\end{definition}

The packet
\[
\texttt{certificates/orientation/rhomred\_definition}
\]
verifies \(\texttt{RHOMRED\_DEFINITION\_VERIFIED}\): it imports the
semiregularity-cosection packet and the surjectivity-obstruction packet,
records the cone formula above for the three retained finite rows, and
records zero definition defect.  Its firewall rows record that
surjectivity, determinant lines, square-root orientations,
quotient orientation, transition compatibility, vanishing cycles, and
protected integration are not supplied by the cone-definition row.

\begin{proposition}[Perfectness of the cosection-reduced self-Ext cone]
\label{prop:rhomred-perfectness}
\textnormal{[proved for the cone complex; determinant line not yet defined]}
Let \(F=\mathcal U_{R,c}\) be the universal perfect complex on
\(X\times\mathfrak M^{\mathrm{rig}}_{R,c}\) supplied by
Proposition~\ref{prop:retained-universal-perfect-complexes}, and let
\(\Theta_F\) be the trace--cosection morphism of
Definition~\ref{def:rhom-red-self-ext-cone}.  Then
\[
\operatorname{RHom}_{\mathrm{red}}(F,F)
=\operatorname{fib}(\Theta_F)
\]
is a perfect complex on the retained finite row.  The perfectness defect
is zero for each of the three retained substacks
\(Mss_{R0,s3},Mss_{R0,s2},Mss_{R0,s1}\).

This proves row \(276\) at the level of the reduced self-Ext cone.  It
does not prove cosection surjectivity, does not construct the determinant
line \(\det\operatorname{RHom}_{\mathrm{red}}(F,F)\), does not construct
a square-root orientation, and does not supply quotient orientation,
transition compatibility, vanishing cycles, or protected integration.
\end{proposition}

\begin{proof}
Let
\[
p:X\times\mathfrak M^{\mathrm{rig}}_{R,c}\longrightarrow
\mathfrak M^{\mathrm{rig}}_{R,c}
\]
be the projection.  Proposition~\ref{prop:retained-universal-perfect-complexes}
gives \(F=\mathcal U_{R,c}\in
\operatorname{Perf}(X\times\mathfrak M^{\mathrm{rig}}_{R,c})\), with
finite Tor amplitude and zero perfection defect.  Hence
\[
R\mathcal Hom(F,F)
\]
is perfect on \(X\times\mathfrak M^{\mathrm{rig}}_{R,c}\).  The morphism
\(p\) is smooth and proper because \(X=K3\times E\) is smooth
projective.  Therefore \(Rp_*R\mathcal Hom(F,F)\), written fibrewise as
\(\operatorname{RHom}_X(F,F)\), is perfect by the Stacks Project
\cite[Tag~0685, Lemma~37.61.13]{StacksProject}.  The target
\[
R\Gamma(X,\mathcal O_X)\oplus H^2(S,\mathcal O_S)[-1]
\]
is a finite-dimensional constant complex, hence perfect.  The stable
\(\infty\)-category of perfect complexes is closed under fibres and
cofibres.  Thus
\[
\operatorname{fib}\!\left(
\operatorname{RHom}_X(F,F)\xrightarrow{\Theta_F}
R\Gamma(X,\mathcal O_X)\oplus H^2(S,\mathcal O_S)[-1]\right)
\]
is perfect.  The row \(276\) packet checks this argument on the three
retained rows by importing the universal-perfect-complex packet and the
cone-definition packet.
\end{proof}

The packet
\[
\texttt{certificates/orientation/rhomred\_perfectness}
\]
verifies \(\texttt{RHOMRED\_PERFECTNESS\_VERIFIED}\): it imports the
retained universal-perfect-complex packet, the retained
derived-enhancement packet, and the cone-definition packet; records
perfect source and perfect target complexes; checks closure of perfect
complexes under fibres for the three retained rows; and records zero
perfectness defect.  Its firewall rows record that determinant lines,
square-root orientations, quotient orientation, transition compatibility,
vanishing cycles, and protected integration are not supplied by the
perfectness row.

\begin{definition}[Determinant line of the reduced self-Ext complex]
\label{def:rhomred-determinant-line}
Let \(F=\mathcal U_{R,c}\) be a retained universal perfect complex and
put
\[
K^{\mathrm{red}}_{F}
:=\operatorname{RHom}_{\mathrm{red}}(F,F).
\]
By Proposition~\ref{prop:rhomred-perfectness}, \(K^{\mathrm{red}}_{F}\)
is perfect.  The reduced determinant line is the Knudsen--Mumford
determinant
\[
\mathcal L^{\det}_{F}
:=\det K^{\mathrm{red}}_{F}
:=\operatorname{Det}_{\mathrm{KM}}\!\left(K^{\mathrm{red}}_{F}\right)
\in\operatorname{Pic}\!\left(\mathfrak M^{\mathrm{rig}}_{R,c}\right)
\]
\cite{KnudsenMumfordDet,DeligneDetCohomology}.  On a local perfect
representative \(K^\bullet\), this line is
\[
\det K^{\mathrm{red}}_{F}
=\bigotimes_i \det(K^i)^{(-1)^i},
\]
with the Knudsen--Mumford transition signs.  It is invariant under
quasi-isomorphism and functorial for pullback of the retained row.  The
determinant-defect rank is zero for
\[
Mss_{R0,s3},\qquad Mss_{R0,s2},\qquad Mss_{R0,s1}.
\]

This definition supplies row \(277\).  It does not choose a square root of
\(\mathcal L^{\det}_{F}\), does not compute \(w_2(\mathcal L^{\det}_F)\),
does not construct quotient orientation or Weyl transport, and does not
supply vanishing cycles or protected integration.
\end{definition}

The packet
\[
\texttt{certificates/orientation/rhomred\_determinant}
\]
verifies \(\texttt{RHOMRED\_DETERMINANT\_VERIFIED}\): it imports the
perfectness packet, applies the Knudsen--Mumford determinant functor to
the three retained perfect reduced self-Ext complexes, records the three
determinant lines, and records zero determinant defect.  Its firewall
rows record that square-root orientations, \(w_2\)-classes, quotient
orientation, transition compatibility, vanishing cycles, and protected
integration are not supplied by the determinant row.

\begin{proposition}[Square-root criterion for the reduced determinant line]
\label{prop:rhomred-square-root-criterion}
\textnormal{[criterion proved; square-root rows not supplied]}
Let \(\mathcal L^{\det}_{F}\) be the determinant line of
Definition~\ref{def:rhomred-determinant-line}.  A retained row
constructs the reduced orientation square root precisely when it supplies
a line
\[
o_F\in\operatorname{Pic}\!\left(\mathfrak M^{\mathrm{rig}}_{R,c}\right)
\]
and an isomorphism
\[
\epsilon_F:o_F^{\otimes 2}\xrightarrow{\sim}\mathcal L^{\det}_{F}.
\]
Equivalently, \(\mathcal L^{\det}_{F}\) lies in the essential image of
the squaring morphism
\[
[2]:\operatorname{Pic}\!\left(\mathfrak M^{\mathrm{rig}}_{R,c}\right)
\longrightarrow
\operatorname{Pic}\!\left(\mathfrak M^{\mathrm{rig}}_{R,c}\right).
\]
When it exists, the set of square roots is a torsor under the
two-torsion subgroup
\(\operatorname{Pic}(\mathfrak M^{\mathrm{rig}}_{R,c})[2]\).

The current retained orientation packet supplies no line \(o_F\), no
isomorphism \(\epsilon_F\), and no zero square-defect row for
\(Mss_{R0,s3},Mss_{R0,s2},Mss_{R0,s1}\).  Therefore row \(278\) is not
discharged by the present finite tables.  The class
\[
w_2(\mathcal L^{\det}_{F})
\]
is the mod-\(2\) obstruction entering row \(279\) and the quotient
orientation gerbes; its vanishing is necessary for a topological square
root, but the algebraic square root itself is the Picard-groupoid datum
\((o_F,\epsilon_F)\).
\end{proposition}

\begin{proof}
The assertion is the definition of a square root of a line bundle in the
Picard groupoid.  Passing to isomorphism classes gives the squaring map
on \(\operatorname{Pic}\); the fibre of this map over
\(\mathcal L^{\det}_{F}\) is either empty or a torsor under the kernel
of squaring, namely the two-torsion subgroup.  The finite-table
assertion is a direct inspection of
\(\texttt{certificates/orientation/k3e\_reduced\_orientation}\): its
\(\texttt{orientation\_lines.csv}\) table has no square-root rows and its
obstruction ledger records the missing square-root and orientation-class
vanishing rows.
\end{proof}

The packet
\[
\texttt{certificates/orientation/rhomred\_square\_root}
\]
verifies \(\texttt{RHOMRED\_SQUARE\_ROOT\_OBSTRUCTION\_VERIFIED}\): it
imports the determinant packet and the reduced-orientation obstruction
ledger, records the Picard-groupoid square-root criterion, checks that
the three determinant lines have no supplied square roots in the current
orientation table, and keeps \(w_2\), quotient orientation, transition
compatibility, vanishing cycles, and protected integration as open
obligations.

\begin{proposition}[Stiefel--Whitney formula for the reduced determinant line]
\label{prop:rhomred-w2-formula}
\textnormal{[formula proved; retained \(w_2\)-values not supplied]}
Let \(\mathcal L^{\det}_{F}\) be the determinant line of
Definition~\ref{def:rhomred-determinant-line}, regarded as a complex
line bundle on the retained analytic stack.  Its second
Stiefel--Whitney class is
\[
w_2\!\left((\mathcal L^{\det}_{F})_{\mathbb R}\right)
=c_1(\mathcal L^{\det}_{F})\bmod 2
\in H^2(\mathfrak M^{\mathrm{rig}}_{R,c};\mathbb F_2).
\]
If \(K^{\mathrm{red}}_{F}\) is represented by a perfect kernel
\(\mathcal K^{\mathrm{red}}_{F}\) on
\(X\times\mathfrak M^{\mathrm{rig}}_{R,c}\) and
\(\pi:X\times\mathfrak M^{\mathrm{rig}}_{R,c}\to
\mathfrak M^{\mathrm{rig}}_{R,c}\), then the determinant-of-cohomology
calculation required for row \(279\) is the degree-one part of
Grothendieck--Riemann--Roch:
\[
c_1(\mathcal L^{\det}_{F})
=
\left[
\pi_*\!\left(
\operatorname{ch}(\mathcal K^{\mathrm{red}}_{F})
\operatorname{td}(T_{\pi})
\right)
\right]_1 .
\]
Thus row \(279\) is discharged only by supplying, for each retained
substack, an integral \(c_1\)-row or an equivalent GRR expansion for
\(\mathcal K^{\mathrm{red}}_{F}\), together with its mod-\(2\)
reduction.

The present determinant and square-root packets supply none of these
rows for \(Mss_{R0,s3},Mss_{R0,s2},Mss_{R0,s1}\).  They define the
three determinant lines and record the absence of square roots, but they
do not compute \(c_1(\mathcal L^{\det}_{F})\), do not reduce it modulo
\(2\), and do not verify \(w_2=0\).
\end{proposition}

\begin{proof}
For a complex vector bundle \(E\), the Stiefel--Whitney classes of the
underlying real bundle satisfy
\(w_{2i}(E_{\mathbb R})=c_i(E)\bmod 2\) and
\(w_{2i+1}(E_{\mathbb R})=0\)
\cite[Theorem~14.10]{MilnorStasheffCharacteristicClasses}; the case of
a complex line gives the displayed formula.  The determinant line is the
Knudsen--Mumford determinant of the perfect complex
\cite{KnudsenMumfordDet,DeligneDetCohomology}; its first Chern class is
computed from determinant of cohomology by the displayed
Grothendieck--Riemann--Roch expression \cite{ToenRRRoot}.  Inspection of
\(\texttt{certificates/orientation/rhomred\_determinant}\) and
\(\texttt{certificates/orientation/rhomred\_square\_root}\) shows that
the retained rows have no \(c_1\)-payload, no GRR expansion row, no
mod-\(2\) reduction row, and no \(w_2\)-value row.
\end{proof}

The packet
\[
\texttt{certificates/orientation/rhomred\_w2}
\]
verifies \(\texttt{RHOMRED\_W2\_FORMULA\_OBSTRUCTION\_VERIFIED}\): it
imports the determinant and square-root packets, records the
\(w_2=c_1\bmod 2\) and determinant-of-cohomology formulae, checks the
three retained determinant lines, and keeps the actual \(c_1\), GRR,
mod-\(2\) reduction, \(w_2\)-vanishing, square-root, quotient
orientation, transition, vanishing-cycle, and protected-integration
rows open.

\begin{proposition}[Vanishing criterion for the reduced orientation gerbe]
\label{prop:alpha-red-vanishing-criterion}
\textnormal{[criterion proved; \(\alpha^{\mathrm{red}}=0\) rows not supplied]}
For a retained charge stratum \(\mathcal C\), define the reduced
orientation gerbe class
\[
\alpha^{\mathrm{red}}_{\mathcal C}
=w_2(\det K^{\mathrm{red}}_{\mathcal C})
+\operatorname{CS}_{\mathcal C}^{*}
w_2(\det\operatorname{coker}\operatorname{CS}_{\mathcal C})
\in H^2(\mathfrak M^{\mathrm{red}}_{\mathcal C};\mathbb F_2).
\]
A retained row proves \(\alpha^{\mathrm{red}}_{\mathcal C}=0\) precisely
when it supplies:
\[
u_{\mathcal C}=w_2(\det K^{\mathrm{red}}_{\mathcal C}),\qquad
v_{\mathcal C}=w_2(\det\operatorname{coker}\operatorname{CS}_{\mathcal C}),
\]
the pullback \(\operatorname{CS}_{\mathcal C}^{*}v_{\mathcal C}\), and
the equality
\[
u_{\mathcal C}+\operatorname{CS}_{\mathcal C}^{*}v_{\mathcal C}=0
\]
in \(H^2(\mathfrak M^{\mathrm{red}}_{\mathcal C};\mathbb F_2)\).
If the cosection-cokernel row proves
\(\operatorname{coker}\operatorname{CS}_{\mathcal C}=0\), then the
second summand is zero and row \(280\) reduces to the \(w_2\)-vanishing
row of Proposition~\ref{prop:rhomred-w2-formula}.  If the cokernel is
not proved zero, its determinant line and mod-\(2\) class are separate
inputs.

The present finite tables do not prove \(\alpha^{\mathrm{red}}=0\) for
\(Mss_{R0,s3},Mss_{R0,s2},Mss_{R0,s1}\).  They do not supply
retained-stratum \(w_2(\det K^{\mathrm{red}})\)-values, do not supply
cosection-cokernel rank-zero rows, do not supply
\(w_2(\det\operatorname{coker}\operatorname{CS})\)-rows, and do not
supply a zero-sum row for \(\alpha^{\mathrm{red}}\).  The null-homotopy
of this class is a stronger datum and belongs to row \(281\).
\end{proposition}

\begin{proof}
The displayed formula is the definition of the reduced degree-two
orientation gerbe in the quotient-orientation tower: the first summand
comes from the determinant line of the reduced self-Ext complex, and the
second is the cosection-cokernel correction pulled back along the
cosection map.  Since the coefficient group is \(\mathbb F_2\), vanishing
of the class is exactly the equality of the two summands with sum zero.
If the cokernel sheaf is zero, its determinant line is the trivial line,
so its second Stiefel--Whitney class is zero.  Inspection of
\(\texttt{certificates/orientation/rhomred\_w2}\),
\(\texttt{certificates/orientation/k3\_cosection\_surjectivity}\), and
\(\texttt{certificates/orientation/k3e\_reduced\_orientation}\) shows
that the determinant \(w_2\)-values, cokernel rank-zero rows, cokernel
\(w_2\)-rows, quotient-Borel \(\alpha^{\mathrm{red}}\)-rows, and
null-homotopy rows are absent.
\end{proof}

The packet
\[
\texttt{certificates/orientation/rhomred\_alpha\_red}
\]
verifies \(\texttt{RHOMRED\_ALPHA\_RED\_OBSTRUCTION\_VERIFIED}\): it
imports the row \(279\) \(w_2\) packet, the cosection-surjectivity
obstruction packet, and the reduced-orientation obstruction ledger;
records the formula and zero-sum criterion for
\(\alpha^{\mathrm{red}}\); checks the three retained strata; and keeps
the determinant \(w_2\), cokernel rank, cokernel \(w_2\), zero-sum,
null-homotopy, quotient orientation, transition, vanishing-cycle, and
protected-integration rows open.

\begin{proposition}[Null-homotopy criterion for the reduced gerbe]
\label{prop:alpha-red-nullhomotopy-criterion}
\textnormal{[criterion proved; null-homotopy rows not supplied]}
Fix a retained Cech--Borel cover of the reduced stratum
\(\mathfrak M^{\mathrm{red}}_{\mathcal C}\).  Let
\[
a^{\mathrm{red}}_{\mathcal C}\in
Z^2(\mathfrak M^{\mathrm{red}}_{\mathcal C};\mathbb F_2)
\]
be a cocycle representative of
\(\alpha^{\mathrm{red}}_{\mathcal C}\).  A row constructs a
null-homotopy of \(\alpha^{\mathrm{red}}_{\mathcal C}\) precisely when
it supplies a \(1\)-cochain
\[
h^{\mathrm{red}}_{\mathcal C}\in
C^1(\mathfrak M^{\mathrm{red}}_{\mathcal C};\mathbb F_2)
\]
and verifies the coboundary equation
\[
\delta h^{\mathrm{red}}_{\mathcal C}=a^{\mathrm{red}}_{\mathcal C}.
\]
Thus row \(281\) is stronger than row \(280\): the equality
\(\alpha^{\mathrm{red}}_{\mathcal C}=0\) proves only that such a
cochain exists after choosing a representative; row \(281\) records the
representative, the cochain, and the zero coboundary-defect row.  When
nonempty, the set of null-homotopies is a torsor under
\(H^1(\mathfrak M^{\mathrm{red}}_{\mathcal C};\mathbb F_2)\).

The present finite tables do not construct a null-homotopy for
\(Mss_{R0,s3},Mss_{R0,s2},Mss_{R0,s1}\).  They do not supply a cocycle
representative of \(\alpha^{\mathrm{red}}\), do not prove the class
zero, do not supply a \(1\)-cochain, and do not verify the coboundary
equation.  No quotient-orientation descent row follows from an empty
quotient-Borel table.
\end{proposition}

\begin{proof}
This is the definition of a null-homotopy in the cochain model for
degree-two \(\mathbb F_2\)-valued Cech--Borel cohomology.  A
degree-two cocycle is null-homotopic precisely when it is a coboundary;
the possible primitives differ by \(1\)-cocycles, and changing a
primitive by a \(1\)-coboundary gives the same trivialization.  Hence the
set of choices is a torsor under \(H^1\).  Inspection of
\(\texttt{certificates/orientation/rhomred\_alpha\_red}\) and
\(\texttt{certificates/orientation/k3e\_reduced\_orientation}\) shows
that no cocycle representative, no zero-class row, no \(1\)-cochain, and
no coboundary-defect-zero row is supplied.
\end{proof}

The packet
\[
\texttt{certificates/orientation/rhomred\_alpha\_red\_nullhomotopy}
\]
verifies
\(\texttt{RHOMRED\_ALPHA\_RED\_NULLHOMOTOPY\_OBSTRUCTION\_VERIFIED}\):
it imports the row \(280\) packet and the reduced-orientation
obstruction ledger, records the cochain criterion
\(\delta h^{\mathrm{red}}=a^{\mathrm{red}}\), checks the three retained
strata, and keeps the cocycle representative, zero-class, \(1\)-cochain,
coboundary-defect, quotient orientation, transition, vanishing-cycle,
and protected-integration rows open.

\begin{proposition}[Computation criterion for the connected free
\(E\)-Borel class]
\label{prop:free-e-borel-class-criterion}
\textnormal{[criterion proved; \(a_1,a_2\) rows not supplied]}
Let \(E^{\mathrm{top}}\simeq T^2\), and use the basis
\[
H^2(BE;\mathbb F_2)=\mathbb F_2u_1\oplus\mathbb F_2u_2
\]
of Lemma~\ref{lem:connected-e-descent}.  For a retained stratum
\(\mathcal C\) with free \(E\)-translation action, the connected-free
quotient-orientation class is the \(E_\infty^{2,0}\)-edge component of
the equivariant degree-two Borel class of the reduced determinant
complex:
\[
\alpha^{E,\mathrm{free}}_{\mathcal C}
=\operatorname{edge}_{2,0}
\bigl(w_2^E(\mathscr D_{\mathcal C})\bigr)
=a_1(\mathcal C)u_1+a_2(\mathcal C)u_2.
\]
A row computes \(\alpha^{E,\mathrm{free}}_{\mathcal C}\) precisely when
it supplies the equivariant Borel representative
\(w_2^E(\mathscr D_{\mathcal C})\), the spectral-sequence edge
calculation including any incoming \(d_2:E_2^{0,1}\to E_2^{2,0}\)
contribution, and the two coefficients
\[
a_1(\mathcal C),a_2(\mathcal C)\in\mathbb F_2.
\]
The vanishing equation \(a_1=a_2=0\) is row \(283\), and a chosen
null-trivialization of that zero class is row \(284\).

The present finite tables do not compute
\(\alpha^{E,\mathrm{free}}\) for
\(Mss_{R0,s3},Mss_{R0,s2},Mss_{R0,s1}\).  The retained
\(E\)-translation packet supplies free actions and quotient stacks, but
the reduced-orientation packet supplies no \(\texttt{free\_E}\)
quotient-Borel rows, no Borel spectral-sequence edge rows, and no
\(a_1,a_2\) coefficient rows.  The equality \(H^1(BE;\mathbb F_2)=0\)
rules out connected degree-one characters; it does not compute the
degree-two class above.
\end{proposition}

\begin{proof}
The cohomology of the connected torus action is
\(H^*(BE;\mathbb F_2)=\mathbb F_2[u_1,u_2]\), \(|u_i|=2\), by
Lemma~\ref{lem:connected-e-descent}.  Hence every connected-free
degree-two edge class has a unique expansion
\(a_1u_1+a_2u_2\).  The Borel fibration spectral sequence computes the
equivariant class from \(H^*_E(\mathcal C;\mathbb F_2)\), and the edge
component in bidegree \((2,0)\) is well-defined only after the relevant
incoming differential has been killed or included in the reported
coefficient row.  Inspection of
\(\texttt{certificates/moduli/retained\_e\_translation\_rigidifications}\)
shows that the free actions are supplied, while inspection of
\(\texttt{certificates/orientation/k3e\_reduced\_orientation}\) shows
that the quotient-Borel table has no \(\texttt{free\_E}\) row and no
coefficient or edge-reduction payload.
\end{proof}

The packet
\[
\texttt{certificates/orientation/rhomred\_free\_e\_borel}
\]
verifies \(\texttt{RHOMRED\_FREE\_E\_BOREL\_OBSTRUCTION\_VERIFIED}\):
it imports the retained \(E\)-translation rigidification packet and the
reduced-orientation obstruction ledger, records the basis and edge-class
criterion for
\(\alpha^{E,\mathrm{free}}=a_1u_1+a_2u_2\), checks the three retained
free \(E\)-actions, and keeps the equivariant Borel representative,
edge-reduction, \(a_1,a_2\) coefficient, vanishing, null-trivialization,
transition, vanishing-cycle, and protected-integration rows open.

\begin{proposition}[Vanishing criterion for the connected free
\(E\)-Borel class]
\label{prop:free-e-borel-vanishing-criterion}
\textnormal{[criterion proved; \(\alpha^{E,\mathrm{free}}=0\) rows not
supplied]}
Let \(\mathcal C\) be one of the retained free \(E\)-translation strata,
and suppose row \(282\) has supplied the connected-free edge class
\[
\alpha^{E,\mathrm{free}}_{\mathcal C}
=a_1(\mathcal C)u_1+a_2(\mathcal C)u_2
\in H^2(BE;\mathbb F_2)
=\mathbb F_2u_1\oplus\mathbb F_2u_2 .
\]
Then row \(283\) proves
\(\alpha^{E,\mathrm{free}}_{\mathcal C}=0\) precisely by supplying the
same coefficient row together with the two equalities
\[
a_1(\mathcal C)=0,\qquad a_2(\mathcal C)=0
\]
in \(\mathbb F_2\).  The vanishing is not implied by
\(H^1(BE;\mathbb F_2)=0\), by freeness of the \(E\)-action, by the
existence of the quotient stack, by a finite-stabilizer calculation, or
by any scalar trace or character value.  A null-trivialization of this
zero class is the separate row \(284\).

The present finite tables do not prove
\(\alpha^{E,\mathrm{free}}=0\).  The row \(282\) packet has no
\(a_1,a_2\) values; the reduced-orientation quotient-Borel table has no
\(\texttt{free\_E}\) row; and no zero-coefficient row is supplied for
\(Mss_{R0,s3},Mss_{R0,s2},Mss_{R0,s1}\).
\end{proposition}

\begin{proof}
The classes \(u_1,u_2\) form a basis of
\(H^2(BE;\mathbb F_2)\) by Lemma~\ref{lem:connected-e-descent}.
Therefore the expansion of
\(\alpha^{E,\mathrm{free}}_{\mathcal C}\) in
Proposition~\ref{prop:free-e-borel-class-criterion} is unique, and the
class is zero if and only if both coordinates are zero.  This is a
basis calculation after the Borel edge class has been computed; it is
not a replacement for that computation.  The packet
\(\texttt{certificates/orientation/rhomred\_free\_e\_borel}\) records
the absence of the coefficient row, and
\(\texttt{certificates/orientation/k3e\_reduced\_orientation}\) records
the empty quotient-Borel table, so the current finite evidence supplies
the criterion but not its hypotheses.
\end{proof}

The packet
\[
\texttt{certificates/orientation/rhomred\_free\_e\_vanishing}
\]
verifies
\(\texttt{RHOMRED\_FREE\_E\_VANISHING\_OBSTRUCTION\_VERIFIED}\): it
imports the row \(282\) packet and the reduced-orientation obstruction
ledger, records the basis-zero criterion
\(\alpha^{E,\mathrm{free}}=0\Longleftrightarrow a_1=a_2=0\), checks the
three retained free \(E\)-actions, and keeps the \(a_1,a_2\)
coefficient, zero-coefficient, null-trivialization, transition,
vanishing-cycle, and protected-integration rows open.

\begin{proposition}[Null-trivialization criterion for the connected
free \(E\)-Borel class]
\label{prop:free-e-borel-nulltrivialization-criterion}
\textnormal{[criterion proved; null-trivialization rows not supplied]}
Fix a cochain model for the connected \(BE\)-edge component used in
Proposition~\ref{prop:free-e-borel-class-criterion}.  Let
\[
a^{E,\mathrm{free}}_{\mathcal C}\in Z^2(BE;\mathbb F_2)
\]
be a cocycle representative of the connected-free edge class
\(\alpha^{E,\mathrm{free}}_{\mathcal C}\).  Row \(284\) constructs a
null-trivialization of
\(\alpha^{E,\mathrm{free}}_{\mathcal C}\) precisely when it supplies a
\(1\)-cochain
\[
h^{E,\mathrm{free}}_{\mathcal C}\in C^1(BE;\mathbb F_2)
\]
and verifies the coboundary equation
\[
\delta h^{E,\mathrm{free}}_{\mathcal C}
=a^{E,\mathrm{free}}_{\mathcal C}.
\]
Thus row \(284\) is stronger than row \(283\): the vanishing
\(\alpha^{E,\mathrm{free}}_{\mathcal C}=0\) supplies the cohomology
condition, while row \(284\) records a representative, a primitive, and
a defect-zero coboundary calculation.  Since
\(H^1(BE;\mathbb F_2)=0\), the connected \(BE\)-edge homotopy class of
such a trivialization has no \(H^1\)-torsor ambiguity after the cochain
model is fixed; this uniqueness does not replace the cochain-level
construction.

The present finite tables do not construct this null-trivialization.
They do not supply a cocycle representative of
\(\alpha^{E,\mathrm{free}}\), do not prove the class zero, do not
supply a \(1\)-cochain, and do not verify a defect-zero coboundary
equation for \(Mss_{R0,s3},Mss_{R0,s2},Mss_{R0,s1}\).
\end{proposition}

\begin{proof}
A null-trivialization of a degree-two \(\mathbb F_2\)-valued Borel
class is, by definition, a primitive for a chosen degree-two cocycle
representative.  Hence the datum is exactly a cochain \(h\) with
\(\delta h=a\), plus the verification that the coboundary defect is
zero.  The equality \(H^1(BE;\mathbb F_2)=0\) from
Lemma~\ref{lem:connected-e-descent} says that two connected-edge
trivializations of the same cocycle are homotopic after fixing the
cochain model; it does not give the cocycle representative, the
primitive, or the coboundary calculation.  Inspection of
\(\texttt{certificates/orientation/rhomred\_free\_e\_vanishing}\) and
\(\texttt{certificates/orientation/k3e\_reduced\_orientation}\) shows
that none of these rows is supplied.
\end{proof}

The packet
\[
\texttt{certificates/orientation/rhomred\_free\_e\_nulltrivialization}
\]
verifies
\(\texttt{RHOMRED\_FREE\_E\_NULLTRIVIALIZATION\_OBSTRUCTION\_VERIFIED}\):
it imports the row \(283\) packet and the reduced-orientation
obstruction ledger, records the cochain criterion
\(\delta h^{E,\mathrm{free}}=a^{E,\mathrm{free}}\), checks the three
retained free \(E\)-actions, and keeps the cocycle representative,
zero-class, \(1\)-cochain, coboundary-defect, quotient orientation,
transition, vanishing-cycle, and protected-integration rows open.

\subsubsection*{The reduced determinant complex}

The following notation belongs to the later orientation rows.  After the
cosection, cone-definition, perfectness, and determinant rows have been
supplied, for an object $\mathcal C$ in the chosen finite stage
write
\[
K^{\mathrm{red}}_{\mathcal C}
=\operatorname{RHom}_{\mathrm{red}}(\mathcal C,\mathcal C),
\]
the cosection-reduced self-Ext complex of
Definition~\ref{def:rhom-red-self-ext-cone}; it is perfect by
Proposition~\ref{prop:rhomred-perfectness}.  Its determinant line is
\[
\mathcal L^{\det}_{\mathcal C}
=\det K^{\mathrm{red}}_{\mathcal C}
\]
by Definition~\ref{def:rhomred-determinant-line}.  The later orientation
row, when supplied, equips the reduced orientation line with a chosen
square root,
\[
o_{\mathcal C}\in\operatorname{Pic}(\mathfrak M^{\mathrm{or}}_X),
\qquad
o_{\mathcal C}^{\otimes 2}\simeq\det K^{\mathrm{red}}_{\mathcal C},
\]
in the sense of \cite{JoyceUpmeier}, and is multiplicative under the
reduced extension correspondences:
$o_{\mathcal C\oplus\mathcal C'}\simeq
o_{\mathcal C}\otimes o_{\mathcal C'}$ compatibly with the Hall
additivity above.  This is the data on which Thom--Sebastiani
transport is supplied for the vanishing-cycle coefficient sheaves.

\subsubsection*{(O1) Strong reduced orientation through the
quotient-orientation tower}

The compact Hall product on $X=K3\times E$ exists only after the
reduced $E$-quotient.  The orientation line must descend through this
quotient, and descent is detected by a tower of mod-$2$
Borel-equivariant gerbes.  For a charge stratum $\mathcal C$:

\begin{enumerate}
\item the reduced gerbe class
\[
\alpha^{\mathrm{red}}_{\mathcal C}
=w_2(\det K^{\mathrm{red}}_{\mathcal C})
+\operatorname{CS}_{\mathcal C}^{*}w_2(\det\operatorname{coker}\operatorname{CS}_{\mathcal C})
\in H^2(\mathfrak M^{\mathrm{red}}_{\mathcal C};\mathbb F_2);
\]
\item the free $E$-class
\[
\alpha^{E,\mathrm{free}}_{\mathcal C}
=a_1(\mathcal C)u_1+a_2(\mathcal C)u_2
\in H^2(BE;\mathbb F_2);
\]
\item for each finite-stabilizer stratum $S$ with stabilizer
$H\subset E_{\mathrm{tors}}$, the equivariant class
$\widetilde\beta^{H}_{\mathcal C,S}\in H^2_H(S;\mathbb F_2)$, whose
point-stabilizer restriction $\beta^H_{\mathcal C,S}\in H^2(BH;\mathbb F_2)$
is read only after the mixed Borel terms, the stabilizer action on
$H^\bullet(S;\mathbb F_2)$, and the spectral-sequence differentials
have vanished.
\end{enumerate}

\begin{definition}[(O1) Strong multiplicative reduced orientation]
\label{def:O1-strong-reduced-orientation}
The (O1) datum at finite HN height $R$ consists of:
the reduced orientation line $o_{\mathcal C}$ with chosen square root;
its multiplicativity under reduced extension correspondences with
chosen Thom--Sebastiani transport;
null-trivializations of $\alpha^{\mathrm{red}}_{\mathcal C}$,
$\alpha^{E,\mathrm{free}}_{\mathcal C}$,
$\widetilde\beta^{H}_{\mathcal C,S}$ and
$\beta^{H}_{\mathcal C,S}$ on every charge stratum
used by the Hall correspondences;
and a choice of zero $H^1(BH;\mathbb F_2)$-linearization character
$\lambda^{H}_{\mathcal C,S}=0$ for each finite stabilizer.

The full-level specialization at $H=E[2]$ is the Klein-four condition
\[
\beta^{E,2}_{\mathcal C,S}
=b_{20}x_1^2+b_{11}x_1x_2+b_{02}x_2^2,
\quad
(b_{20},b_{11},b_{02})=(r_1,r_1+r_2+r_3,r_2),
\]
\[
\operatorname{res}_{\langle e_i\rangle}\beta^{E,2}_{\mathcal C,S}
=r_i\tau_i^2,
\]
with vanishing $r_1=r_2=r_3=0$ on every retained two-torsion
stratum.  For full-level $H=E[N]$ with $2^a\parallel N$, $a\ge2$, the
two-primary residual takes value in
$H^2(B(\Z/2^a)^2;\mathbb F_2)=\mathbb F_2\langle y_1,x_1x_2,y_2\rangle$
and contains the mixed coefficient $A_{12}^{(N)}$ \emph{not} detected
by cyclic restrictions.
\end{definition}

(O1) is therefore not a single Picard-line assertion: it is the
ordinary square root, the null-trivialization of the reduced and
equivariant degree-two gerbes, and the choice of zero degree-one
linearization characters for the actual finite stabilizers.  At
finite HN height this is the cocycle system of
\cite{JoyceUpmeier} restricted to $K3\times E$.

After finite trivialisations, the \((O1)\) datum is recorded by the
orientation proof tables
\[
\mathrm{orientation\_lines},\quad
\mathrm{ts\_multiplicativity},\quad
\mathrm{quotient\_borel},\quad
\mathrm{finite\_stabilizers}.
\]
The first table records the determinant square root and the vanishing
orientation class.  The second records Thom--Sebastiani
multiplicativity and flag-pentagon defects.  The third records the
four quotient Cech--Borel classes
\(\alpha^{\mathrm{red}}\), \(\alpha^{E,\mathrm{free}}\),
\(\widetilde\beta^H\), and \(\lambda^H\), with chosen
null-trivialisations and edge-reduction status.  The fourth records
the finite-stabilizer edge calculation, including the Klein-four
coefficients, the two-primary coefficients
\((A_1^{(H)},A_{12}^{(H)},A_2^{(H)})\), and the two degree-one
linearization coefficients.  All ranks and coefficients in these rows
must vanish.  A row using only cyclic restrictions is not an
orientation row.

\begin{proposition}[Computation criterion for the finite-stabilizer
Borel class]
\label{prop:finite-stabilizer-beta-class-criterion}
\textnormal{[criterion proved; \(\beta^H_{\mathcal C,S}\) rows not
supplied]}
Let \(S\) be a retained stratum with residual finite stabilizer
\(H\subset E_{\mathrm{tors}}\).  Row \(285\) computes the
finite-stabilizer class
\[
\beta^H_{\mathcal C,S}\in H^2(BH;\mathbb F_2)
\]
precisely when it supplies:
\begin{enumerate}
\item a finite-stabilizer stratum row identifying \(H\) and its
two-primary type;
\item a Cech--Borel representative
\[
\widetilde\beta^H_{\mathcal C,S}
=w_2^H(\mathscr D_{\mathcal C}|_S)
\in H^2_H(S;\mathbb F_2);
\]
\item an edge-reduction row killing or accounting for the ordinary
term in \(H^2(S;\mathbb F_2)^H\), the mixed term in
\(H^1(BH;H^1(S;\mathbb F_2))\), and the incoming Borel
spectral-sequence differentials;
\item the classifying-space coordinates of the resulting edge class.
\end{enumerate}
After these rows are supplied, the computation is the following group
cohomology calculation.  If \(H\) has odd order, and also for
\(H=1\), positive-degree \(H^*(BH;\mathbb F_2)\) vanishes.  If
\(H\simeq E[2]\), then
\[
\beta^H_{\mathcal C,S}
=b_{20}x_1^2+b_{11}x_1x_2+b_{02}x_2^2.
\]
If the two-primary part is \((\mathbb Z/2^a)^2\), \(a\ge2\), then
\[
\beta^H_{\mathcal C,S}
=A_1^{(H)}y_1+A_{12}^{(H)}x_1x_2+A_2^{(H)}y_2.
\]
The coefficient \(A_{12}^{(H)}\) is a rank-two stabilizer coefficient;
it is not read from cyclic order-two restrictions.  Row \(286\), not
row \(285\), is the assertion that the computed class vanishes for
every retained \(H\).  Rows \(288\)--\(289\) concern the independent
degree-one linearization character.

The present finite tables do not compute
\(\beta^H_{\mathcal C,S}\).  The retained rigidification-inertia packet
records finite residual groups \(H_{R0,s3},H_{R0,s2},H_{R0,s1}\) of
order \(1\), but it does not supply finite closed inertia
stratifications, Cech--Borel representatives, edge-reduction rows, or
finite-stabilizer orientation rows.  The reduced-orientation packet has
an empty \(\texttt{finite\_stabilizers}\) table and an empty
\(\texttt{quotient\_borel}\) table.
\end{proposition}

\begin{proof}
The Borel spectral sequence for \(S\to EH\times_H S\to BH\) filters
\(H^2_H(S;\mathbb F_2)\) by the classifying-space edge
\(H^2(BH;\mathbb F_2)\), the mixed term
\(H^1(BH;H^1(S;\mathbb F_2))\), and the ordinary term
\(H^2(S;\mathbb F_2)^H\), with the relevant incoming differentials.
Only after the non-edge terms are killed or explicitly accounted for
does the restriction of \(\widetilde\beta^H\) define
\(\beta^H\in H^2(BH;\mathbb F_2)\).  Lemma~\ref{lem:finite-stabilizer-e-descent}
and Lemmas~\ref{lem:G-Klein-four}--\ref{lem:G-two-primary-stabilizer}
give the displayed coordinates: transfer kills odd stabilizers,
\((\mathbb Z/2)^2\) gives the quadratic basis
\(x_1^2,x_1x_2,x_2^2\), and \((\mathbb Z/2^a)^2\), \(a\ge2\), gives
the basis \(y_1,x_1x_2,y_2\).  Inspection of
\(\texttt{certificates/moduli/retained\_rigidification\_inertia}\) and
\(\texttt{certificates/orientation/k3e\_reduced\_orientation}\) shows
that only residual finite group input is present; the edge and
coefficient rows are absent.
\end{proof}

The packet
\[
\texttt{certificates/orientation/rhomred\_finite\_stabilizer\_beta}
\]
verifies
\(\texttt{RHOMRED\_FINITE\_STABILIZER\_BETA\_OBSTRUCTION\_VERIFIED}\):
it imports the residual-inertia packet and the reduced-orientation
obstruction ledger, records the odd, Klein-four, and two-primary
coordinate criteria for row \(285\), checks the three retained residual
groups, and keeps the finite closed inertia stratum, Cech--Borel
representative, edge-reduction, \(\beta^H\)-coefficient, vanishing,
null-trivialization, linearization, transition, vanishing-cycle, and
protected-integration rows open.

\begin{proposition}[Vanishing criterion for the finite-stabilizer
Borel class]
\label{prop:finite-stabilizer-beta-vanishing-criterion}
\textnormal{[criterion proved; \(\beta^H_{\mathcal C,S}=0\) rows not
supplied]}
Assume row \(285\) has computed
\[
\beta^H_{\mathcal C,S}\in H^2(BH;\mathbb F_2)
\]
for every retained finite-stabilizer stratum.  Row \(286\) proves
\[
\beta^H_{\mathcal C,S}=0
\]
for every retained \(H\) precisely when it verifies the corresponding
zero condition in the classifying-space basis for every retained row:
\[
H=1\ \text{or}\ |H|\ \text{odd}:\quad H^2(BH;\mathbb F_2)=0
\quad\text{after edge reduction,}
\]
\[
H\simeq E[2]:\quad b_{20}=b_{11}=b_{02}=0,
\]
and, for two-primary rank two stabilizers,
\[
H_{(2)}\simeq(\mathbb Z/2^a)^2,\ a\ge2:
\quad A_1^{(H)}=A_{12}^{(H)}=A_2^{(H)}=0.
\]
The condition is imposed on every retained finite-stabilizer row.  A
cyclic-only calculation does not prove the two-primary case, because it
does not see \(A_{12}^{(H)}\).  Row \(287\) is the subsequent choice of
compatible null-trivializations, and rows \(288\)--\(289\) concern the
independent \(H^1(BH;\mathbb F_2)\)-linearization character.

The present finite tables do not prove the displayed vanishing.  The
row \(285\) packet records residual finite groups of order \(1\), but
it records no finite-stabilizer orientation row, no edge-reduction row,
no \(\beta^H\)-value, and no zero row.  The group-cohomology fact
\(H^2(B1;\mathbb F_2)=0\) is a zero template for a supplied edge row;
it is not itself the missing finite-stabilizer orientation datum.
\end{proposition}

\begin{proof}
After the edge-reduction hypotheses of
Proposition~\ref{prop:finite-stabilizer-beta-class-criterion}, the
class is an element of \(H^2(BH;\mathbb F_2)\).  Transfer kills this
group for odd \(H\), and it is zero for \(H=1\).  Lemma~\ref{lem:G-Klein-four}
identifies the \(E[2]\) basis
\(x_1^2,x_1x_2,x_2^2\), so vanishing is equivalent to the three
quadratic coefficients being zero.  Lemma~\ref{lem:G-two-primary-stabilizer}
identifies the two-primary rank-two basis \(y_1,x_1x_2,y_2\), so
vanishing is equivalent to the three displayed coefficients being
zero.  The coefficient \(A_{12}^{(H)}\) is independent of cyclic
order-two restrictions.  Inspection of
\(\texttt{certificates/orientation/rhomred\_finite\_stabilizer\_beta}\)
and \(\texttt{certificates/orientation/k3e\_reduced\_orientation}\)
shows that the current finite evidence supplies the criterion but not
its hypotheses.
\end{proof}

The packet
\[
\texttt{certificates/orientation/rhomred\_finite\_stabilizer\_beta\_vanishing}
\]
verifies
\(\texttt{RHOMRED\_FINITE\_STABILIZER\_BETA\_VANISHING\_OBSTRUCTION\_VERIFIED}\):
it imports the row \(285\) packet and the reduced-orientation
obstruction ledger, records the odd/trivial, Klein-four, and
two-primary zero criteria for row \(286\), checks the three retained
residual groups, and keeps the \(\beta^H\)-value, zero-coordinate,
null-trivialization, linearization, transition, vanishing-cycle, and
protected-integration rows open.

\begin{proposition}[Compatible null-trivialization criterion for
finite-stabilizer Borel classes]
\label{prop:finite-stabilizer-beta-nulltrivialization-criterion}
\textnormal{[criterion proved; compatible null-trivialization rows not
supplied]}
Assume rows \(285\)--\(286\) have supplied, for every retained
finite-stabilizer stratum \(S\) with residual stabilizer \(H\), a
Cech--Borel cocycle
\[
a^H_{\mathcal C,S}\in Z^2(BH;\mathbb F_2)
\]
representing \(\beta^H_{\mathcal C,S}\), and have verified
\(\beta^H_{\mathcal C,S}=0\).  Row \(287\) is the following stronger
cochain-level datum.  It must supply a \(1\)-cochain
\[
h^H_{\mathcal C,S}\in C^1(BH;\mathbb F_2)
\]
with
\[
\delta h^H_{\mathcal C,S}=a^H_{\mathcal C,S}.
\]
For every retained subgroup inclusion
\(\iota\colon K\hookrightarrow H\) it must also verify
\[
\operatorname{res}_{\iota}a^H_{\mathcal C,S}=a^K_{\mathcal C,S}
\]
and supply a comparison \(0\)-cochain
\[
q^{\iota}_{\mathcal C,S}\in C^0(BK;\mathbb F_2)
\]
such that
\[
\operatorname{res}_{\iota}h^H_{\mathcal C,S}
-h^K_{\mathcal C,S}
=\delta q^{\iota}_{\mathcal C,S}.
\]
For every chain \(L\xrightarrow{\jmath}K\xrightarrow{\iota}H\), the
chosen comparisons must have zero defect
\[
q^{\iota\circ\jmath}_{\mathcal C,S}
=q^{\jmath}_{\mathcal C,S}
+\operatorname{res}_{\jmath}q^{\iota}_{\mathcal C,S}
\]
in the chosen Cech--Borel cochain model.  This is the subgroup
restriction compatibility demanded in row \(287\).  It is not the
linearization character of rows \(288\)--\(289\), and it is not implied
by a scalar trace, by the residual group order, or by the abstract
vanishing of \(H^2(BH;\mathbb F_2)\) unless the representative
\(a^H_{\mathcal C,S}\), the cochain \(h^H_{\mathcal C,S}\), and all
restriction comparisons are supplied.

The present finite tables do not construct row \(287\).  They record
three residual groups of order \(1\), but they record no
finite-stabilizer orientation row, no Cech--Borel representative
\(a^H_{\mathcal C,S}\), no verified zero row from row \(286\), no
\(1\)-cochain \(h^H_{\mathcal C,S}\), no retained subgroup-inclusion
table, no comparison \(0\)-cochain, and no chain-coherence row.  Thus
row \(287\) remains an open cochain-level obstruction, even in the
trivial residual-group template.
\end{proposition}

\begin{proof}
A null-trivialization of the representative cocycle
\(a^H_{\mathcal C,S}\) is, by definition, a degree-one cochain whose
coboundary is that representative.  Therefore the cohomological
statement \(\beta^H_{\mathcal C,S}=0\) is only the existence statement;
row \(287\) asks for a chosen primitive.  Restriction along
\(\iota\colon K\hookrightarrow H\) is a cochain map, so
\(\delta(\operatorname{res}_{\iota}h^H-h^K)=0\) once
\(\operatorname{res}_{\iota}a^H=a^K\).  Compatibility requires this
closed \(1\)-cochain to be killed by a chosen \(0\)-cochain
\(q^\iota\), and functoriality over a two-step chain is exactly the
displayed defect-zero equation.  These are cochain equalities, not
statements about the final scalar shadow.  Inspection of
\(\texttt{certificates/orientation/rhomred\_finite\_stabilizer\_beta\_vanishing}\),
\(\texttt{certificates/orientation/k3e\_reduced\_orientation}\), and
\(\texttt{certificates/orientation/transition\_orientation\_preservation}\)
shows that the required cochains and subgroup-restriction rows are not
present.
\end{proof}

The packet
\[
\texttt{certificates/orientation/rhomred\_finite\_stabilizer\_beta\_nulltrivialization}
\]
verifies
\(\texttt{RHOMRED\_FINITE\_STABILIZER\_BETA\_NULLTRIVIALIZATION\_OBSTRUCTION\_VERIFIED}\):
it imports the row \(286\) packet and the reduced-orientation
obstruction ledger, records the cochain and subgroup-restriction
criteria for row \(287\), checks the three retained residual groups,
and keeps the \(\beta^H\)-cocycle, zero row, \(1\)-cochain,
coboundary-defect, subgroup-inclusion, comparison \(0\)-cochain,
chain-coherence, linearization, transition, vanishing-cycle, and
protected-integration rows open.

\begin{proposition}[Computation criterion for the finite-stabilizer
linearization character]
\label{prop:finite-stabilizer-linearization-character-criterion}
\textnormal{[criterion proved; linearization-character rows not
supplied]}
Fix a retained finite-stabilizer stratum \(S\) with residual stabilizer
\(H\).  After the row \(287\) null-trivialization has killed the
degree-two Borel gerbe, the remaining finite-stabilizer descent datum
is a character
\[
\lambda^H_{\mathcal C,S}\in H^1(BH;\mathbb F_2)
=\operatorname{Hom}(H,\mathbb F_2).
\]
Row \(288\) computes this class precisely when it supplies the
orientation-line action of \(H\) on the already trivialized Pfaffian
orientation line, an ordered generator basis for the two-primary part
of \(H\), and the sign values of this action on the chosen generators.
For \(H=1\), and more generally for odd \(H\), the target group
\(H^1(BH;\mathbb F_2)\) is zero after the retained stabilizer row and
orientation-line action have been supplied.  For
\(H\simeq E[2]\simeq(\mathbb Z/2)^2\), with basis \(x_1,x_2\) dual to
the chosen generators, the computation is
\[
\lambda^H_{\mathcal C,S}=\lambda_1x_1+\lambda_2x_2,
\qquad
\rho_1=\lambda_1,\quad \rho_2=\lambda_2,\quad
\rho_3=\lambda_1+\lambda_2
\]
on the three non-zero cyclic subgroups.  For
\(H_{(2)}\simeq(\mathbb Z/2^a)^2\), \(a\ge2\), it is
\[
\lambda^H_{\mathcal C,S}
=\lambda^{(H)}_1x_1+\lambda^{(H)}_2x_2
\in H^1(BH;\mathbb F_2),
\]
where \(x_i\) are the mod-\(2\) degree-one classes dual to the chosen
two-primary generators.  Row \(288\) is only this computation of
\(\lambda^H_{\mathcal C,S}\); row \(289\) is the separate assertion
that the computed character is zero.

The present finite tables do not compute row \(288\).  The determinant
anchor translation-weight packet records the class-level restriction
\(h\mapsto\chi_\eta h\), but it marks coherent stabilizer
linearization as not certified.  The reduced-orientation packet has no
finite-stabilizer rows, no quotient Borel rows, and no
orientation-line action rows.  The row \(287\) packet has no
null-trivializations and records \(\texttt{linearization\_computed}=
\texttt{false}\) on every retained residual group.  Thus the equality
\(H^1(B1;\mathbb F_2)=0\), cyclic quadratic restrictions, and
translation-weight formulas are templates or separate data; none is a
computed value of \(\lambda^H_{\mathcal C,S}\) for row \(288\).
\end{proposition}

\begin{proof}
Once the degree-two gerbe has been trivialized, two finite-stabilizer
linearizations of the same orientation line differ by a sign character
of \(H\).  This torsor is \(H^1(BH;\mathbb F_2)\), so a computation of
the residual class is exactly a computation of the sign character on
generators.  Transfer gives \(H^1(BH;\mathbb F_2)=0\) for odd \(H\),
and the trivial group has no positive cohomology.  Lemma~\ref{lem:G-Klein-four}
identifies the \(E[2]\) basis and the three cyclic restrictions
\(\rho_1,\rho_2,\rho_3\).  Lemma~\ref{lem:G-two-primary-stabilizer}
identifies the degree-one basis \(x_1,x_2\) for the two-primary rank
two case.  These are coordinate computations for a supplied action on
the orientation line.  Inspection of
\(\texttt{certificates/orientation/rhomred\_finite\_stabilizer\_beta\_nulltrivialization}\),
\(\texttt{certificates/orientation/k3e\_reduced\_orientation}\), and
\(\texttt{certificates/hybrid/determinant\_anchor\_translation\_weight}\)
shows that the action rows, generator-basis rows, and character values
are not present.
\end{proof}

The packet
\[
\texttt{certificates/orientation/rhomred\_finite\_stabilizer\_linearization\_character}
\]
verifies
\(\texttt{RHOMRED\_FINITE\_STABILIZER\_LINEARIZATION\_CHARACTER\_OBSTRUCTION\_VERIFIED}\):
it imports the row \(287\) packet, the reduced-orientation obstruction
ledger, and the determinant-anchor translation-weight packet, records
the \(H^1(BH;\mathbb F_2)\) coordinate criterion for row \(288\),
checks the three retained residual groups, and keeps the
orientation-line action, generator-basis, character-value,
\(H^1\)-coordinate, row \(289\) vanishing, transition,
vanishing-cycle, and protected-integration rows open.

\begin{proposition}[Vanishing criterion for the finite-stabilizer
linearization character]
\label{prop:finite-stabilizer-linearization-vanishing-criterion}
\textnormal{[criterion proved; zero-linearization rows not supplied]}
Assume row \(288\) has computed
\[
\lambda^H_{\mathcal C,S}\in H^1(BH;\mathbb F_2)
\]
for every retained finite-stabilizer stratum.  Row \(289\) proves
\[
\lambda^H_{\mathcal C,S}=0
\]
for every retained \(H\) precisely when it verifies the corresponding
zero coordinate row in the \(H^1(BH;\mathbb F_2)\)-basis:
\[
H=1\ \text{or}\ |H|\ \text{odd}:\quad
H^1(BH;\mathbb F_2)=0
\quad\text{after the retained action row is supplied,}
\]
\[
H\simeq E[2]:\quad \lambda_1=\lambda_2=0
\quad\text{equivalently}\quad
\rho_1=\rho_2=\rho_3=0
\]
with \(\rho_3=\lambda_1+\lambda_2\), and, for two-primary rank two
stabilizers,
\[
H_{(2)}\simeq(\mathbb Z/2^a)^2,\ a\ge2:
\quad \lambda^{(H)}_1=\lambda^{(H)}_2=0.
\]
This is a degree-one vanishing statement.  It is not implied by
\(\beta^H_{\mathcal C,S}=0\), by a row \(287\)
null-trivialization of the degree-two gerbe, by the determinant-anchor
translation-weight formula, or by class invariance in
\(\operatorname{Pic}\).  Row \(290\) records this independence as a
separate firewall.

The present finite tables do not prove row \(289\).  They record no
orientation-line action, no generator-basis row, no character values,
and no computed \(H^1\)-coordinate for \(\lambda^H_{\mathcal C,S}\).
The row \(288\) packet records three order-one residual group
templates, but the corresponding \(\lambda\)-values are
\(\texttt{missing}\), with
\(\texttt{lambda\_character\_computed}=\texttt{false}\) and
\(\texttt{lambda\_vanishing\_verified}=\texttt{false}\).  Hence the
zero \(H^1(B1;\mathbb F_2)\) template cannot be promoted to a
zero-linearization row for the retained orientation problem.
\end{proposition}

\begin{proof}
Once row \(288\) has produced an element of \(H^1(BH;\mathbb F_2)\),
vanishing is a coordinate calculation.  The groups \(H=1\) and odd
\(H\) have zero mod-\(2\) degree-one cohomology.  Lemma~\ref{lem:G-Klein-four}
identifies the \(E[2]\) character as
\(\lambda_1x_1+\lambda_2x_2\) and the cyclic restrictions as
\(\rho_1=\lambda_1\), \(\rho_2=\lambda_2\),
\(\rho_3=\lambda_1+\lambda_2\); all three cyclic restrictions vanish
if and only if \(\lambda_1=\lambda_2=0\).  Lemma~\ref{lem:G-two-primary-stabilizer}
identifies the two-primary degree-one basis \(x_1,x_2\), so vanishing
is equivalent to the two displayed coefficients being zero.  These
statements require the computed character from row \(288\).  Inspection
of
\(\texttt{certificates/orientation/rhomred\_finite\_stabilizer\_linearization\_character}\)
and \(\texttt{certificates/orientation/k3e\_reduced\_orientation}\)
shows that no such character value or zero-coordinate row is present.
\end{proof}

The packet
\[
\texttt{certificates/orientation/rhomred\_finite\_stabilizer\_linearization\_vanishing}
\]
verifies
\(\texttt{RHOMRED\_FINITE\_STABILIZER\_LINEARIZATION\_VANISHING\_OBSTRUCTION\_VERIFIED}\):
it imports the row \(288\) packet and the reduced-orientation
obstruction ledger, records the \(H^1(BH;\mathbb F_2)\) zero criteria
for row \(289\), checks the three retained residual groups, and keeps
the computed-character, zero-coordinate, all-retained-coverage,
transition, vanishing-cycle, and protected-integration rows open.

\begin{proposition}[Gerbe vanishing does not imply zero
finite-stabilizer linearization]
\label{prop:finite-stabilizer-gerbe-linearization-independence}
\textnormal{[firewall proved]}
The implication
\[
\beta^H_{\mathcal C,S}=0\quad\Longrightarrow\quad
\lambda^H_{\mathcal C,S}=0
\]
is false as a finite-stabilizer descent principle.  For
\(H\simeq E[2]\), the equations
\[
b_{20}=b_{11}=b_{02}=0
\]
kill the degree-two class
\(\beta^H_{\mathcal C,S}\in H^2(BH;\mathbb F_2)\).  The degree-one
class
\[
\lambda^H_{\mathcal C,S}=\lambda_1x_1+\lambda_2x_2
\in H^1(BH;\mathbb F_2)
\]
is independent; for example \((\lambda_1,\lambda_2)=(1,0)\) gives
\[
\rho_1=1,\qquad \rho_2=0,\qquad \rho_3=1,
\]
so the gerbe vanishes and the linearization does not.  For
\(H_{(2)}\simeq(\mathbb Z/2^a)^2\), \(a\ge2\), the same separation is
\[
A_1^{(H)}=A_{12}^{(H)}=A_2^{(H)}=0,\qquad
(\lambda^{(H)}_1,\lambda^{(H)}_2)=(1,0).
\]
Thus rows \(286\)--\(287\) cannot replace rows \(288\)--\(289\).  If
\(H=1\) or \(|H|\) is odd, then \(H^1(BH;\mathbb F_2)=0\); this is a
separate group-cohomological zero target, not a consequence of
\(\beta^H_{\mathcal C,S}=0\).

The present finite tables obey this firewall.  The row \(289\) packet
records \(\texttt{beta\_vanishing\_implies\_lambda\_vanishing}=
\texttt{false}\) and
\(\texttt{nulltrivialization\_implies\_lambda\_vanishing}=
\texttt{false}\); it also records missing \(\lambda\)-values and no
zero-linearization rows.  Hence no scalar trace, determinant-anchor
translation weight, gerbe vanishing row, or null-trivialization row is
used as a substitute for the degree-one linearization calculation.
\end{proposition}

\begin{proof}
The cohomological degrees already separate the two obstructions.
Lemma~\ref{lem:G-Klein-four} gives
\[
H^*(BE[2];\mathbb F_2)=\mathbb F_2[x_1,x_2],\qquad |x_i|=1.
\]
The finite gerbe lives in degree \(2\), with coordinates
\((b_{20},b_{11},b_{02})\), while the residual linearization lives in
degree \(1\), with coordinates \((\lambda_1,\lambda_2)\).  These are
coordinates on different graded pieces, so the zero locus of the
degree-two coordinates contains every value of the degree-one
coordinates.  The displayed choice \((\lambda_1,\lambda_2)=(1,0)\)
is an explicit counterexample to the implication.  Lemma~\ref{lem:G-two-primary-stabilizer}
gives the same degree separation for
\((\mathbb Z/2^a)^2\): the gerbe coordinates are
\((A_1^{(H)},A_{12}^{(H)},A_2^{(H)})\) in degree \(2\), and the
linearization coordinates are
\((\lambda^{(H)}_1,\lambda^{(H)}_2)\) in degree \(1\).  The trivial
and odd cases have zero degree-one target by transfer or by
triviality; that is not an implication from the degree-two class.
\end{proof}

The packet
\[
\texttt{certificates/orientation/rhomred\_finite\_stabilizer\_gerbe\_linearization\_independence}
\]
verifies
\(\texttt{RHOMRED\_FINITE\_STABILIZER\_GERBE\_LINEARIZATION\_INDEPENDENCE\_VERIFIED}\):
it imports the row \(289\) packet and the reduced-orientation
obstruction ledger, records explicit \(E[2]\) and two-primary
counterexamples to the false implication, records the trivial/odd
zero-target exception, and checks that no current orientation packet
uses a gerbe, scalar, or determinant-anchor row as a zero-linearization
row.

\begin{proposition}[\(E\lbrack2\rbrack\) Klein-four finite-stabilizer class]
\label{prop:finite-stabilizer-klein-four-class}
\textnormal{[local computation proved; retained \(E[2]\)-edge row not
supplied]}
Let \(H=E[2]\simeq(\mathbb Z/2)^2\), choose generators \(e_1,e_2\),
and let \(x_1,x_2\in H^1(BH;\mathbb F_2)\) be the dual basis.  After
the finite-stabilizer Cech--Borel class on a retained \(BE[2]\)-stratum
has been reduced to the classifying-space edge, its degree-two part is
uniquely of the form
\[
\beta^{E,2}_{\mathcal C,S}
=b_{20}x_1^2+b_{11}x_1x_2+b_{02}x_2^2
\in H^2(BE[2];\mathbb F_2).
\]
This is row \(291\).  The symbols \(b_{20},b_{11},b_{02}\) are the
coordinates of the supplied edge class in the ordered basis
\[
x_1^2,\quad x_1x_2,\quad x_2^2.
\]
Rows \(292\)--\(294\) are the subsequent restriction computations
which express these coordinates through the three cyclic restrictions
\(r_1,r_2,r_3\).  Row \(291\) does not by itself supply a retained
\(E[2]\)-stratum, an equivariant Cech--Borel representative, an
edge-reduction row, or the values of the three coefficients.
\end{proposition}

\begin{proof}
For a complex elliptic curve, \(E[2]\simeq(\mathbb Z/2)^2\).  The
classifying-space cohomology is
\[
H^*(BE[2];\mathbb F_2)=\mathbb F_2[x_1,x_2],\qquad |x_i|=1.
\]
Therefore its degree-two part has basis
\[
H^2(BE[2];\mathbb F_2)
=\mathbb F_2\langle x_1^2,x_1x_2,x_2^2\rangle.
\]
An edge-reduced finite-stabilizer Borel class is an element of this
three-dimensional vector space, hence has the unique displayed
coordinate expansion.  The proof is a computation in
\(H^*(B(\mathbb Z/2)^2;\mathbb F_2)\); it does not assert that the
current finite tables have supplied the geometric edge class whose
coordinates are \(b_{20},b_{11},b_{02}\).
\end{proof}

The packet
\[
\texttt{certificates/orientation/rhomred\_finite\_stabilizer\_klein\_four\_class}
\]
verifies
\(\texttt{RHOMRED\_FINITE\_STABILIZER\_KLEIN\_FOUR\_CLASS\_VERIFIED}\):
it imports the finite-stabilizer beta packet, the row \(290\) firewall,
and the reduced-orientation obstruction ledger, records the
\(H^2(BE[2];\mathbb F_2)\) basis and the displayed Klein-four edge
class, and keeps the retained \(E[2]\)-stratum, edge-reduction,
coefficient-value, cyclic-restriction, vanishing, linearization,
transition, vanishing-cycle, and protected-integration rows open.

\begin{proposition}[First cyclic restriction of the \(E\lbrack2\rbrack\)
Klein-four class]
\label{prop:finite-stabilizer-klein-four-b20}
\textnormal{[local computation proved; first cyclic restriction row
not supplied]}
Keep the notation of
Proposition~\ref{prop:finite-stabilizer-klein-four-class}.  Let
\(\iota_1\colon \langle e_1\rangle\hookrightarrow E[2]\) be the first
cyclic subgroup and let
\[
t\in H^1(B\langle e_1\rangle;\mathbb F_2)
\]
be its degree-one generator.  The restriction map satisfies
\[
\iota_1^*(x_1)=t,\qquad \iota_1^*(x_2)=0.
\]
Therefore
\[
\iota_1^*\beta^{E,2}_{\mathcal C,S}=b_{20}t^2.
\]
If the first cyclic restriction is written
\[
\iota_1^*\beta^{E,2}_{\mathcal C,S}=r_1t^2,
\]
then
\[
b_{20}=r_1.
\]
This is row \(292\).  It is the first coefficient extraction from the
row \(291\) class.  It does not supply the retained \(E[2]\)-edge row,
the first cyclic restriction row, or the value of \(r_1\).
\end{proposition}

\begin{proof}
The inclusion \(\iota_1\) sends the generator of
\(\langle e_1\rangle\) to \(e_1\).  Hence the dual class \(x_1\)
restricts to the generator \(t\), while \(x_2\) restricts to zero.
Applying \(\iota_1^*\) to
\[
\beta^{E,2}_{\mathcal C,S}
=b_{20}x_1^2+b_{11}x_1x_2+b_{02}x_2^2
\]
gives \(b_{20}t^2\).  Since
\(H^2(B\langle e_1\rangle;\mathbb F_2)=\mathbb F_2\langle t^2\rangle\),
equality with \(r_1t^2\) is equivalent to equality of coefficients,
namely \(b_{20}=r_1\).
\end{proof}

The packet
\[
\texttt{certificates/orientation/rhomred\_finite\_stabilizer\_klein\_four\_b20}
\]
verifies
\(\texttt{RHOMRED\_FINITE\_STABILIZER\_KLEIN\_FOUR\_B20\_VERIFIED}\):
it imports the row \(291\) Klein-four class packet, records the first
cyclic restriction map \(x_1\mapsto t,\ x_2\mapsto0\), verifies the
coefficient identity \(b_{20}=r_1\), and keeps the retained edge row,
the first cyclic restriction value, the \(b_{11}\) and \(b_{02}\)
restriction computations, vanishing, linearization, transition,
vanishing-cycle, and protected-integration rows open.

\begin{proposition}[Diagonal cyclic restriction of the \(E\lbrack2\rbrack\)
Klein-four class]
\label{prop:finite-stabilizer-klein-four-b11}
\textnormal{[local computation proved; cyclic restriction rows not
supplied]}
Keep the notation of
Propositions~\ref{prop:finite-stabilizer-klein-four-class}
and~\ref{prop:finite-stabilizer-klein-four-b20}.  Let
\[
\iota_2\colon \langle e_2\rangle\hookrightarrow E[2],
\qquad
\iota_3\colon \langle e_1+e_2\rangle\hookrightarrow E[2]
\]
be the second and diagonal cyclic subgroups, and let \(t\) denote the
degree-one generator on either cyclic subgroup.  The restriction maps
satisfy
\[
\iota_2^*(x_1)=0,\qquad \iota_2^*(x_2)=t,
\]
and
\[
\iota_3^*(x_1)=t,\qquad \iota_3^*(x_2)=t.
\]
Thus
\[
\iota_2^*\beta^{E,2}_{\mathcal C,S}=b_{02}t^2,\qquad
\iota_3^*\beta^{E,2}_{\mathcal C,S}
=(b_{20}+b_{11}+b_{02})t^2.
\]
If the three cyclic restrictions are written
\[
\iota_i^*\beta^{E,2}_{\mathcal C,S}=r_it^2,\qquad i=1,2,3,
\]
where \(i=3\) denotes the diagonal subgroup, then
\[
b_{11}=r_1+r_2+r_3
\]
in \(\mathbb F_2\).  This is row \(293\).  It does not supply the
retained \(E[2]\)-edge row, the three cyclic restriction rows, the
values of \(r_1,r_2,r_3\), or the value of \(b_{11}\).  Row \(294\)
records the second-square extraction \(b_{02}=r_2\) as its own ledger
row.
\end{proposition}

\begin{proof}
The inclusion \(\iota_2\) kills the dual class \(x_1\) and sends
\(x_2\) to the cyclic generator \(t\).  The inclusion \(\iota_3\)
sends the cyclic generator to \(e_1+e_2\), so both dual classes
restrict to \(t\).  Applying these restrictions to
\[
\beta^{E,2}_{\mathcal C,S}
=b_{20}x_1^2+b_{11}x_1x_2+b_{02}x_2^2
\]
gives the two displayed formulae.  The first cyclic restriction gives
\(b_{20}=r_1\) by
Proposition~\ref{prop:finite-stabilizer-klein-four-b20}; the second
restriction gives \(b_{02}=r_2\); and the diagonal restriction gives
\[
b_{20}+b_{11}+b_{02}=r_3.
\]
Substitution yields \(r_1+b_{11}+r_2=r_3\).  Since the coefficient
field is \(\mathbb F_2\), this is equivalent to
\[
b_{11}=r_1+r_2+r_3 .
\]
All equalities here are identities in the cohomology of cyclic
classifying spaces.  They do not assert that the current finite
tables have supplied the geometric cyclic restrictions or their
values.
\end{proof}

The packet
\[
\texttt{certificates/orientation/rhomred\_finite\_stabilizer\_klein\_four\_b11}
\]
verifies
\(\texttt{RHOMRED\_FINITE\_STABILIZER\_KLEIN\_FOUR\_B11\_VERIFIED}\):
it imports the row \(291\) and row \(292\) packets, records the second
and diagonal cyclic restriction maps, verifies the mixed coefficient
identity \(b_{11}=r_1+r_2+r_3\), and keeps the retained edge row,
the cyclic restriction values, the standalone \(b_{02}\) row,
vanishing, linearization, transition, vanishing-cycle, and
protected-integration rows open.

\begin{proposition}[Second cyclic restriction of the \(E\lbrack2\rbrack\)
Klein-four class]
\label{prop:finite-stabilizer-klein-four-b02}
\textnormal{[local computation proved; second cyclic restriction row
not supplied]}
Keep the notation of
Proposition~\ref{prop:finite-stabilizer-klein-four-class}.  Let
\(\iota_2\colon \langle e_2\rangle\hookrightarrow E[2]\) be the second
cyclic subgroup and let
\[
t\in H^1(B\langle e_2\rangle;\mathbb F_2)
\]
be its degree-one generator.  The restriction map satisfies
\[
\iota_2^*(x_1)=0,\qquad \iota_2^*(x_2)=t.
\]
Therefore
\[
\iota_2^*\beta^{E,2}_{\mathcal C,S}=b_{02}t^2.
\]
If the second cyclic restriction is written
\[
\iota_2^*\beta^{E,2}_{\mathcal C,S}=r_2t^2,
\]
then
\[
b_{02}=r_2.
\]
This is row \(294\).  It is the standalone second-square coefficient
extraction from the row \(291\) class.  It does not supply the
retained \(E[2]\)-edge row, the second cyclic restriction row, or the
value of \(r_2\).
\end{proposition}

\begin{proof}
The inclusion \(\iota_2\) sends the generator of
\(\langle e_2\rangle\) to \(e_2\).  Hence \(x_1\) restricts to zero
and \(x_2\) restricts to the generator \(t\).  Applying \(\iota_2^*\)
to
\[
\beta^{E,2}_{\mathcal C,S}
=b_{20}x_1^2+b_{11}x_1x_2+b_{02}x_2^2
\]
gives \(b_{02}t^2\).  Since
\(H^2(B\langle e_2\rangle;\mathbb F_2)=\mathbb F_2\langle t^2\rangle\),
equality with \(r_2t^2\) is equivalent to \(b_{02}=r_2\).  This is a
classifying-space calculation only; it does not assert that the
finite-stabilizer tables have supplied the retained cyclic
restriction or its value.
\end{proof}

The packet
\[
\texttt{certificates/orientation/rhomred\_finite\_stabilizer\_klein\_four\_b02}
\]
verifies
\(\texttt{RHOMRED\_FINITE\_STABILIZER\_KLEIN\_FOUR\_B02\_VERIFIED}\):
it imports the row \(291\) and row \(293\) packets, records the second
cyclic restriction map \(x_1\mapsto0,\ x_2\mapsto t\), verifies the
coefficient identity \(b_{02}=r_2\), and keeps the retained edge row,
the second cyclic restriction value, vanishing, linearization,
transition, vanishing-cycle, and protected-integration rows open.

\begin{proposition}[Klein-four zero-restriction criterion]
\label{prop:finite-stabilizer-klein-four-zero-restrictions}
\textnormal{[criterion proved; zero restriction rows not supplied]}
Keep the notation of
Propositions~\ref{prop:finite-stabilizer-klein-four-class},
\ref{prop:finite-stabilizer-klein-four-b20},
\ref{prop:finite-stabilizer-klein-four-b11}, and
\ref{prop:finite-stabilizer-klein-four-b02}.  Let
\[
\iota_i^*\beta^{E,2}_{\mathcal C,S}=r_it^2,\qquad i=1,2,3,
\]
where \(i=1\) is the subgroup \(\langle e_1\rangle\), \(i=2\) is
\(\langle e_2\rangle\), and \(i=3\) is
\(\langle e_1+e_2\rangle\).  Then
\[
\beta^{E,2}_{\mathcal C,S}=0
\quad\Longleftrightarrow\quad
r_1=r_2=r_3=0.
\]
Equivalently, the three cyclic zero restrictions are necessary and
sufficient for the vanishing of the classifying-space edge of the
\(E[2]\)-finite-stabilizer Borel class.  This is the algebraic content
of row \(295\).  The present retained finite-stabilizer tables do not
yet supply the three cyclic restriction rows or the values
\(r_1,r_2,r_3\), so this proposition is not a geometric proof that
those values vanish on the retained \(K3\times E\) strata.
\end{proposition}

\begin{proof}
Rows \(292\)--\(294\) give
\[
b_{20}=r_1,\qquad b_{02}=r_2,\qquad
b_{11}=r_1+r_2+r_3
\]
in \(\mathbb F_2\).  If \(r_1=r_2=r_3=0\), then
\(b_{20}=b_{02}=b_{11}=0\), hence
\[
\beta^{E,2}_{\mathcal C,S}
=b_{20}x_1^2+b_{11}x_1x_2+b_{02}x_2^2
\]
vanishes in \(H^2(BE[2];\mathbb F_2)\).  Conversely, if
\(\beta^{E,2}_{\mathcal C,S}=0\), then the coefficients in the basis
\(x_1^2,x_1x_2,x_2^2\) vanish.  Thus \(r_1=b_{20}=0\),
\(r_2=b_{02}=0\), and
\[
r_3=b_{20}+b_{11}+b_{02}=0.
\]
The implication is a classifying-space calculation.  A proof of the
geometric assertion \(r_1=r_2=r_3=0\) still requires retained cyclic
restriction rows, or an equivalent edge-reduced vanishing row, in the
finite-stabilizer orientation tables.
\end{proof}

The packet
\[
\texttt{certificates/orientation/rhomred\_finite\_stabilizer\_klein\_four\_zero\_restrictions}
\]
verifies
\(\texttt{RHOMRED\_FINITE\_STABILIZER\_KLEIN\_FOUR\_ZERO\_RESTRICTIONS\_OBSTRUCTION\_VERIFIED}\):
it imports the row \(291\)--\(294\) packets, proves the equivalence
between the three zero cyclic restrictions and vanishing of the
Klein-four classifying-space edge, and records that the retained
cyclic zero rows, the \(r_i\)-values, finite-stabilizer beta
vanishing, quotient orientation, transition compatibility,
vanishing-cycle, and protected-integration rows remain open.

\begin{proposition}[Even two-primary finite-stabilizer class]
\label{prop:finite-stabilizer-two-primary-class}
\textnormal{[local computation proved; retained two-primary row not
supplied]}
Let \(N\ge4\) be even and write \(2^a\parallel N\), \(a\ge2\).  The
two-primary part of the \(N\)-torsion stabilizer of a complex elliptic
curve is
\[
E[N]_{(2)}\simeq(\mathbb Z/2^a)^2.
\]
Choose generators and let
\[
x_1,x_2\in H^1(B(\mathbb Z/2^a)^2;\mathbb F_2),
\qquad
y_1,y_2\in H^2(B(\mathbb Z/2^a)^2;\mathbb F_2)
\]
be the corresponding exterior and polynomial generators.  Then
\[
H^*\!\left(B(\mathbb Z/2^a)^2;\mathbb F_2\right)
=\Lambda(x_1,x_2)\otimes\mathbb F_2[y_1,y_2],
\qquad |x_i|=1,\quad |y_i|=2,
\]
and
\[
H^2\!\left(B(\mathbb Z/2^a)^2;\mathbb F_2\right)
=\mathbb F_2\langle y_1,x_1x_2,y_2\rangle.
\]
After the finite-stabilizer Cech--Borel class on a retained
two-primary stratum has been reduced to the classifying-space edge,
its degree-two part is uniquely of the form
\[
\beta^{E,N}_{\mathcal C,S}
=A_1^{(N)}y_1+A_{12}^{(N)}x_1x_2+A_2^{(N)}y_2.
\]
This is row \(296\).  The coefficients
\(A_1^{(N)},A_{12}^{(N)},A_2^{(N)}\) are coordinates of the supplied
edge class in the ordered basis \(y_1,x_1x_2,y_2\).  Row \(296\) does
not supply a retained two-primary stratum, an equivariant
Cech--Borel representative, an edge-reduction row, or the values of
the three coefficients.  Row \(297\) is the separate mixed-term
detection check for \(A_{12}^{(N)}\).
\end{proposition}

\begin{proof}
For a complex elliptic curve, \(E[N]\simeq(\mathbb Z/N)^2\).  Primary
decomposition gives the displayed two-primary subgroup.  For
\(a\ge2\),
\[
H^*(B\mathbb Z/2^a;\mathbb F_2)=\Lambda(x)\otimes\mathbb F_2[y],
\qquad |x|=1,\quad |y|=2.
\]
The K\"unneth formula gives
\[
H^*\!\left(B(\mathbb Z/2^a)^2;\mathbb F_2\right)
=\Lambda(x_1,x_2)\otimes\mathbb F_2[y_1,y_2].
\]
Its degree-two part has basis \(y_1,x_1x_2,y_2\).  Therefore an
edge-reduced two-primary finite-stabilizer Borel class is a unique
linear combination of these three basis vectors.  The argument is
only a classifying-space computation; it does not assert that the
current finite tables have supplied the retained edge class whose
coordinates are \(A_1^{(N)},A_{12}^{(N)},A_2^{(N)}\).
\end{proof}

The packet
\[
\texttt{certificates/orientation/rhomred\_finite\_stabilizer\_two\_primary\_class}
\]
verifies
\(\texttt{RHOMRED\_FINITE\_STABILIZER\_TWO\_PRIMARY\_CLASS\_VERIFIED}\):
it imports the finite-stabilizer beta packet and the reduced
orientation obstruction ledger, records the
\(\Lambda(x_1,x_2)\otimes\mathbb F_2[y_1,y_2]\) cohomology ring,
the degree-two basis \(y_1,x_1x_2,y_2\), and the displayed
two-primary edge-class form, and keeps the retained stratum,
edge-reduction, coefficient-value, mixed-term detection, vanishing,
linearization, transition, vanishing-cycle, and protected-integration
rows open.

\begin{proposition}[Detection of the two-primary mixed term]
\label{prop:finite-stabilizer-two-primary-mixed-term}
\textnormal{[local detection proved; retained \(A_{12}\)-value not
supplied]}
Keep the notation of
Proposition~\ref{prop:finite-stabilizer-two-primary-class}.  Define
the coefficient functional
\[
\pi_{12}\colon
H^2\!\left(B(\mathbb Z/2^a)^2;\mathbb F_2\right)\longrightarrow
\mathbb F_2
\]
by
\[
\pi_{12}(y_1)=0,\qquad
\pi_{12}(x_1x_2)=1,\qquad
\pi_{12}(y_2)=0.
\]
Then
\[
\pi_{12}\!\left(\beta^{E,N}_{\mathcal C,S}\right)=A_{12}^{(N)}.
\]
For every cyclic subgroup \(C\subset(\mathbb Z/2^a)^2\), the
restriction map satisfies
\[
\operatorname{res}^{(\mathbb Z/2^a)^2}_{C}(x_1x_2)=0.
\]
Thus the mixed coefficient \(A_{12}^{(N)}\) is detected only by the
rank-two classifying-space edge, or by an equivalent rank-two
coefficient row; it is not determined by cyclic restrictions.  This
is row \(297\).  It does not supply the retained two-primary edge row
or the value of \(A_{12}^{(N)}\).
\end{proposition}

\begin{proof}
The first assertion is the definition of the coordinate functional in
the ordered basis \(y_1,x_1x_2,y_2\).  It remains to check the cyclic
restriction statement.  Let \(C\subset(\mathbb Z/2^a)^2\) be cyclic.
If \(|C|=2\), then \(C\subset 2(\mathbb Z/2^a)^2\) because \(a\ge2\);
therefore the two mod-\(2\) characters \(x_1,x_2\) restrict trivially
to \(C\).  If \(|C|\ge4\), then
\[
H^*(BC;\mathbb F_2)=\Lambda(u)\otimes\mathbb F_2[v],
\qquad |u|=1,\quad |v|=2,
\]
and both \(x_1\) and \(x_2\) restrict to scalar multiples of \(u\).
Hence their product restricts to a scalar multiple of \(u^2=0\).
Thus \(x_1x_2\) restricts to zero on every cyclic subgroup.  The
coefficient \(A_{12}^{(N)}\) can therefore not be recovered from
cyclic restrictions of the degree-two class.
\end{proof}

The packet
\[
\texttt{certificates/orientation/rhomred\_finite\_stabilizer\_two\_primary\_mixed\_term\_detection}
\]
verifies
\(\texttt{RHOMRED\_FINITE\_STABILIZER\_TWO\_PRIMARY\_MIXED\_TERM\_DETECTION\_VERIFIED}\):
it imports the row \(296\) two-primary class packet, records the
rank-two coefficient functional \(\pi_{12}\), proves that cyclic
subgroup restrictions annihilate \(x_1x_2\), and keeps the retained
edge row, \(A_{12}^{(N)}\)-value, vanishing, linearization,
transition, vanishing-cycle, and protected-integration rows open.

\begin{proposition}[Odd-order finite-stabilizer transfer]
\label{prop:finite-stabilizer-odd-transfer}
\textnormal{[local transfer proved; retained odd edge row not
supplied]}
Let \(H\subset E[N]\) be a finite stabilizer of odd order.  Then
\[
H^i(BH;\mathbb F_2)=0,\qquad i>0.
\]
Consequently, after the finite-stabilizer Cech--Borel class on a
retained odd-order stratum has been reduced to the
classifying-space edge,
\[
\beta^H_{\mathcal C,S}\in H^2(BH;\mathbb F_2)
\]
vanishes.  The residual degree-one stabilizer character also vanishes
after edge reduction, since \(H^1(BH;\mathbb F_2)=0\).  This is row
\(298\).  It proves the odd-order target calculation; it does not
supply a retained odd-order stratum, an equivariant Cech--Borel
representative, an edge-reduction row, an all-retained-stabilizer
coverage row, or the compatible null-trivialisations required for
(O1).
\end{proposition}

\begin{proof}
Since \(|H|\) is odd, multiplication by \(|H|\) is the identity on
\(\mathbb F_2\).  The restriction-transfer identity in group
cohomology gives multiplication by \(|H|\) on
\(H^i(BH;\mathbb F_2)\), while restriction to the point is zero in
positive degree.  Thus every positive-degree class in
\(H^*(BH;\mathbb F_2)\) is zero.  Equivalently, the averaging
idempotent makes the invariants functor exact because \(|H|\) is
invertible in \(\mathbb F_2\).  In particular, the degree-two
classifying-space edge of the finite-stabilizer Borel class and the
degree-one residual character both vanish once the geometric data have
been reduced to those targets.
\end{proof}

The packet
\[
\texttt{certificates/orientation/rhomred\_finite\_stabilizer\_odd\_transfer}
\]
verifies
\(\texttt{RHOMRED\_FINITE\_STABILIZER\_ODD\_TRANSFER\_VERIFIED}\):
it imports the finite-stabilizer beta and beta-vanishing obstruction
packets, records the odd-order transfer calculation
\(H^{>0}(BH;\mathbb F_2)=0\), and keeps the retained odd stratum,
edge-reduction, all-retained coverage, null-trivialisation,
transition, vanishing-cycle, and protected-integration rows open.

\begin{proposition}[Finite-stabilizer orientation detection;
\textit{proved after edge reduction}]
\label{prop:finite-stabilizer-orientation-detection}
Let \(S\) be a retained finite-stabilizer stratum with stabilizer
\(H\subset E_{\mathrm{tors}}\).  Suppose the equivariant
Cech--Borel orientation class
\(\widetilde\beta^H_{\mathcal C,S}\in H^2_H(S;\F_2)\) has been
reduced to its classifying-space edge, i.e. its ordinary component in
\(H^2(S;\F_2)^H\), mixed component in
\(H^1(BH;H^1(S;\F_2))\), and incoming spectral-sequence differential
terms vanish with chosen null-trivialisations.  Then the finite
stabilizer part of (O1) is exactly the following finite calculation.

If \(H\) has odd order, the mod-\(2\) classifying-space edge vanishes.
If \(H\simeq E[2]\simeq(\Z/2)^2\), write
\[
\beta^H_{\mathcal C,S}
=b_{20}x_1^2+b_{11}x_1x_2+b_{02}x_2^2,\qquad
\lambda^H_{\mathcal C,S}=\lambda_1x_1+\lambda_2x_2.
\]
Then quotient orientation on this stratum is equivalent to
\[
b_{20}=b_{11}=b_{02}=\lambda_1=\lambda_2=0,
\]
or equivalently to the three cyclic quadratic restrictions and the
three cyclic linearization restrictions all vanishing, as in
Lemma~\ref{lem:G-Klein-four}.  If the two-primary part of \(H\) is
\((\Z/2^a)^2\), \(a\ge2\), write
\[
\beta^H_{\mathcal C,S}
=A_1^{(H)}y_1+A_{12}^{(H)}x_1x_2+A_2^{(H)}y_2,\qquad
\lambda^H_{\mathcal C,S}=\lambda_1^{(H)}x_1+\lambda_2^{(H)}x_2.
\]
Then quotient orientation is equivalent to the five equations
\[
A_1^{(H)}=A_{12}^{(H)}=A_2^{(H)}
=\lambda_1^{(H)}=\lambda_2^{(H)}=0.
\]
The mixed coefficient \(A_{12}^{(H)}\) is not detected by cyclic
order-two restrictions.  Therefore a finite-stabilizer row which lists
only cyclic restrictions is not an (O1) quotient-orientation row unless
the edge-reduction and mixed-coefficient computations are supplied.
\end{proposition}

\begin{proof}
The edge-reduction hypothesis identifies the relevant obstruction with
an element of \(H^2(BH;\F_2)\), and the residual linearization with an
element of \(H^1(BH;\F_2)\).  For odd \(H\), transfer kills positive
degree cohomology with \(\F_2\)-coefficients.  For \(E[2]\), the
coordinate calculation is Lemma~\ref{lem:G-Klein-four}: the three
nonzero cyclic restrictions invert the quadratic coefficients, and the
linearization character is an independent \(H^1\)-torsor.  For
\((\Z/2^a)^2\), \(a\ge2\), Lemma~\ref{lem:G-two-primary-stabilizer}
gives
\[
H^2(B(\Z/2^a)^2;\F_2)
=\F_2\langle y_1,x_1x_2,y_2\rangle.
\]
The term \(x_1x_2\) pulls back trivially to cyclic order-two
subgroups, so it must be computed on the rank-two stabilizer.  These
are exactly the equations in the statement.
\end{proof}

\begin{proposition}[Direct-sum orientation multiplicativity criterion]
\label{prop:orientation-direct-sum-multiplicativity}
\textnormal{[Picard-groupoid criterion proved; direct-sum
Thom--Sebastiani row not supplied]}
Let \(\mathcal C,\mathcal D\) be retained objects at a fixed finite
HN height.  Suppose that the cosection-reduced self-Ext construction
has been supplied on the direct-sum chart and is compatible with the
block decomposition
\[
K^{\mathrm{red}}_{\mathcal C\oplus\mathcal D}\simeq
K^{\mathrm{red}}_{\mathcal C}\oplus
K^{\mathrm{red}}_{\mathcal D}\oplus
K^{\mathrm{red}}_{\mathcal C,\mathcal D}\oplus
K^{\mathrm{red}}_{\mathcal D,\mathcal C}.
\]
Then the determinant functor gives
\[
\mathcal L^{\det}_{\mathcal C\oplus\mathcal D}\simeq
\mathcal L^{\det}_{\mathcal C}\otimes
\mathcal L^{\det}_{\mathcal D}\otimes
\det K^{\mathrm{red}}_{\mathcal C,\mathcal D}\otimes
\det K^{\mathrm{red}}_{\mathcal D,\mathcal C}.
\]
Direct-sum multiplicativity of orientations is therefore not the
displayed determinant identity.  It is the additional
Picard-groupoid datum consisting of square roots
\(o_{\mathcal C}^{\otimes2}\simeq\mathcal L^{\det}_{\mathcal C}\),
\(o_{\mathcal D}^{\otimes2}\simeq\mathcal L^{\det}_{\mathcal D}\),
\(o_{\mathcal C\oplus\mathcal D}^{\otimes2}\simeq
\mathcal L^{\det}_{\mathcal C\oplus\mathcal D}\), a chosen
Calabi--Yau hyperbolic trivialization
\[
\det K^{\mathrm{red}}_{\mathcal C,\mathcal D}\otimes
\det K^{\mathrm{red}}_{\mathcal D,\mathcal C}
\simeq q_{\mathcal C,\mathcal D}^{\otimes2},
\]
and an isomorphism
\[
m_{\oplus,\mathcal C,\mathcal D}:
o_{\mathcal C\oplus\mathcal D}
\xrightarrow{\sim}
o_{\mathcal C}\otimes o_{\mathcal D}\otimes
q_{\mathcal C,\mathcal D}
\]
which is symmetric in \((\mathcal C,\mathcal D)\), functorial under
pullback to the retained direct-sum and two-step flag charts, and has
zero Thom--Sebastiani pentagon defect.  Row \(299\) is proved for a
retained pair exactly when this datum is supplied for that pair and
recorded as a zero-defect row in
\(\mathrm{ts\_multiplicativity}\).

In the present finite packet the table
\(\mathrm{orientation\_lines}\) has no square-root rows and
\(\mathrm{ts\_multiplicativity}\) has no direct-sum isomorphism row.
Thus the current status of row \(299\) is the criterion above together
with an open obstruction, not a proof of direct-sum multiplicativity.
\end{proposition}

\begin{proof}
The first displayed isomorphism is the standard block decomposition of
derived endomorphisms of a direct sum, after restricting to a chart
where the reduced cone construction has been made compatible with
direct sum.  The Knudsen--Mumford determinant functor is multiplicative
for finite direct sums of perfect complexes, hence gives the
determinant-line formula.  The off-diagonal factors are not killed by
this formula; they must be paired by the Calabi--Yau duality and
trivialized as a square before one can compare square roots.

A reduced orientation is a square root in the Picard groupoid, and a
multiplicative reduced orientation is therefore the isomorphism
\(m_{\oplus,\mathcal C,\mathcal D}\) together with the symmetry,
pullback, and pentagon coherences stated above.  These are exactly the
entries carried by a direct-sum Thom--Sebastiani row.  The current
orientation ledger records the square-root row and the
Thom--Sebastiani multiplicativity row as missing open obligations, and
the corresponding tables are empty.  Hence the proposition proves the
criterion and identifies the missing data; it does not assert that the
missing data have been constructed.
\end{proof}

The packet
\[
\texttt{certificates/orientation/rhomred\_orientation\_direct\_sum\_multiplicativity}
\]
verifies
\(\texttt{RHOMRED\_ORIENTATION\_DIRECT\_SUM\_MULTIPLICATIVITY\_OBSTRUCTION\_VERIFIED}\):
it imports the reduced determinant, square-root obstruction, and
reduced-orientation ledger packets, records the direct-sum
Picard-groupoid criterion, checks that
\(\mathrm{orientation\_lines}\) and
\(\mathrm{ts\_multiplicativity}\) are empty, and keeps the square-root,
cross-term hyperbolic trivialization, direct-sum
Thom--Sebastiani isomorphism, pentagon, extension, quotient,
transition, vanishing-cycle, and protected-integration rows open.

\begin{proposition}[Extension orientation multiplicativity criterion]
\label{prop:orientation-extension-multiplicativity}
\textnormal{[criterion proved; extension Thom--Sebastiani row not
supplied]}
Let
\[
\mathfrak M^{\mathrm{red}}_{R,\widehat c}\times
\mathfrak M^{\mathrm{red}}_{R,\widehat c'}
\xleftarrow{\ p=(p_1,p_2)\ }
\overline{\mathfrak E}^{\mathrm{red}}_{R,\widehat c,\widehat c'}
\xrightarrow{\ q\ }
\mathfrak M^{\mathrm{red}}_{R,\widehat c+\widehat c'}
\]
be a retained compactified extension correspondence.  Row \(300\)
is an orientation row only if it supplies a Joyce--Upmeier
multiplicativity isomorphism
\[
\operatorname{TS}^{o}_{\mathfrak E}:
p_1^*o_{R,\widehat c}\otimes p_2^*o_{R,\widehat c'}
\xrightarrow{\sim}q^*o_{R,\widehat c+\widehat c'}
\]
whose square is the cosection-reduced determinant comparison on
\(\overline{\mathfrak E}^{\mathrm{red}}_{R,\widehat c,\widehat c'}\).
Equivalently, after tensoring with the BBDJS reduced
Thom--Sebastiani isomorphism for vanishing cycles, it gives the
coefficient-system isomorphism
\[
p^*\bigl((\Phi^{\mathrm{red}}_{R,\widehat c}\otimes
o_{R,\widehat c})\boxtimes
(\Phi^{\mathrm{red}}_{R,\widehat c'}\otimes
o_{R,\widehat c'})\bigr)
\xrightarrow{\sim}
q^*(\Phi^{\mathrm{red}}_{R,\widehat c+\widehat c'}\otimes
o_{R,\widehat c+\widehat c'}).
\]
It must be functorial on overlaps, compatible with quotient descent,
and its pullbacks to every retained two-step flag stack must have zero
Thom--Sebastiani pentagon defect.  Thus row \(300\) is proved exactly
by a zero-defect row in \(\mathrm{ts\_multiplicativity}\) with
\[
(\mathrm{extension\_stack\_id},\mathrm{left\_line\_id},
\mathrm{right\_line\_id},\mathrm{target\_line\_id},
\mathrm{ts\_isomorphism\_id}).
\]

The retained extension-closure rows record only that certain middle
HN types stay inside the finite window.  The present orientation table
\(\mathrm{ts\_multiplicativity}\) is empty, and the mixed
Thom--Sebastiani packet is conditional on supplied orientation
transport.  Consequently row \(300\) is presently a criterion with an
open extension-orientation obstruction.
\end{proposition}

\begin{proof}
The Hall product pulls the external coefficient system on
\(\mathfrak M^{\mathrm{red}}_{R,\widehat c}\times
\mathfrak M^{\mathrm{red}}_{R,\widehat c'}\) to the extension stack
and pushes it along \(q\).  The BBDJS Thom--Sebastiani theorem
identifies the reduced vanishing-cycle factors after the
reduced d-critical charts and the extension potential comparison have
been supplied.  That theorem does not choose square roots of the
orientation gerbe.  Joyce--Upmeier multiplicativity is the additional
Picard-groupoid isomorphism
\(\operatorname{TS}^{o}_{\mathfrak E}\) whose square is the determinant
comparison for the exact-triangle self-Ext complex.

Associativity of the oriented Hall product requires the same
isomorphism to satisfy the two-step flag pentagon; quotient descent
and HN transition compatibility are separate rows.  The current
certificate
\(\texttt{certificates/orientation/k3e\_reduced\_orientation}\) has no
rows in \(\mathrm{orientation\_lines}\) or
\(\mathrm{ts\_multiplicativity}\), and its obstruction ledger records
both \(\mathrm{ts\_multiplicativity}\) and
\(\mathrm{ts\_pentagon}\) as missing open obligations.  Retained
extension closure and conditional mixed Thom--Sebastiani transport
therefore do not prove row \(300\).
\end{proof}

The packet
\[
\texttt{certificates/orientation/rhomred\_orientation\_extension\_multiplicativity}
\]
verifies
\(\texttt{RHOMRED\_ORIENTATION\_EXTENSION\_MULTIPLICATIVITY\_OBSTRUCTION\_VERIFIED}\):
it imports retained extension closure, the reduced-orientation ledger,
the square-root obstruction packet, the direct-sum multiplicativity
criterion, and the mixed Thom--Sebastiani packet; it checks that the
extension-closure rows do not construct extension stacks or
orientations, that the mixed Thom--Sebastiani packet is conditional on
orientation transport, and that
\(\mathrm{orientation\_lines}\) and
\(\mathrm{ts\_multiplicativity}\) are empty.

\begin{proposition}[Thom--Sebastiani compatibility for orientation
lines]
\label{prop:orientation-thom-sebastiani-compatibility}
\textnormal{[compatibility criterion proved; orientation-line
Thom--Sebastiani row not supplied]}
Let \(\operatorname{TS}^{o}_{\mathfrak E}\) be a candidate
orientation-line isomorphism on a retained extension correspondence as
in Proposition~\ref{prop:orientation-extension-multiplicativity}.
It is compatible with Thom--Sebastiani if and only if the following
finite row is supplied.

\begin{enumerate}
\item The square-root diagram commutes, equivalently
\[
q^*\sigma_{\widehat c+\widehat c'}\circ
(\operatorname{TS}^{o}_{\mathfrak E})^{\otimes2}
=
\det(\operatorname{TS}^{\mathrm{red}}_{\mathfrak E})\circ
(\sigma_{\widehat c}\otimes\sigma_{\widehat c'}),
\]
as morphisms from
\(\bigl(p_1^*o_{R,\widehat c}\otimes
p_2^*o_{R,\widehat c'}\bigr)^{\otimes2}\) to
\(q^*\mathcal L^{\det}_{R,\widehat c+\widehat c'}\).
\item The oriented coefficient-system map obtained by tensoring
\(\operatorname{TS}^{o}_{\mathfrak E}\) with the BBDJS reduced
Thom--Sebastiani isomorphism is the row used by the Hall product:
\[
\operatorname{TS}^{\Phi\otimes o}_{\mathfrak E}
=\operatorname{TS}^{\Phi,\mathrm{red}}_{\mathfrak E}\otimes
\operatorname{TS}^{o}_{\mathfrak E}.
\]
\item The chart-overlap cocycle, quotient descent cocycle, and
two-step flag pentagon defect of this oriented
Thom--Sebastiani row vanish.
\end{enumerate}
Thus row \(301\) is not a new scalar identity and not merely the BBDJS
vanishing-cycle Thom--Sebastiani theorem.  It is the zero-defect
compatibility row saying that the Joyce--Upmeier orientation
isomorphism is the square root of the determinant Thom--Sebastiani
comparison and that its tensor with the vanishing-cycle
Thom--Sebastiani map is the coefficient-system transport.

The present packets do not supply such a row.  The hybrid
Thom--Sebastiani packets are conditional on supplied orientation
transport, the four-input pentagon packet treats the orientation
pentagon as a functorial input, and
\(\mathrm{orientation\_lines}\) and
\(\mathrm{ts\_multiplicativity}\) are empty.
\end{proposition}

\begin{proof}
An orientation line is a square root of the determinant line in the
Picard groupoid.  Therefore compatibility of an orientation
Thom--Sebastiani isomorphism is exactly compatibility after squaring:
the square of \(\operatorname{TS}^{o}_{\mathfrak E}\) must be the
determinant-line comparison induced by the reduced exact-triangle
self-Ext complex.  This proves the first condition.

BBDJS Thom--Sebastiani identifies the vanishing-cycle factors after
the reduced d-critical chart and potential comparison have been
supplied.  It does not choose the Joyce--Upmeier square root.  The
oriented Hall coefficient system is
\(\Phi^{\mathrm{red}}\otimes o\), so its Thom--Sebastiani transport is
the tensor product of the vanishing-cycle transport and the
orientation-line transport.  This gives the second condition.

The same row must be invariant under chart changes, descend through
the quotient, and satisfy the two-step flag pentagon; otherwise the
oriented Hall product depends on choices or fails associativity.  The
current orientation ledger records the relevant orientation-line,
Thom--Sebastiani, and pentagon rows as missing open obligations.  The
mixed and four-input Thom--Sebastiani packets use orientation
compatibility as a hypothesis.  Hence the proposition proves the
criterion and records the missing data, not an unconditional
compatibility theorem.
\end{proof}

The packet
\[
\texttt{certificates/orientation/rhomred\_orientation\_thom\_sebastiani\_compatibility}
\]
verifies
\(\texttt{RHOMRED\_ORIENTATION\_THOM\_SEBASTIANI\_COMPATIBILITY\_OBSTRUCTION\_VERIFIED}\):
it imports the row \(300\) extension-multiplicativity packet, the
reduced-orientation ledger, the mixed Thom--Sebastiani packet, and the
four-input pentagon packet; it records that the BBDJS
Thom--Sebastiani rows and conditional pentagon rows do not construct
the missing orientation-line compatibility row.

\begin{proposition}[Orientation descent through the \(E\)-quotient]
\label{prop:orientation-e-quotient-descent}
\textnormal{[descent criterion proved; quotient-orientation rows not
supplied]}
Let \(o_{\mathcal C}\) be a reduced orientation line on a retained
stratum of the \(K3\times E\) reduced moduli stack.  The line descends
through the quotient by \(E\) precisely when the following finite
Picard-groupoid datum is supplied.
\begin{enumerate}
\item a square-root row for \(o_{\mathcal C}^{\otimes2}\simeq
\det K^{\mathrm{red}}_{\mathcal C}\);
\item a null-trivialization of the ordinary reduced gerbe
\(\alpha^{\mathrm{red}}_{\mathcal C}\);
\item a Cech--Borel representative of the connected free \(E\)-class
\[
\alpha^{E,\mathrm{free}}_{\mathcal C}
=a_1(\mathcal C)u_1+a_2(\mathcal C)u_2\in H^2(BE;\mathbb F_2)
\]
with \(a_1=a_2=0\) and a chosen null-trivialization;
\item for every retained finite-stabilizer stratum \(S\), an
edge-reduced equivariant class
\[
\widetilde\beta^H_{\mathcal C,S}\rightsquigarrow
\beta^H_{\mathcal C,S}\in H^2(BH;\mathbb F_2),
\]
a zero row for \(\beta^H_{\mathcal C,S}\), and compatible
null-trivialisations under retained subgroup restrictions;
\item a computed residual linearization character
\[
\lambda^H_{\mathcal C,S}\in H^1(BH;\mathbb F_2)
\]
with \(\lambda^H_{\mathcal C,S}=0\) for every retained stabilizer;
\item compatibility of these choices on overlaps and on the quotient
presentation of the retained extension and flag charts.
\end{enumerate}
Equivalently, row \(302\) is proved by a populated
\(\mathrm{quotient\_borel}\) table and a populated
\(\mathrm{finite\_stabilizers}\) table whose obstruction ranks and
degree-one characters vanish.  The descent object is the descended
line \(o_{\mathcal C}/E\) on the quotient stack together with the
chosen quotient cocycle
\[
\mathfrak q_{\mathcal C}=
(\alpha^{\mathrm{red}}_{\mathcal C},
\alpha^{E,\mathrm{free}}_{\mathcal C},
\{(\beta^H_{\mathcal C,S},\lambda^H_{\mathcal C,S})\}_{S,H})
\simeq0.
\]

The present packets do not supply this datum.  They record criteria
for the reduced gerbe, connected free \(E\)-class, finite-stabilizer
classes, and residual linearization characters, but
\(\mathrm{orientation\_lines}\), \(\mathrm{quotient\_borel}\), and
\(\mathrm{finite\_stabilizers}\) are empty.  Thus row \(302\) is
currently a descent criterion with an open quotient-orientation
obstruction, not a descended orientation.
\end{proposition}

\begin{proof}
For a line on an \(E\)-stack to descend to the quotient stack, its
equivariant gerbe and stabilizer linearizations must be trivialized.
In the reduced orientation problem the ordinary component is
\(\alpha^{\mathrm{red}}_{\mathcal C}\).  The connected part of the
\(E\)-action contributes the \(H^2(BE;\mathbb F_2)\) edge class
\(\alpha^{E,\mathrm{free}}_{\mathcal C}\).  Finite stabilizers
contribute the Borel classes
\(\widetilde\beta^H_{\mathcal C,S}\), whose classifying-space edges
\(\beta^H_{\mathcal C,S}\) are legitimate targets only after the
mixed Borel and differential terms have been killed.  After the
degree-two gerbes vanish, a residual degree-one stabilizer character
\(\lambda^H_{\mathcal C,S}\) remains; nonzero \(\lambda^H\) changes
the descended orientation by a sign character and therefore prevents
descent.

The six listed rows are exactly the data required to trivialize these
obstructions and make the equivariant orientation line descend.  The
compatibility clauses are needed because descent is a statement in the
Picard groupoid on the quotient presentation, not a collection of
unrelated pointwise vanishings.  The current reduced-orientation
fixture records all quotient-Borel and finite-stabilizer rows as
missing open obligations, and the specialized packets for
\(\alpha^{\mathrm{red}}\), \(\alpha^{E,\mathrm{free}}\),
\(\beta^H\), and \(\lambda^H\) are obstruction ledgers.  Hence the
criterion is proved, while the descended orientation is not
constructed.
\end{proof}

The packet
\[
\texttt{certificates/orientation/rhomred\_orientation\_e\_quotient\_descent}
\]
verifies
\(\texttt{RHOMRED\_ORIENTATION\_E\_QUOTIENT\_DESCENT\_OBSTRUCTION\_VERIFIED}\):
it imports the square-root, reduced-gerbe nullhomotopy, connected
free \(E\)-nulltrivialization, finite-stabilizer
nulltrivialization, finite-stabilizer linearization-vanishing, and
reduced-orientation ledger packets; it checks that the quotient
tables are empty and records the exact missing descent rows.

\begin{proposition}[Orientation descent and the Hall product]
\label{prop:orientation-descent-hall-product}
\textnormal{[product-descent criterion proved; Hall-product orientation
rows not supplied]}
Fix a finite height \(R\) and a retained compactified extension
correspondence
\[
\mathfrak M_{R,\widehat c}\times\mathfrak M_{R,\widehat c'}
\xleftarrow{\ p\ }
\overline{\mathfrak E}_{R,\widehat c,\widehat c'}
\xrightarrow{\ q\ }
\mathfrak M_{R,\widehat c+\widehat c'} .
\]
Put \(K_{R,\widehat c}:=\Phi^{\mathrm{red}}_{R,\widehat c}\otimes
o_{R,\widehat c}\).  The statement that quotient-orientation descent
commutes with the Hall product is the commutative square
\[
\begin{tikzcd}
Q_{E,R}\,q_!\operatorname{TS}^{\Phi\otimes o}_R p^*
\ar[r,"\theta^Q_{\mu,R}"]
\ar[d,"Q_{E,R}m_R"']&
\overline q_!\,
\overline{\operatorname{TS}}^{\Phi\otimes o}_R
\overline p^*(Q_{E,R}\boxtimes Q_{E,R})
\ar[d,"\overline m_R"]\\
Q_{E,R}K_{R,\widehat c+\widehat c'}
\ar[r,"\theta^Q_{K,R}"']&
\overline K_{R,\widehat c+\widehat c'}
\end{tikzcd}
\]
in the Borel--Moore correspondence category on the quotient stack,
where \(\theta^Q_{\mu,R}\) is the quotient-after-correspondence
operation comparison and
\(\theta^Q_{K,R}\colon Q_{E,R}K_{R,\widehat c+\widehat c'}
\simeq\overline K_{R,\widehat c+\widehat c'}\) is the descended
coefficient-system identification.

This square is proved by supplying the following finite datum.
\begin{enumerate}
\item quotient-orientation rows for the two input strata, the
extension stack, and the output stratum, with the quotient cocycles
\[
p^*(\mathfrak q_{R,\widehat c}\boxtimes
\mathfrak q_{R,\widehat c'})
\quad\text{and}\quad
q^*\mathfrak q_{R,\widehat c+\widehat c'}
\]
identified by a zero Picard-groupoid defect row;
\item the oriented Thom--Sebastiani row
\[
\operatorname{TS}^{\Phi\otimes o}_R
=\operatorname{TS}^{\Phi,\mathrm{red}}_R\otimes
\operatorname{TS}^{o}_R
\]
whose square is the reduced determinant Thom--Sebastiani comparison
and whose quotient-descent defect is zero;
\item external-product descent
\[
Q_{E,R}(K_{R,\widehat c}\boxtimes K_{R,\widehat c'})
\simeq
\overline K_{R,\widehat c}\boxtimes\overline K_{R,\widehat c'}
\]
on the product quotient;
\item compact-support pushforward descent for \(q_!\), with proper
base-change and projection-formula defect ranks zero;
\item the quotient-after-correspondence comparison
\(\theta^Q_{\mu,R}\) for the supplied binary Hall product row;
\item compatibility of the preceding rows with the two-step flag
pentagon, so that product descent is associative after quotient.
\end{enumerate}
The current packets do not supply this datum.  Row \(302\) supplies no
quotient-orientation rows, row \(301\) supplies no oriented
Thom--Sebastiani row, and the compact Hall source packet supplies no
geometric Hall product matrix.  The
quotient-after-correspondence packet constructs \(\theta^Q_{\mu,R}\)
only for operation rows with descent interfaces already supplied.  Thus
row \(303\) is presently a product-descent criterion, not a proof that
orientation descent commutes with the Hall product.
\end{proposition}

\begin{proof}
The Hall product is the pull--transport--push composite
\[
m_R=q_!\operatorname{TS}^{\Phi\otimes o}_R p^* .
\]
A quotient comparison for this composite is an isomorphism of
composites, not a scalar identity.  Therefore it decomposes into three
functorial requirements: descent of \(p^*\) and of the external product
coefficient system, descent of the oriented Thom--Sebastiani
transport, and descent of \(q_!\) through compact supports.  The first
requirement is exactly the Picard-groupoid equality of the pulled-back
quotient cocycles on the extension correspondence.  The second is the
orientation-line row of Proposition~\ref{prop:orientation-thom-sebastiani-compatibility}
tensored with the reduced BBDJS Thom--Sebastiani transport.  The third
is the compact-support pushforward comparison used by
Proposition~\ref{prop:quotient-after-correspondence-theta-mu-comparison}.

If these rows are supplied, the naturality of pullback, the descended
Thom--Sebastiani isomorphism, and compact-support pushforward descent
identify
\[
Q_{E,R}q_!\operatorname{TS}^{\Phi\otimes o}_R p^*
\quad\text{with}\quad
\overline q_!\overline{\operatorname{TS}}^{\Phi\otimes o}_R
\overline p^*(Q_{E,R}\boxtimes Q_{E,R}),
\]
which is the displayed square.  Associativity after quotient is the
same statement on the retained two-step flag stack.  The present
tables fail before this proof can be applied: the quotient-orientation
tables are empty, the oriented Thom--Sebastiani row is absent, and the
source Hall product rows are not populated.  Hence only the criterion
is proved.
\end{proof}

The packet
\[
\texttt{certificates/orientation/rhomred\_orientation\_hall\_product\_descent}
\]
verifies
\(\texttt{RHOMRED\_ORIENTATION\_HALL\_PRODUCT\_DESCENT\_OBSTRUCTION\_VERIFIED}\):
it imports the row \(302\) quotient-orientation descent packet, the
row \(301\) oriented Thom--Sebastiani packet, the
quotient-after-correspondence \(\theta^Q_{\mu}\) packet, and the
transition Hall-product obstruction packet; it records that the current
tree has comparison maps for already supplied operation rows but no
quotient-orientation Hall-product descent row.

\begin{proposition}[Orientation descent and the Hall coproduct]
\label{prop:orientation-descent-hall-coproduct}
\textnormal{[coproduct-descent criterion proved; Hall-coproduct
orientation rows not supplied]}
Fix a finite height \(R\) and the same retained compactified
correspondence, read in the splitting direction
\[
\mathfrak M_{R,\widehat c+\widehat c'}
\xleftarrow{\ q\ }
\overline{\mathfrak E}_{R,\widehat c,\widehat c'}
\xrightarrow{\ p\ }
\mathfrak M_{R,\widehat c}\times\mathfrak M_{R,\widehat c'} .
\]
With \(K_{R,\widehat c}=\Phi^{\mathrm{red}}_{R,\widehat c}\otimes
o_{R,\widehat c}\), quotient-orientation descent commutes with the
Hall coproduct precisely when the pull--inverse-transport--push
comparison
\[
\begin{tikzcd}
Q_{E,R}\,p_!(\operatorname{TS}^{\Phi\otimes o}_R)^{-1}q^*
\ar[r,"\theta^Q_{\Delta,R}"]
\ar[d,"Q_{E,R}\Delta_R"']&
\overline p_!\,
(\overline{\operatorname{TS}}^{\Phi\otimes o}_R)^{-1}
\overline q^*Q_{E,R}
\ar[d,"\overline\Delta_R"]\\
Q_{E,R}(K_{R,\widehat c}\boxtimes K_{R,\widehat c'})
\ar[r,"\theta^Q_{K\boxtimes K,R}"']&
\overline K_{R,\widehat c}\boxtimes
\overline K_{R,\widehat c'}
\end{tikzcd}
\]
commutes on the quotient Borel--Moore correspondence category.

This square is proved by supplying the following finite datum.
\begin{enumerate}
\item quotient-orientation rows for the source stratum, the splitting
stack, and the two output strata, with the cocycle comparison
\[
q^*\mathfrak q_{R,\widehat c+\widehat c'}
=
p^*(\mathfrak q_{R,\widehat c}\boxtimes
\mathfrak q_{R,\widehat c'})
\]
in the Picard groupoid;
\item the inverse oriented Thom--Sebastiani row
\[
(\operatorname{TS}^{\Phi\otimes o}_R)^{-1}
=
(\operatorname{TS}^{\Phi,\mathrm{red}}_R)^{-1}\otimes
(\operatorname{TS}^{o}_R)^{-1}
\]
with determinant-square and quotient-descent defects zero;
\item external-product descent for the two-output coefficient system
\(K_{R,\widehat c}\boxtimes K_{R,\widehat c'}\);
\item compact-support pushforward descent for \(p_!\), with proper
base-change and projection-formula defect ranks zero;
\item a supplied collision-coproduct row \(\Delta_R\), finite
coproduct matrix \(D_R\), counit row, and zero coassociativity and
counit defects;
\item the quotient-after-correspondence comparison
\(\theta^Q_{\Delta,R}\) and compatibility with the bialgebra square
relating \(\Delta_R\) and \(m_R\).
\end{enumerate}
The current packets do not supply this datum.  Row \(302\) supplies no
quotient-orientation rows, row \(301\) supplies no oriented
Thom--Sebastiani row, the compact Hall source packet supplies no
\(D\)-entries or counit rows, and the
quotient-after-correspondence coproduct packet applies only to already
supplied quotient-presented coproduct and bialgebra rows.  Thus row
\(304\) is presently a coproduct-descent criterion, not a proof that
orientation descent commutes with the Hall coproduct.
\end{proposition}

\begin{proof}
The collision coproduct is the pull--inverse-transport--push composite
\[
\Delta_R=p_!(\operatorname{TS}^{\Phi\otimes o}_R)^{-1}q^* .
\]
For this composite to descend, the quotient cocycle on the source
object must agree, after pullback to the splitting correspondence, with
the external product of the two output quotient cocycles.  This is the
Picard-groupoid equality displayed above.  The coefficient transport is
then the inverse of the oriented Thom--Sebastiani isomorphism, not the
product orientation row alone.  Finally the exceptional pushforward is
along \(p\), so the six-functor comparison is the \(p_!\)-descent row
for the splitting correspondence.

Given those rows, functoriality of pullback, inverse
Thom--Sebastiani descent, and compact-support descent identify
\[
Q_{E,R}p_!(\operatorname{TS}^{\Phi\otimes o}_R)^{-1}q^*
\quad\text{with}\quad
\overline p_!
(\overline{\operatorname{TS}}^{\Phi\otimes o}_R)^{-1}
\overline q^*Q_{E,R},
\]
which is the displayed square.  Coassociativity and counitality after
quotient are the same statement on the retained two-step flag and
zero-object rows.  The present tables fail before the criterion can be
applied: the quotient-orientation tables are empty, the oriented
Thom--Sebastiani row is absent, and the source coproduct and counit
rows are not populated.  Hence only the criterion is proved.
\end{proof}

The packet
\[
\texttt{certificates/orientation/rhomred\_orientation\_hall\_coproduct\_descent}
\]
verifies
\(\texttt{RHOMRED\_ORIENTATION\_HALL\_COPRODUCT\_DESCENT\_OBSTRUCTION\_VERIFIED}\):
it imports the row \(302\) quotient-orientation descent packet, the
row \(301\) oriented Thom--Sebastiani packet, the row \(303\)
product-descent packet, the quotient-after-correspondence
\(\theta^Q_{\Delta}\) packet, and the transition Hall-coproduct
obstruction packet; it records that the current tree has conditional
coproduct comparison rows but no quotient-orientation Hall-coproduct
descent row.

\begin{proposition}[Orientation descent and primitive projection]
\label{prop:orientation-descent-primitive-projection}
\textnormal{[primitive-projection descent criterion proved; primitive
orientation rows not supplied]}
Fix \(R\).  Let
\[
K_R=\Phi^{\mathrm{red}}_R\otimes o_R
\]
be the oriented vanishing-cycle coefficient system on the retained
finite Hall object \(\mathcal H_R(K_R)\).  Write
\[
\widetilde\Delta_R=\Delta_R-\eta_R\epsilon_R\boxtimes\mathrm{id}
-\mathrm{id}\boxtimes\eta_R\epsilon_R,\qquad
P_R(K_R)=\ker\widetilde\Delta_R\cap\ker\epsilon_R
\]
and let
\(\pi_R^{\Prim}\colon\mathcal H_R(K_R)\to P_R(K_R)\)
be the supplied primitive idempotent.  Denote the quotient objects by
\(\overline{\mathcal H}_R(\overline K_R)\),
\(\overline P_R(\overline K_R)\), and
\(\overline\pi_R^{\Prim}\).  Orientation descent commutes with
primitive projection precisely when the idempotent square
\[
\begin{tikzcd}
Q_{E,R}\mathcal H_R(K_R)
\ar[r,"\theta^Q_{H,R}"]
\ar[d,"Q_{E,R}\pi_R^{\Prim}"']&
\overline{\mathcal H}_R(\overline K_R)
\ar[d,"\overline\pi_R^{\Prim}"]\\
Q_{E,R}P_R(K_R)
\ar[r,"\theta^{Q,\Prim}_{o,R}"']&
\overline P_R(\overline K_R)
\end{tikzcd}
\]
commutes in the idempotent-complete quotient Borel--Moore
correspondence category.

This square is proved by supplying the following finite datum.
\begin{enumerate}
\item quotient-orientation rows for the ambient Hall object, the
zero-object stratum, and every splitting correspondence entering
\(\widetilde\Delta_R\);
\item the Hall-coproduct orientation-descent datum of
Proposition~\ref{prop:orientation-descent-hall-coproduct}, together
with unit and counit orientation descent;
\item finite source and quotient primitive kernel rows
\[
P_R(K_R)=\ker(\widetilde\Delta_R,\epsilon_R),\qquad
\overline P_R(\overline K_R)
=\ker(\widetilde{\overline\Delta}_R,\overline\epsilon_R),
\]
with inclusion and projection matrices for the idempotents
\(\pi_R^{\Prim}\) and \(\overline\pi_R^{\Prim}\);
\item the quotient-after-correspondence primitive comparison
\[
\theta^{Q,\Prim}_{R}\colon Q_{E,R}P_R(K_R)
\xrightarrow{\sim}\overline P_R(\overline K_R)
\]
with zero reduced-coproduct and counit kernel-intertwining defects;
\item the projection-commutator identity
\[
\overline\pi_R^{\Prim}\theta^Q_{H,R}
=\theta^{Q,\Prim}_{o,R}Q_{E,R}\pi_R^{\Prim};
\]
\item for the pro-object, transition-compatible primitive matrices and
zero composition defects for \(R'\to R\).
\end{enumerate}
The present packets do not supply this datum.  The quotient
primitive-projection packet proves compatibility of
\(\theta^Q_{\mu,R}\) only for already supplied finite primitive kernel
and projection rows.  The compact Hall source packet supplies no
\(D\)-entries, counit rows, primitive kernels, or primitive projection
matrices; the transition primitive-subspace packet records empty
transition tables.  Hence row \(305\) is a primitive-projection descent
criterion, not a proof that orientation descent commutes with primitive
projection.
\end{proposition}

\begin{proof}
The primitive subobject is the kernel of the map
\[
(\widetilde\Delta_R,\epsilon_R)\colon
\mathcal H_R(K_R)\longrightarrow
\mathcal H_R(K_R)\boxtimes\mathcal H_R(K_R)\oplus\mathbf 1 .
\]
If the coproduct, unit, counit, and orientation coefficient systems
descend through \(Q_{E,R}\), then applying \(Q_{E,R}\) carries this
kernel diagram to the quotient kernel diagram.  The supplied kernel
rows make this statement finite: the reduced-coproduct and counit
matrices intertwine with the quotient comparison, so
\(Q_{E,R}P_R(K_R)\) identifies with
\(\overline P_R(\overline K_R)\).

The primitive projection is then an idempotent splitting of this kernel
inside the idempotent-complete correspondence category.  Two morphisms
from \(Q_{E,R}\mathcal H_R(K_R)\) to
\(\overline P_R(\overline K_R)\) are equal once their compositions with
the inclusion into \(\overline{\mathcal H}_R(\overline K_R)\) agree.
That equality is exactly the kernel-intertwining identity for
\(\widetilde\Delta_R\) and \(\epsilon_R\), together with the
projection-commutator row.  Transition compatibility is a separate
strictness condition needed only after passing from one finite height to
the pro-object.

The available data stop before this criterion can be applied.  Row
\(304\) does not provide an actual oriented Hall-coproduct descent row,
the compact source has no coproduct or counit matrices, and the
primitive transition packet has no primitive kernel or projection
matrices.  Thus only the criterion is proved.
\end{proof}

The packet
\[
\texttt{certificates/orientation/rhomred\_orientation\_primitive\_projection\_descent}
\]
verifies
\(\texttt{RHOMRED\_ORIENTATION\_PRIMITIVE\_PROJECTION\_DESCENT\_OBSTRUCTION\_VERIFIED}\):
it imports row \(302\), row \(304\), the
quotient-after-correspondence primitive-projection packet, the
transition primitive-subspace obstruction packet, and the compact Hall
source ledger; it records that the present tree has a conditional
\(\theta^{Q,\Prim}\) comparison but no compact Hall primitive matrices
and no orientation primitive-projection descent row.

\subsubsection*{(O1)$^+$ Weyl-equivariant transport along
$W^{(2)}(\Lambda_{II}^{2,1})$}

The previous criterion isolates the data needed for descent through
the reduced $E$-quotient.  Once such quotient-orientation data are
supplied, the Igusa structure still requires compatibility with the
Weyl group of the BKM target.  For each of the three type-II simple
reflections $s_{\delta_i}$ ($i=1,2,3$), the $K3\times E$ side carries
a wall transport
\[
\tau_i:o_{\mathcal C}\longrightarrow s_{\delta_i}^{*}o_{\mathcal C}.
\]
The word \emph{coherent} is a hypothesis.  At finite height, the
first object is the finite partial action groupoid $\mathcal G_R$
whose objects are retained oriented quotient strata and whose arrows
are the retained type-II wall transports; a choice of
determinant-line lifts gives a $2$-cocycle
$c_{o,R}\in Z^2(\mathcal G_R;\underline{\mathbb F}_2)$.  Only on a
complete $W^{(2)}(\Lambda_{II}^{2,1})$-orbit with constant
coefficients does this reduce to the group cohomology class
\[
[c_o]\in H^2\bigl(W^{(2)}(\Lambda_{II}^{2,1});\mathbb F_2\bigr).
\]

\begin{definition}[(O1)$^+$ Weyl-equivariant determinant-line
transport]
\label{def:O1-plus-weyl-transport}
The (O1)$^+$ datum is the chosen set of lifts $\{\tau_i\}$
satisfying $[c_o]=0$ with chosen $1$-cochain killing $c_o$, the
normalization $\tau_i^2=1$, and the transport of the
quotient-orientation cocycle
\[
\mathfrak q_{\mathcal C}=\bigl(\alpha^{\mathrm{red}}_{\mathcal C},
\alpha^{E,\mathrm{free}}_{\mathcal C},
\{(\beta^H_{\mathcal C},\lambda^H_{\mathcal C})\}_{H\in
\mathcal H(\mathcal C)}\bigr)
\]
along $\tau_i$ between charge strata over Weyl-related Gram degrees.
The torsor defect of this transport is
\[
\omega_{i,\mathcal C}\in
H^1(\mathfrak M^{\mathrm{red}}_{s_{\delta_i}\mathcal C};\mathbb F_2)
\oplus
\bigoplus_{H\in\mathcal H(\mathcal C)}H^1(BH;\mathbb F_2),
\]
required to vanish.  The Coxeter presentation of
$W^{(2)}(\Lambda_{II}^{2,1})$ then assembles these lifts into an honest
action on the oriented Pfaffian line over the wall-complement of the
reduced compact Hall sector.
\end{definition}

(O1)$^+$ is a determinant-line monodromy datum.  No scalar formula
enters; the transports must move the entire null-trivialization
datum, not only the underlying line.  The (O1)$^+$ data comprise
$[c_o]=0$, the chosen cochain, the unitarity $\tau_i^2=1$, and the
torsor vanishing $\omega_{i,\mathcal C}=0$ on every retained orbit.
In finite coordinates this is recorded by
\[
\mathrm{weyl\_lifts},\quad
\mathrm{coxeter\_cocycles},\quad
\mathrm{transitions},\quad
\mathrm{scalar\_firewall}.
\]
The Weyl-lift rows carry \(\tau_i^2=1\), torsor-defect zero, and
quotient-cocycle transport zero.  The Coxeter rows carry the vanishing
projective cocycle, the cochain killing it, and the type-II Coxeter
defect.  Transition rows carry strictness and \(R^1\varprojlim=0\) for
orientation lines, quotient data, and Weyl lifts.  The scalar firewall
excludes scalar trace, squared determinant, OP scalar branch, and
Maass-character values as orientation data.

\begin{proposition}[Weyl wall transport maps]
\label{prop:weyl-wall-transport-tau-construction}
\textnormal{[wall-transport construction criterion proved; \(\tau_i\)
rows not supplied]}
Fix \(R\), a retained quotient stratum \(S\), and a type-II simple
reflection \(s_{\delta_i}\), \(i\in\{1,2,3\}\).  A Weyl wall transport
\[
\tau_{i,R,S}\colon o_{R,S}\longrightarrow
s_{\delta_i}^{*}o_{R,s_{\delta_i}S}
\]
is constructed by a finite \(\texttt{weyl\_lifts}\) row with the
following payload.
\begin{enumerate}
\item a verified type-II generator row for \(s_{\delta_i}\) and a
retained wall correspondence carrying \(S\) to
\(s_{\delta_i}S\);
\item source and target Joyce--Upmeier orientation lines
\[
o_{R,S}^{\otimes2}\simeq\det K^{\mathrm{red}}_{R,S},
\qquad
o_{R,s_{\delta_i}S}^{\otimes2}
\simeq\det K^{\mathrm{red}}_{R,s_{\delta_i}S};
\]
\item a Picard-groupoid lift matrix for \(\tau_{i,R,S}\), whose square
with the determinant transport commutes:
\[
(\tau_{i,R,S})^{\otimes2}\colon
o_{R,S}^{\otimes2}\longrightarrow
s_{\delta_i}^{*}o_{R,s_{\delta_i}S}^{\otimes2}
\]
agrees with the determinant transport
\(\det K^{\mathrm{red}}_{R,S}\to
s_{\delta_i}^{*}\det K^{\mathrm{red}}_{R,s_{\delta_i}S}\);
\item transport of the quotient-orientation cocycle
\(\mathfrak q_{R,S}\) to
\(s_{\delta_i}^{*}\mathfrak q_{R,s_{\delta_i}S}\), including the
finite-stabilizer Borel classes and the zero linearization
characters;
\item vanishing torsor and quotient-cocycle transport defects
\[
\omega_{i,R,S}=0,\qquad
\operatorname{defect}_{q,i,R,S}=0 .
\]
\end{enumerate}
The condition \(\tau_{i,R,S}^2=1\) is a separate involutivity row; it
is not part of this construction criterion.

The current data prove only the target and profile prerequisites: the
type-II chamber packet supplies the three simple reflections, and the
quotient-after-correspondence wall-profile packet proves that the three
profiles survive the \(E\)-quotient.  The reduced-orientation packet
has no orientation-line rows and no \(\texttt{weyl\_lifts}\) rows, and
the wall-atlas packet has no retained wall-object rows.  Thus row
\(306\) is presently a Weyl-wall transport criterion, not a
construction of the maps \(\tau_i\).
\end{proposition}

\begin{proof}
The map \(\tau_{i,R,S}\) is a Picard-groupoid isomorphism, not a scalar
sign.  Hence the first input is an actual source and target orientation
line over Weyl-related quotient strata.  The determinant square-root
condition forces the displayed square to commute, so the lift preserves
the Joyce--Upmeier square root of the reduced determinant complex.
The quotient-orientation datum contains the reduced gerbe, the free
\(E\)-Borel class, the finite-stabilizer Borel classes, and the
linearization characters; transporting only the underlying line would
lose this data.  The vanishing of \(\omega_{i,R,S}\) and of the
quotient-cocycle defect is exactly the assertion that the entire
quotient-orientation package moves across the wall.

The type-II generator and profile-survival rows do not construct this
Picard isomorphism.  They identify which reflections and quotient
profiles must be used.  The construction itself starts only when the
orientation-line rows, retained wall-object rows, and Weyl-lift rows
are populated.  Those tables are empty in the current packets, so the
criterion is the strongest proved statement.
\end{proof}

The packet
\[
\texttt{certificates/orientation/rhomred\_weyl\_wall\_transport\_tau}
\]
verifies
\(\texttt{RHOMRED\_WEYL\_WALL\_TRANSPORT\_TAU\_OBSTRUCTION\_VERIFIED}\):
it imports the blocked reduced-orientation ledger, the type-II chamber
generator packet, the quotient wall-profile survival packet, and the
blocked wall-atlas ledger; it records that the three target reflections
and quotient profiles are available, but no source orientation-line,
wall-object, or \(\tau_i\)-transport rows are supplied.

\begin{proposition}[Involutivity of the Weyl wall transport]
\label{prop:weyl-wall-transport-tau-square}
\textnormal{[involutivity criterion proved; \(\tau_i^2\) rows not
supplied]}
Fix \(R\), a retained quotient stratum \(S\), and a supplied Weyl wall
transport
\[
\tau_{i,R,S}\colon o_{R,S}\longrightarrow
s_{\delta_i}^{*}o_{R,s_{\delta_i}S}.
\]
The identity \(\tau_{i,R,S}^2=1\) is the commutativity of the
Picard-groupoid square
\[
\begin{tikzcd}
o_{R,S}\ar[r,"\tau_{i,R,S}"]\ar[dr,"\mathrm{id}"']&
s_{\delta_i}^{*}o_{R,s_{\delta_i}S}
\ar[d,"s_{\delta_i}^{*}\tau_{i,R,s_{\delta_i}S}"]\\
&s_{\delta_i}^{*}s_{\delta_i}^{*}o_{R,S}
\ar[d,"\kappa_{i,R,S}"]\\
&o_{R,S}
\end{tikzcd}
\]
where \(\kappa_{i,R,S}\) is the supplied identification of the
double-reflected orientation line with \(o_{R,S}\) over the equality
\(s_{\delta_i}^2=1\) on the retained wall correspondence.

It is proved by a finite \(\texttt{weyl\_lifts}\) involutivity row with
the following data:
\begin{enumerate}
\item a row constructing both lifts
\(\tau_{i,R,S}\) and
\(\tau_{i,R,s_{\delta_i}S}\) as in
Proposition~\ref{prop:weyl-wall-transport-tau-construction};
\item the target Coxeter row \(s_{\delta_i}^2=1\) for the type-II
reflection, and the corresponding double-wall source correspondence;
\item a determinant-square compatibility row showing that
\((s_{\delta_i}^{*}\tau_{i,R,s_{\delta_i}S})\circ\tau_{i,R,S}\)
squares to the identity transport of
\(\det K^{\mathrm{red}}_{R,S}\);
\item transport of the quotient-orientation cocycle around the
two-edge loop with zero reduced-gerbe, free-\(E\), finite-stabilizer,
and linearization defects;
\item a finite matrix row
\[
\operatorname{defect}_{\tau^2,i,R,S}=0
\]
recording that the composite equals \(\mathrm{id}_{o_{R,S}}\), not
only its square on the determinant line.
\end{enumerate}

The present data do not supply these rows.  Row \(306\) records only
the criterion for constructing \(\tau_i\).  The type-II chamber packet
records the target relation \(s_{\delta_i}^2=1\), and the Maass
character packet records
\(\nu_{\Delta_5}(s_{\delta_i})^2=1\); neither constructs an
orientation-line lift nor proves its square is the identity.  Thus row
\(307\) is presently an involutivity criterion, not a proof of
\(\tau_i^2=1\).
\end{proposition}

\begin{proof}
The reflection identity \(s_{\delta_i}^2=1\) is an identity in the
target type-II Coxeter group.  The lift \(\tau_{i,R,S}\) lives in a
Picard groupoid over the retained source stratum.  Therefore
involutivity is a statement about the two-edge source loop
\[
S\longrightarrow s_{\delta_i}S\longrightarrow S
\]
and about the induced automorphism of the orientation line over that
loop.

The determinant-square row removes the first possible defect: after
squaring, the composite lift acts trivially on the determinant complex.
This is not enough, because a square root can still differ from the
identity by the residual sign automorphism of the orientation line and
by the quotient-orientation torsor.  The quotient-cocycle transport
rows remove the Borel and finite-stabilizer defects.  The final
\(\tau^2\)-defect row asserts that the remaining line automorphism is
the identity.  These are finite matrix and Picard-groupoid checks, not
consequences of the scalar character value.

Since the current \(\texttt{weyl\_lifts}\) table is empty, none of the
composites in the displayed square is supplied.  The available target
and automorphic rows identify the relation to be lifted; they do not
lift it.  Hence only the criterion is proved.
\end{proof}

The packet
\[
\texttt{certificates/orientation/rhomred\_weyl\_tau\_square}
\]
verifies
\(\texttt{RHOMRED\_WEYL\_TAU\_SQUARE\_OBSTRUCTION\_VERIFIED}\):
it imports the row \(306\) transport criterion, the blocked
reduced-orientation ledger, the type-II chamber packet, and the Maass
character packet; it records that the target reflection relation and
the scalar square value are available, but no \(\tau_i\)-lift or
\(\tau_i^2\)-defect row is supplied.

\begin{definition}[Finite partial action groupoid]
\label{def:finite-partial-action-groupoid-GR}
\textnormal{[groupoid datum specified; finite groupoid rows not
supplied]}
At height \(R\), the finite partial action groupoid
\(\mathcal G_R\) is the following finite datum.
\begin{enumerate}
\item The objects are retained quotient-orientation strata \(S\) for
which the orientation line \(o_{R,S}\), quotient cocycle
\(\mathfrak q_{R,S}\), and retained support condition are supplied.
\item The generating arrows are the retained type-II wall transports
\[
g_{i,R,S}\colon S\longrightarrow s_{\delta_i}S
\]
for \(i\in\{1,2,3\}\), defined only when \(S\) and
\(s_{\delta_i}S\) both lie in the retained finite window and the row
\(\tau_{i,R,S}\) of
Proposition~\ref{prop:weyl-wall-transport-tau-construction} is
supplied.
\item The identity arrow \(\mathrm{id}_S\), inverse arrow
\[
g_{i,R,S}^{-1}=g_{i,R,s_{\delta_i}S},
\]
and partial compositions
\[
g_{i_k,R,S_{k-1}}\circ\cdots\circ g_{i_1,R,S_0}
\]
are included exactly when every intermediate stratum \(S_j\) is
retained and every wall-transport row in the word is supplied.
\item The structure tables record source, target, identity, inverse,
and associativity defects; all are required to have defect rank zero.
\end{enumerate}
The coefficient system for the orientation cocycle on
\(\mathcal G_R\) is not forced to be constant.  It is the local system
whose fibre at \(S\) is the automorphism group of the complete
quotient-orientation package over \(S\).  Only on a complete retained
orbit with identified fibres may one replace it by the constant
\(\underline{\mathbb F}_2\)-system.

The present data do not populate \(\mathcal G_R\).  The type-II
chamber packet gives the three target generators and the no-finite-braid
relation among distinct generators, and the quotient wall-profile
packet gives three surviving profiles.  The reduced-orientation packet
has no retained object rows, no Weyl-lift rows, and no
\(\texttt{coxeter\_cocycles}\) rows.  Thus row \(308\) specifies the
finite partial action groupoid datum, but the current tree contains no
finite \(\mathcal G_R\) object, arrow, inverse, or composition tables.
\end{definition}

The packet
\[
\texttt{certificates/orientation/rhomred\_weyl\_partial\_action\_groupoid}
\]
verifies
\(\texttt{RHOMRED\_WEYL\_PARTIAL\_ACTION\_GROUPOID\_OBSTRUCTION\_VERIFIED}\):
it imports the target type-II chamber packet, the row \(306\) and row
\(307\) Weyl-lift criteria, the quotient wall-profile survival packet,
and the blocked reduced-orientation ledger; it records that the target
generators are available but the finite groupoid object, arrow,
inverse, and composition rows are not supplied.

\begin{definition}[Orientation projective cocycle]
\label{def:orientation-projective-cocycle-coR}
\textnormal{[cocycle formula specified; finite \(c_{o,R}\) rows not
supplied]}
Assume the finite partial action groupoid \(\mathcal G_R\) of
Definition~\ref{def:finite-partial-action-groupoid-GR} has been
populated, and that every arrow \(g:S\to T\) in \(\mathcal G_R\) is
equipped with a determinant-line lift
\[
\tau_g:o_{R,S}\longrightarrow g^*o_{R,T}
\]
transporting the full quotient-orientation package.  For every
composable pair
\[
S\xrightarrow{g}T\xrightarrow{h}U
\]
with composite \(hg\), the orientation projective cocycle is the
unique element
\[
c_{o,R}(h,g;S)\in
\operatorname{Aut}_{\mathfrak Q^{\mathrm{or}}_{R,S}}(o_{R,S})
\simeq \mathbb F_2
\]
defined by
\[
g^*\tau_h\circ\tau_g
=(-1)^{c_{o,R}(h,g;S)}\,\tau_{hg}.
\]
Equivalently, \(c_{o,R}\) is the finite table measuring the defect
between composition of chosen Weyl determinant-line lifts and the
chosen lift of the partial product in \(\mathcal G_R\).

The row is valid only after the following data are supplied:
\begin{enumerate}
\item object, arrow, inverse, and composition rows for
\(\mathcal G_R\);
\item determinant-line lifts for \(g\), \(h\), and \(hg\), including
the quotient-cocycle transport rows;
\item an identified coefficient fibre
\(\operatorname{Aut}_{\mathfrak Q^{\mathrm{or}}_{R,S}}(o_{R,S})\) at
the source object;
\item a cocycle matrix row recording \(c_{o,R}(h,g;S)\) for every
composable pair;
\item a normalization row
\[
c_{o,R}(g,\mathrm{id}_S;S)=
c_{o,R}(\mathrm{id}_T,g;S)=0;
\]
\item a cocycle-defect row
\[
dc_{o,R}(k,h,g;S)=0
\]
for every composable triple in the partial groupoid.
\end{enumerate}
The present packets do not provide these rows.  The row \(308\)
groupoid packet is empty, the \(\texttt{weyl\_lifts}\) table is empty,
and the \(\texttt{coxeter\_cocycles}\) table is empty.  Thus row
\(309\) defines the orientation \(2\)-cocycle formula, but no finite
cocycle table is currently populated.
\end{definition}

The packet
\[
\texttt{certificates/orientation/rhomred\_weyl\_orientation\_cocycle}
\]
verifies
\(\texttt{RHOMRED\_WEYL\_ORIENTATION\_COCYCLE\_OBSTRUCTION\_VERIFIED}\):
it imports the row \(308\) partial-action-groupoid packet, the row
\(306\) and row \(307\) lift criteria, and the blocked
reduced-orientation ledger; it records that the formula for
\(c_{o,R}\) is fixed but that the finite composable-pair, lift, and
cocycle-defect rows are absent.

\begin{proposition}[Vanishing criterion for the orientation cocycle
class]
\label{prop:orientation-cocycle-class-vanishing-criterion}
\textnormal{[criterion proved; cohomology-vanishing rows not supplied]}
Let \(\mathcal G_R\), the coefficient system
\(\underline{\mathbb F}_{2,o}\), and the normalized cocycle
\(c_{o,R}\in Z^2(\mathcal G_R;\underline{\mathbb F}_{2,o})\) be
supplied as in Definition~\ref{def:orientation-projective-cocycle-coR}.
Choose ordered finite bases for the normalized groupoid cochain spaces
\[
C^1_R=C^1(\mathcal G_R;\underline{\mathbb F}_{2,o}),\qquad
C^2_R=C^2(\mathcal G_R;\underline{\mathbb F}_{2,o}),
\]
write the coboundary as a matrix
\[
d^1_R:C^1_R\longrightarrow C^2_R,
\]
and write the cocycle as a column vector \(v(c_{o,R})\in C^2_R\).
Then
\[
[c_{o,R}]=0\in
H^2(\mathcal G_R;\underline{\mathbb F}_{2,o})
\]
if and only if
\[
\operatorname{rank}_{\mathbb F_2}(d^1_R)
=
\operatorname{rank}_{\mathbb F_2}\bigl[d^1_R\mid v(c_{o,R})\bigr].
\]
Equivalently, the augmented linear system
\[
d^1_R b=v(c_{o,R})
\]
is soluble over \(\mathbb F_2\).  This is the row \(310\) criterion.
A supplied rank equality proves the vanishing of the cohomology class;
row \(311\) is the subsequent choice of a particular solution \(b\).

The criterion is not triggered by the target Coxeter graph, by the
Maass character, by the relation \(s_{\delta_i}^2=1\), or by an
abstract calculation of \(H^2(W^{(2)};\mathbb F_2)\) before the
retained groupoid, coefficient system, finite orbit, and cocycle
vector have been supplied.  The present finite tables contain no
values of \(c_{o,R}\), no normalized cochain-complex bases, no
coboundary matrix, and no rank certificate.  Hence row \(310\) remains
an open cohomological obstruction.
\end{proposition}

\begin{proof}
For any finite groupoid with coefficients in a local system of
\(\mathbb F_2\)-lines, the normalized cochain complex is a finite
cochain complex of \(\mathbb F_2\)-vector spaces.  A closed
\(2\)-cochain represents zero in \(H^2\) exactly when it lies in
\(\operatorname{im}d^1_R\).  After ordered bases are fixed, membership
in the image of a matrix is equivalent to equality of the rank of the
matrix and the rank of the matrix augmented by the target vector.
This proves the displayed criterion.  Inspection of
\(\texttt{certificates/orientation/rhomred\_weyl\_orientation\_cocycle}\)
and
\(\texttt{certificates/orientation/k3e\_reduced\_orientation/coxeter\_cocycles.csv}\)
shows that the required cocycle, complex, coboundary, and rank rows are
not present.
\end{proof}

The packet
\[
\texttt{certificates/orientation/rhomred\_weyl\_orientation\_cocycle\_class\_vanishing}
\]
verifies
\(\texttt{RHOMRED\_WEYL\_ORIENTATION\_COCYCLE\_CLASS\_VANISHING\_OBSTRUCTION\_VERIFIED}\):
it imports the row \(309\) cocycle packet and the blocked
reduced-orientation ledger, records the finite rank criterion for
\([c_{o,R}]=0\), and keeps the cocycle-value, closedness,
cochain-complex, coboundary-matrix, rank-certificate, and witness
cochain rows open.

\begin{proposition}[Cochain trivialization criterion for the
orientation cocycle]
\label{prop:orientation-cocycle-cochain-trivialization-criterion}
\textnormal{[criterion proved; killing \(1\)-cochain not supplied]}
Assume the data of
Proposition~\ref{prop:orientation-cocycle-class-vanishing-criterion}
are supplied and that row \(310\) has produced a rank certificate for
\([c_{o,R}]=0\).  A row \(311\) cochain trivialization is a named
vector
\[
b_R\in C^1(\mathcal G_R;\underline{\mathbb F}_{2,o})
\]
such that, in the ordered bases used for row \(310\),
\[
d^1_R b_R=v(c_{o,R}).
\]
Equivalently, it is a chosen lift of \(v(c_{o,R})\) through the
coboundary map \(d^1_R\), together with a defect row
\[
\operatorname{def}(b_R)
:=d^1_R b_R-v(c_{o,R})=0
\]
over \(\mathbb F_2\).  The class vanishing in row \(310\) proves that
such \(b_R\) exists; row \(311\) chooses one, records its coordinates,
and verifies the displayed equation.  Compatibility of these choices
under \(R\to R'\) is row \(312\), not part of row \(311\).

The present finite tables do not construct row \(311\).  The row
\(309\) cocycle values are absent, the row \(310\) rank-certificate
rows are absent, and no vector \(b_R\), gauge normalization, or
coboundary-defect row is supplied.
\end{proposition}

\begin{proof}
Once ordered bases are fixed, a \(1\)-cochain is a column vector in the
finite \(\mathbb F_2\)-vector space \(C^1_R\).  The condition that it
kill the projective cocycle is the single matrix equation
\(d^1_R b_R=v(c_{o,R})\).  Therefore the construction of \(b_R\) is
stronger than the rank equality of row \(310\): it supplies an actual
preimage, not only the statement that a preimage exists.  The residual
freedom is the addition of a \(1\)-cocycle in
\(\ker d^1_R\), so a complete finite row must also record either a
gauge convention or the torsor in which the choice lives.  Inspection
of the row \(309\) and row \(310\) packets shows that the input
cocycle, the coboundary matrix, the rank certificate, and the witness
vector are not present.
\end{proof}

The packet
\[
\texttt{certificates/orientation/rhomred\_weyl\_orientation\_cocycle\_cochain\_trivialization}
\]
verifies
\(\texttt{RHOMRED\_WEYL\_ORIENTATION\_COCYCLE\_COCHAIN\_TRIVIALIZATION\_OBSTRUCTION\_VERIFIED}\):
it imports the row \(309\) cocycle packet and the row \(310\)
class-vanishing packet, records the equation defining the killing
\(1\)-cochain, and keeps the cochain-vector, gauge-normalization, and
coboundary-defect rows open.

\begin{proposition}[Transition compatibility for the Coxeter cochain]
\label{prop:orientation-cocycle-cochain-transition-compatibility}
\textnormal{[criterion proved; transition rows not supplied]}
Let \(R'\le R\).  Suppose rows \(309\)--\(311\) have supplied
\(\mathcal G_R,\mathcal G_{R'}\), the coefficient systems, cocycles,
cochain complexes, and chosen cochains \(b_R\) and \(b_{R'}\).  A row
\(312\) transition compatibility datum consists of cochain maps
\[
\rho^i_{RR'}:
C^i(\mathcal G_{R'};\underline{\mathbb F}_{2,o})
\longrightarrow
C^i(\mathcal G_R;\underline{\mathbb F}_{2,o}),
\qquad i=1,2,
\]
induced by the finite-stage transition \(t_{RR'}\), such that
\[
d^1_R\rho^1_{RR'}=\rho^2_{RR'}d^1_{R'},\qquad
\rho^2_{RR'}v(c_{o,R'})=v(c_{o,R}),
\]
and
\[
\rho^1_{RR'}(b_{R'})=b_R.
\]
Equivalently, the transition defect
\[
\Delta^b_{RR'}:=b_R-\rho^1_{RR'}(b_{R'})
\in C^1(\mathcal G_R;\underline{\mathbb F}_{2,o})
\]
must be zero in the chosen gauge.  For a chain \(R''\le R'\le R\),
the cochain maps must also satisfy
\[
\rho^1_{RR''}=\rho^1_{RR'}\rho^1_{R'R''},\qquad
\rho^2_{RR''}=\rho^2_{RR'}\rho^2_{R'R''}
\]
on the retained cochain bases.  Thus the \(b_R\) form a strict inverse
system of cochain trivializations.  Compatibility only up to an
unrecorded \(1\)-cocycle is not enough; the comparison cocycle or the
gauge fixing must be part of the finite row.

The present tables do not prove row \(312\).  No row \(311\) cochain
vectors are supplied, the transition-orientation packet has no Coxeter
cochain transition row, and the reduced-orientation transition table
is empty.
\end{proposition}

\begin{proof}
The first displayed equation says that the transition maps form a
cochain map.  The second says that the projective cocycle at level
\(R'\) pulls back to the projective cocycle at level \(R\).  If
\(d^1_{R'}b_{R'}=v(c_{o,R'})\), these two equations imply
\[
d^1_R\rho^1_{RR'}(b_{R'})
=\rho^2_{RR'}d^1_{R'}b_{R'}
=\rho^2_{RR'}v(c_{o,R'})
=v(c_{o,R}).
\]
Thus both \(\rho^1_{RR'}(b_{R'})\) and \(b_R\) kill \(c_{o,R}\).
Row \(312\) asks that the chosen lifts agree, not merely that both
exist.  The defect \(\Delta^b_{RR'}\) is therefore a closed
\(1\)-cochain; strict compatibility is exactly the assertion
\(\Delta^b_{RR'}=0\) in the fixed gauge, together with the two-step
composition equations for the \(\rho^i\).  Inspection of the row
\(311\) packet, the transition-orientation packet, and
\(\texttt{certificates/orientation/k3e\_reduced\_orientation/transitions.csv}\)
shows that none of these finite rows are present.
\end{proof}

The packet
\[
\texttt{certificates/orientation/rhomred\_weyl\_orientation\_cocycle\_cochain\_transition}
\]
verifies
\(\texttt{RHOMRED\_WEYL\_ORIENTATION\_COCYCLE\_COCHAIN\_TRANSITION\_OBSTRUCTION\_VERIFIED}\):
it imports the row \(311\) cochain packet and the transition-orientation
ledger, records the equations defining compatibility under
\(R\to R'\), and keeps the transition cochain-map, cocycle-pullback,
cochain-defect, gauge, and two-step coherence rows open.

\begin{proposition}[Computation criterion for the Weyl torsor defects]
\label{prop:weyl-orientation-torsor-defect-computation}
\textnormal{[criterion proved; torsor-defect rows not supplied]}
Fix \(R\), a retained quotient stratum \(S\), and a simple type-II
reflection \(s_{\delta_i}\).  Suppose the row \(306\) lift
\(\tau_{i,R,S}\) and the quotient-orientation packages
\[
\mathfrak q_{R,S},\qquad
\mathfrak q_{R,s_{\delta_i}S}
\]
are supplied by explicit Cech--Borel representatives, null
trivialisations, finite-stabilizer edge representatives, and
linearization characters.  Let
\[
\Theta^1_{i,R,S}\colon
\mathcal Q^1_{R,S}\longrightarrow
\mathcal Q^1_{R,s_{\delta_i}S}
\]
be the degree-one cochain map induced by \(\tau_{i,R,S}\) on the
quotient-orientation complex
\[
\mathcal Q^\bullet_{R,S}:=
C^\bullet(\mathfrak M^{\mathrm{red}}_{R,S};\mathbb F_2)
\oplus
\bigoplus_{H\in\mathcal H(R,S)}C^\bullet(BH;\mathbb F_2),
\]
with the connected \(E\)-free summand already killed by its row \(284\)
null-trivialization.  If
\(\eta_{R,S}\in\mathcal Q^1_{R,S}\) and
\(\eta_{R,s_{\delta_i}S}\in\mathcal Q^1_{R,s_{\delta_i}S}\) denote the
chosen quotient-orientation null trivialisations, the torsor-defect
cochain is
\[
\Delta^\omega_{i,R,S}:=
\Theta^1_{i,R,S}(\eta_{R,S})-\eta_{R,s_{\delta_i}S}.
\]
After the degree-two quotient representatives have been transported
with zero defect, \(\Delta^\omega_{i,R,S}\) is closed.  The row \(313\)
torsor defect is the cohomology class
\[
\omega_{i,R,S}:=
[\Delta^\omega_{i,R,S}]
\in
H^1(\mathfrak M^{\mathrm{red}}_{R,s_{\delta_i}S};\mathbb F_2)
\oplus
\bigoplus_{H\in\mathcal H(R,s_{\delta_i}S)}
H^1(BH;\mathbb F_2),
\]
recorded in a finite basis for the displayed \(H^1\)-space.  Thus row
\(313\) computes the defect vector; row \(314\) is the separate
assertion that this vector is zero for every retained \(i,R,S\).

The present finite tables do not compute row \(313\).  They contain no
row \(306\) Weyl lifts, no quotient Borel representatives, no
finite-stabilizer transport rows, no linearization transport rows, no
\(\mathcal Q^\bullet\)-basis rows, and no torsor-defect vector rows.
\end{proposition}

\begin{proof}
The quotient orientation is not a single line: it is the orientation
line together with null trivialisations of the reduced gerbe, the
free-\(E\) Borel class, the finite-stabilizer Borel classes, and the
zero linearization characters.  Transport along \(\tau_{i,R,S}\)
therefore gives cochain maps on the Cech--Borel complexes carrying
these representatives.  If the degree-two representatives are
transported with zero defect, the difference between the transported
source null-trivialization and the target null-trivialization has zero
coboundary; hence it defines the displayed \(H^1\)-class.  Coordinates
of this class in a chosen finite basis are exactly the torsor-defect
computation.  A target Weyl relation, Maass-character value, or local
Pfaffian sign has no such cochain representative and cannot compute
\(\omega_{i,R,S}\).  Inspection of the reduced-orientation and
row \(306\) packets shows that the required rows are absent.
\end{proof}

The packet
\[
\texttt{certificates/orientation/rhomred\_weyl\_orientation\_torsor\_defects}
\]
verifies
\(\texttt{RHOMRED\_WEYL\_ORIENTATION\_TORSOR\_DEFECTS\_OBSTRUCTION\_VERIFIED}\):
it imports the row \(306\) Weyl-wall transport criterion, the row
\(312\) cochain-transition packet, and the blocked reduced-orientation
ledger; it records the finite cochain formula for
\(\omega_{i,R,S}\) and keeps the quotient-orientation representative,
transport, \(H^1\)-basis, and defect-vector rows open.

\begin{proposition}[Vanishing criterion for the Weyl torsor defects]
\label{prop:weyl-orientation-torsor-defect-vanishing}
\textnormal{[criterion proved; zero-defect rows not supplied]}
Assume the row \(313\) data of
Proposition~\ref{prop:weyl-orientation-torsor-defect-computation} have
been supplied for every retained triple \((i,R,S)\).  Thus each
\(\Delta^\omega_{i,R,S}\) is a closed cochain in the target
quotient-orientation complex and has coordinates
\[
v(\omega_{i,R,S})\in
H^1(\mathfrak M^{\mathrm{red}}_{R,s_{\delta_i}S};\mathbb F_2)
\oplus
\bigoplus_{H\in\mathcal H(R,s_{\delta_i}S)}
H^1(BH;\mathbb F_2)
\]
in a fixed finite basis.  Row \(314\) is proved exactly when the
finite vanishing table contains, for every retained triple
\((i,R,S)\), the following four rows:
\[
d\Delta^\omega_{i,R,S}=0,\qquad
v(\omega_{i,R,S})=0,\qquad
\operatorname{rank}\bigl(v(\omega_{i,R,S})\bigr)=0,
\qquad
\mathrm{cov}\{(i,R,S)\}=\mathrm{cov}_{\mathrm{ret}} .
\]
Equivalently, the row \(313\) defect vectors must be present, closed,
computed in the displayed \(H^1\)-basis, and zero with complete
coverage over the retained Weyl-wall atlas.  Under these hypotheses
the Weyl lift \(\tau_{i,R,S}\) preserves the quotient-orientation
torsor, and \(\omega_{i,R,S}=0\) for all retained \(i,R,S\).

The present finite tables do not prove row \(314\).  The row \(313\)
packet contains no torsor-defect vector rows, no \(H^1\)-basis rows,
no zero-vector certificates, and no coverage rows for retained
triples.  Therefore the vanishing of \(\omega_{i,R,S}\) remains an
open orientation obligation.
\end{proposition}

\begin{proof}
In the fixed \(H^1\)-basis, a closed defect cochain represents the
zero torsor class if and only if its coordinate vector is zero.  Thus
the four displayed finite checks are precisely the statement that
each class \(\omega_{i,R,S}\) vanishes, together with the assertion
that no retained type-II wall has been omitted.  This is the needed
source-side orientation statement: it says that the transported null
trivialisation equals the target null trivialisation in the quotient
orientation torsor.  A target Weyl relation, a Maass character, a
Pfaffian wall sign, or an empty defect table is not a coordinate
calculation in the displayed \(H^1\)-space.  Inspection of the row
\(313\) packet shows that the required defect vectors and zero
certificates are not supplied.
\end{proof}

The packet
\[
\texttt{certificates/orientation/rhomred\_weyl\_orientation\_torsor\_defect\_vanishing}
\]
verifies
\(\texttt{RHOMRED\_WEYL\_ORIENTATION\_TORSOR\_DEFECT\_VANISHING\_OBSTRUCTION\_VERIFIED}\):
it imports the row \(313\) torsor-defect computation criterion,
records the finite zero-vector criterion for row \(314\), and keeps
the computed defect-vector, closedness, zero-certificate, and retained
coverage rows open.

\begin{proposition}[Weyl preservation criterion for the reduced gerbe]
\label{prop:weyl-alpha-red-preservation-criterion}
\textnormal{[criterion proved; preservation rows not supplied]}
Fix \(R\), a retained quotient stratum \(S\), and a simple type-II
reflection \(s_{\delta_i}\).  Suppose the row \(306\) lift
\(\tau_{i,R,S}\) is supplied, and suppose row \(280\) supplies
Cech--Borel cocycles
\[
a^{\mathrm{red}}_{R,S}\in Z^2(\mathcal B_{R,S};\mathbb F_2),
\qquad
a^{\mathrm{red}}_{R,s_{\delta_i}S}
\in Z^2(\mathcal B_{R,s_{\delta_i}S};\mathbb F_2)
\]
representing
\(\alpha^{\mathrm{red}}_{R,S}\) and
\(\alpha^{\mathrm{red}}_{R,s_{\delta_i}S}\).  Let
\[
\Theta^{\bullet,\mathrm{red}}_{i,R,S}\colon
C^\bullet(\mathcal B_{R,S};\mathbb F_2)
\longrightarrow
C^\bullet(\mathcal B_{R,s_{\delta_i}S};\mathbb F_2)
\]
be the cochain map induced by the Weyl transport on the reduced
Cech--Borel groupoids.  The row \(315\) preservation defect is
\[
\Delta^{\mathrm{red}}_{i,R,S}:=
\Theta^{2,\mathrm{red}}_{i,R,S}(a^{\mathrm{red}}_{R,S})
-a^{\mathrm{red}}_{R,s_{\delta_i}S}.
\]
Weyl transport preserves the reduced gerbe on this retained wall
exactly when there is a finite cochain
\[
u^{\mathrm{red}}_{i,R,S}\in
C^1(\mathcal B_{R,s_{\delta_i}S};\mathbb F_2)
\]
with
\[
d u^{\mathrm{red}}_{i,R,S}=\Delta^{\mathrm{red}}_{i,R,S};
\]
equivalently,
\[
[\Delta^{\mathrm{red}}_{i,R,S}]=0
\in H^2(\mathcal B_{R,s_{\delta_i}S};\mathbb F_2).
\]
The finite row proving row \(315\) must record the source and target
\(\alpha^{\mathrm{red}}\)-cocycles, the \(\Theta^{2,\mathrm{red}}\)
matrix, the transported defect vector, the coboundary witness
\(u^{\mathrm{red}}_{i,R,S}\), a zero class-rank certificate, and
coverage of every retained triple \((i,R,S)\).  If strict
representative preservation is imposed as the chosen gauge, the same
row must record the stronger condition
\(\Delta^{\mathrm{red}}_{i,R,S}=0\).

The present finite tables do not prove row \(315\).  The row \(280\)
packet contains no computed \(\alpha^{\mathrm{red}}\)-values, the row
\(281\) packet contains no Cech--Borel null-homotopies, the row \(306\)
packet contains no Weyl-lift cochain maps, and row \(314\) has not
proved vanishing of the quotient-orientation torsor defects.
\end{proposition}

\begin{proof}
Because \(\Theta^{\bullet,\mathrm{red}}_{i,R,S}\) is required to be a
cochain map, the displayed defect is closed whenever the source and
target representatives are closed and the degree-three transport
defect is zero.  The transported reduced gerbe class equals the target
reduced gerbe class precisely when this closed defect is a coboundary.
A chosen cochain \(u^{\mathrm{red}}_{i,R,S}\) with
\(d u^{\mathrm{red}}_{i,R,S}=\Delta^{\mathrm{red}}_{i,R,S}\) is the
finite certificate of equality in \(H^2\).  Complete retained-wall
coverage is part of the statement because \(O1^+\) is a Weyl-equivariant
orientation assertion, not a single-wall computation.  A target Weyl
relation, a scalar multiplier, a Pfaffian sign, or the uncomputed
vanishing of the full torsor defect is not a cochain calculation for
\(\alpha^{\mathrm{red}}\).  Inspection of the row \(280\), row \(281\),
row \(306\), and row \(314\) packets shows that the needed rows are
absent.
\end{proof}

The packet
\[
\texttt{certificates/orientation/rhomred\_weyl\_alpha\_red\_preservation}
\]
verifies
\(\texttt{RHOMRED\_WEYL\_ALPHA\_RED\_PRESERVATION\_OBSTRUCTION\_VERIFIED}\):
it imports the reduced-gerbe, reduced-gerbe null-homotopy, Weyl-wall
transport, and torsor-defect vanishing packets; it records the finite
cochain criterion for row \(315\); and it keeps the source-target
\(\alpha^{\mathrm{red}}\)-cocycle, transport matrix, coboundary
witness, zero class-rank, and retained coverage rows open.

\begin{proposition}[Weyl preservation criterion for the connected
free \(E\)-Borel class]
\label{prop:weyl-alpha-e-free-preservation-criterion}
\textnormal{[criterion proved; preservation rows not supplied]}
Fix \(R\), a retained free \(E\)-quotient stratum \(S\), and a simple
type-II reflection \(s_{\delta_i}\).  Suppose the row \(306\) lift
\(\tau_{i,R,S}\) is supplied and suppose row \(282\) supplies
connected \(BE\)-edge representatives
\[
a^{E,\mathrm{free}}_{R,S},\qquad
a^{E,\mathrm{free}}_{R,s_{\delta_i}S}
\in Z^2(BE;\mathbb F_2)
\]
with coordinates
\[
\alpha^{E,\mathrm{free}}_{R,S}=a_1(R,S)u_1+a_2(R,S)u_2,\qquad
\alpha^{E,\mathrm{free}}_{R,s_{\delta_i}S}
=a_1(R,s_{\delta_i}S)u_1+a_2(R,s_{\delta_i}S)u_2
\]
in the basis of Lemma~\ref{lem:connected-e-descent}.  Let
\[
M^{E}_{i,R,S}\in \operatorname{GL}_2(\mathbb F_2)
\]
be the induced map of the Weyl transport on
\(H^2(BE;\mathbb F_2)=\mathbb F_2u_1\oplus\mathbb F_2u_2\), computed
from the cochain-level transport of the connected \(E\)-translation
rigidification.  The row \(316\) preservation defect is the coordinate
vector
\[
\Delta^{E,\mathrm{free}}_{i,R,S}:=
M^{E}_{i,R,S}
\begin{pmatrix}a_1(R,S)\\ a_2(R,S)\end{pmatrix}
-
\begin{pmatrix}a_1(R,s_{\delta_i}S)\\
a_2(R,s_{\delta_i}S)\end{pmatrix}
\in\mathbb F_2^2 .
\]
Weyl transport preserves the connected free \(E\)-Borel class on this
retained wall exactly when
\[
\Delta^{E,\mathrm{free}}_{i,R,S}=0.
\]
Equivalently, in a cochain model one may record a degree-one witness
for the difference between the transported source representative and
the target representative; after passing to the connected \(BE\)-edge,
this is the same as the displayed zero coordinate vector, since
\(u_1,u_2\) form a basis of \(H^2(BE;\mathbb F_2)\).

The finite row proving row \(316\) must therefore contain the source
and target \(a_1,a_2\) coefficient rows, the induced matrix
\(M^E_{i,R,S}\), the transported coordinate vector, the zero-vector
certificate, the underlying cochain-transport defect if strict
representatives are used, and coverage of every retained triple
\((i,R,S)\).  The present finite tables do not prove row \(316\):
rows \(282\), \(283\), and \(284\) supply no connected \(BE\)-edge
coefficients, zero checks, or null-trivialisations; row \(306\)
supplies no Weyl cochain map; and row \(315\) has not proved
preservation of \(\alpha^{\mathrm{red}}\).
\end{proposition}

\begin{proof}
The connected part of the quotient orientation is controlled by the
degree-two edge of \(H^*_E(-;\mathbb F_2)\).  In degree two the
connected \(BE\)-edge has the fixed basis \(u_1,u_2\).  Therefore a
Weyl transport preserves \(\alpha^{E,\mathrm{free}}\) if and only if
the transported coordinate vector equals the target coordinate vector.
The displayed equation is exactly this equality in
\(\mathbb F_2^2\).  If the comparison is made before passage to the
edge, the difference of cocycle representatives must be recorded as a
coboundary; its edge coordinate is the same defect vector.  Freeness of
the \(E\)-action, \(H^1(BE;\mathbb F_2)=0\), target Weyl relations,
Maass-character values, and Pfaffian signs do not compute the
\(H^2(BE)\) coordinate transport.  Inspection of the row \(282\), row
\(283\), row \(284\), row \(306\), and row \(315\) packets shows that
the required finite rows are absent.
\end{proof}

The packet
\[
\texttt{certificates/orientation/rhomred\_weyl\_alpha\_e\_free\_preservation}
\]
verifies
\(\texttt{RHOMRED\_WEYL\_ALPHA\_E\_FREE\_PRESERVATION\_OBSTRUCTION\_VERIFIED}\):
it imports the connected free \(E\)-Borel computation, vanishing,
null-trivialization, Weyl-wall transport, and reduced-gerbe
preservation packets; it records the finite coordinate criterion for
row \(316\); and it keeps the source-target \(a_1,a_2\) rows, induced
\(H^2(BE)\)-matrix, zero-vector certificate, cochain-transport defect,
and retained coverage rows open.

\begin{proposition}[Weyl preservation criterion for the finite-stabilizer
Borel classes]
\label{prop:weyl-finite-stabilizer-beta-preservation-criterion}
\textnormal{[criterion proved; preservation rows not supplied]}
Fix \(R\), a retained finite-stabilizer stratum \(S\), and a simple
type-II reflection \(s_{\delta_i}\).  Suppose the row \(306\) lift
\(\tau_{i,R,S}\) is supplied and identifies the source and target
finite stabilizers by an isomorphism
\[
\varphi_{i,R,S}\colon H_{R,S}\xrightarrow{\sim}
H_{R,s_{\delta_i}S}.
\]
Suppose row \(285\) has supplied classifying-space edge representatives
\[
\beta^{H_{R,S}}_{R,S}\in H^2(BH_{R,S};\mathbb F_2),\qquad
\beta^{H_{R,s_{\delta_i}S}}_{R,s_{\delta_i}S}
\in H^2(BH_{R,s_{\delta_i}S};\mathbb F_2),
\]
with coordinates in the appropriate odd, Klein-four, or two-primary
basis of Proposition~\ref{prop:finite-stabilizer-beta-class-criterion}.
Let
\[
M^{H}_{i,R,S}\colon
H^2(BH_{R,S};\mathbb F_2)\longrightarrow
H^2(BH_{R,s_{\delta_i}S};\mathbb F_2)
\]
be the matrix induced by \(B\varphi_{i,R,S}\) on the chosen
classifying-space bases.  The row \(317\) preservation defect is
\[
\Delta^{\beta}_{i,R,S}:=
M^{H}_{i,R,S}\,v(\beta^{H_{R,S}}_{R,S})
-v(\beta^{H_{R,s_{\delta_i}S}}_{R,s_{\delta_i}S}).
\]
Weyl transport preserves the finite-stabilizer Borel class on this
retained wall exactly when
\[
\Delta^{\beta}_{i,R,S}=0
\]
for every retained triple \((i,R,S)\).  If the comparison is made at
the Cech--Borel representative level, the row must also record the
transport of \(\widetilde\beta^H\), the edge-reduction compatibility,
and a zero cochain-level transport defect before taking the displayed
classifying-space coordinate vector.

Thus a row proving \(317\) must contain the source and target
finite-stabilizer orientation rows, the stabilizer isomorphism
\(\varphi_{i,R,S}\), the induced \(H^2(BH)\)-matrix, the transported
\(\beta^H\)-coordinate vector, a zero-vector certificate, edge-reduction
compatibility, and complete retained-wall coverage.  The present
finite tables do not prove row \(317\): row \(285\) has no
finite-stabilizer orientation or edge rows, row \(286\) has no
\(\beta^H=0\) rows, row \(287\) has no compatible
null-trivializations, row \(306\) has no Weyl lift on stabilizers, and
row \(316\) has not proved preservation of \(\alpha^{E,\mathrm{free}}\).
\end{proposition}

\begin{proof}
The finite-stabilizer component of the quotient orientation is the
classifying-space edge of the \(H\)-equivariant Borel class
\(\widetilde\beta^H\).  A Weyl lift can preserve this component only
after it identifies the source and target stabilizer groups and carries
the source edge representative to the target edge representative.  In
the fixed group-cohomology bases this condition is precisely the
vanishing of the displayed coordinate defect.  For odd or trivial
stabilizers the target space is zero only after the finite-stabilizer
edge row has been supplied; for rank-two two-primary stabilizers the
mixed coefficient is part of the basis and cannot be read from cyclic
restrictions.  A residual group order, a scalar trace, a Maass
character, a Pfaffian sign, or preservation of the connected free
\(E\)-component is not a finite-stabilizer edge transport calculation.
Inspection of the row \(285\), row \(286\), row \(287\), row \(306\),
and row \(316\) packets shows that the required finite rows are absent.
\end{proof}

The packet
\[
\texttt{certificates/orientation/rhomred\_weyl\_finite\_stabilizer\_beta\_preservation}
\]
verifies
\(\texttt{RHOMRED\_WEYL\_FINITE\_STABILIZER\_BETA\_PRESERVATION\_OBSTRUCTION\_VERIFIED}\):
it imports the finite-stabilizer Borel computation, vanishing,
null-trivialization, Weyl-wall transport, and connected free
\(E\)-preservation packets; it records the finite edge-coordinate
criterion for row \(317\); and it keeps the stabilizer isomorphism,
source-target \(\beta^H\)-coordinates, \(H^2(BH)\)-matrix,
edge-transport defect, zero-vector certificate, and retained coverage
rows open.

\begin{proposition}[Weyl preservation criterion for the zero
finite-stabilizer linearization]
\label{prop:weyl-finite-stabilizer-lambda-zero-preservation-criterion}
\textnormal{[criterion proved; linearization-transport rows not
supplied]}
Fix \(R\), a retained finite-stabilizer stratum \(S\), and a simple
type-II reflection \(s_{\delta_i}\).  Suppose the row \(306\) lift
\(\tau_{i,R,S}\) is supplied and identifies the source and target
finite stabilizers by an isomorphism
\[
\varphi_{i,R,S}\colon H_{R,S}\xrightarrow{\sim}
H_{R,s_{\delta_i}S}.
\]
Suppose row \(288\) has supplied the finite-stabilizer
linearization characters
\[
\lambda^{H_{R,S}}_{R,S}\in H^1(BH_{R,S};\mathbb F_2),
\qquad
\lambda^{H_{R,s_{\delta_i}S}}_{R,s_{\delta_i}S}
\in H^1(BH_{R,s_{\delta_i}S};\mathbb F_2),
\]
with coordinates in the bases of
Proposition~\ref{prop:finite-stabilizer-linearization-character-criterion},
and suppose row \(289\) has supplied the corresponding zero rows.
Let
\[
N^{H}_{i,R,S}\colon
H^1(BH_{R,S};\mathbb F_2)\longrightarrow
H^1(BH_{R,s_{\delta_i}S};\mathbb F_2)
\]
be the matrix induced by \(B\varphi_{i,R,S}\) on the chosen
degree-one classifying-space bases.  The row \(318\)
linearization-transport defect is
\[
\Delta^{\lambda}_{i,R,S}:=
N^H_{i,R,S}\,v(\lambda^{H_{R,S}}_{R,S})
-v(\lambda^{H_{R,s_{\delta_i}S}}_{R,s_{\delta_i}S}).
\]
Weyl transport preserves the zero finite-stabilizer linearization on
this retained wall exactly when
\[
\Delta^{\lambda}_{i,R,S}=0
\]
for every retained triple \((i,R,S)\).  If the comparison is made at
the orientation-line action level, the row must also record the
transport of the sign action of \(H_{R,S}\) to
\(H_{R,s_{\delta_i}S}\) and a zero action-transport defect before
passing to \(H^1(BH;\mathbb F_2)\).

Thus a row proving \(318\) must contain the source and target
\(\lambda^H\)-coordinate rows, the stabilizer isomorphism
\(\varphi_{i,R,S}\), the induced \(H^1(BH)\)-matrix, the transported
coordinate vector, a rank-zero defect row, orientation-line action
transport, and complete retained-wall coverage.  The present finite
tables do not prove row \(318\): row \(288\) has no orientation-line
action rows, generator-basis rows, character values, or computed
\(\lambda^H\)-coordinates; row \(289\) has no zero-coordinate rows;
row \(306\) has no Weyl lift on finite stabilizers; and row \(317\)
has not proved preservation of the degree-two finite-stabilizer Borel
classes.
\end{proposition}

\begin{proof}
After the degree-two finite-stabilizer gerbe has been trivialized, the
remaining descent ambiguity of the orientation line is the sign
character
\[
\lambda^H\in H^1(BH;\mathbb F_2)=\operatorname{Hom}(H,\mathbb F_2).
\]
A Weyl lift can preserve the zero character only after it identifies
the source and target stabilizer groups and transports the source sign
action to the target sign action.  In the chosen \(H^1(BH)\)-bases,
this condition is exactly the vanishing of the displayed coordinate
defect.  If rows \(288\) and \(289\) have supplied zero characters on
both sides, the class-level equality follows from functoriality of
group cohomology, but the transport still has no mathematical meaning
without the stabilizer isomorphism, the induced degree-one matrix, and
the action-level comparison.  The identity
\(H^1(B1;\mathbb F_2)=0\), a vanishing \(\beta^H\)-class, a
null-trivialization of \(\beta^H\), a connected \(E\)-Borel
calculation, a Maass-character value, a Pfaffian wall sign, or a scalar
trace does not compute this degree-one transport.  Inspection of the
row \(288\), row \(289\), row \(306\), and row \(317\) packets shows
that the required finite rows are absent.
\end{proof}

The packet
\[
\texttt{certificates/orientation/rhomred\_weyl\_finite\_stabilizer\_lambda\_zero\_preservation}
\]
verifies
\(\texttt{RHOMRED\_WEYL\_FINITE\_STABILIZER\_LAMBDA\_ZERO\_PRESERVATION\_OBSTRUCTION\_VERIFIED}\):
it imports the finite-stabilizer linearization-character,
linearization-vanishing, Weyl-wall transport, and finite-stabilizer
Borel-preservation packets; it records the degree-one coordinate
criterion for row \(318\); and it keeps the stabilizer isomorphism,
source-target \(\lambda^H\)-coordinates, \(H^1(BH)\)-matrix,
orientation-line action transport, coordinate defect, rank-zero row,
and retained coverage rows open.

\begin{definition}[Orientation character of the strictified Weyl
transport]
\label{def:orientation-character-epsilon-o}
\textnormal{[definition specified; finite character rows not supplied]}
Assume rows \(306\)--\(318\) have supplied a strictified
\(W^{(2)}(\Lambda_{II}^{2,1})\)-transport of the complete
quotient-oriented determinant-line package on a retained finite
height \(R\): the finite partial action groupoid \(\mathcal G_R\),
the lifts \(\tau_{i,R,S}\), the cochain \(b_R\) killing the
projective cocycle \(c_{o,R}\), transition compatibility of \(b_R\),
vanishing torsor defects, and preservation of
\(\alpha^{\mathrm{red}}\), \(\alpha^{E,\mathrm{free}}\),
\(\beta^H\), and \(\lambda^H=0\).  For each simple type-II reflection
\(s_{\delta_i}\) and retained source stratum \(S\), let
\[
\sigma^o_{i,R,S}\in\{\pm1\}
\]
be the residual sign of the strictified lift on the oriented
one-dimensional Pfaffian square-root line after the target line and
the complete quotient-orientation package have been identified with
the transported source package.  A row defining the finite-stage
orientation character must prove that \(\sigma^o_{i,R,S}\) is
independent of \(S\) on every retained connected wall component and is
compatible under \(R'\le R\).  Write the resulting sign as
\(\sigma^o_{i,R}\).

The finite-stage orientation character is then the homomorphism
\[
\epsilon_{o,R}\colon W^{(2)}(\Lambda_{II}^{2,1})\longrightarrow
\{\pm1\}
\]
defined on a reduced word
\[
w=s_{\delta_{i_1}}\cdots s_{\delta_{i_m}}
\]
by
\[
\epsilon_{o,R}(w):=
\prod_{a=1}^{m}\sigma^o_{i_a,R}.
\]
The type-II Coxeter presentation has only the relations
\[
s_{\delta_i}^2=1,\qquad i=1,2,3,
\]
and no braid relations between distinct simple reflections.  Hence the
displayed formula is well-defined once the generator signs are
supplied, since \((\sigma^o_{i,R})^2=1\).  A compatible inverse system
of finite characters defines
\[
\epsilon_o=\varprojlim_R\epsilon_{o,R}.
\]

This is a definition from source-side orientation transport.  It is
not the automorphic Maass character
\(\nu_{\Delta_5}\), not the determinant character of the target Coxeter
group, not the OP scalar branch, and not a consequence of the scalar
trace.  Row \(320\) computes the generator values
\(\epsilon_o(s_{\delta_i})\); row \(321\) compares the resulting
character with the determinant character; row \(322\) compares it with
\(\nu_{\Delta_5}\).  The present finite tables do not define
\(\epsilon_{o,R}\): row \(312\) has no compatible cochain \(b_R\), row
\(314\) has no zero torsor-defect rows, rows \(315\)--\(318\) do not
prove preservation of the four quotient-orientation components, and no
generator-sign or transition-compatible character row is supplied.
\end{definition}

The packet
\[
\texttt{certificates/orientation/rhomred\_weyl\_orientation\_character\_definition}
\]
verifies
\(\texttt{RHOMRED\_WEYL\_ORIENTATION\_CHARACTER\_DEFINITION\_OBSTRUCTION\_VERIFIED}\):
it imports the Coxeter-cochain transition, torsor-defect vanishing,
and four component-preservation packets; it records the word-evaluation
definition of \(\epsilon_{o,R}\); it imports the automorphic Maass
character only as a firewall; and it keeps the strict Weyl action,
generator-sign, stratum-independence, transition-compatibility,
finite-character, and inverse-limit character rows open.

\begin{proposition}[Simple-wall sign criterion for the orientation
character]
\label{prop:orientation-character-simple-wall-sign-criterion}
\textnormal{[local sign formula proved; finite simple-wall sign rows
not supplied]}
Assume the finite-stage orientation character of
Definition~\ref{def:orientation-character-epsilon-o} has been supplied.
Fix a simple type-II reflection \(s_{\delta_i}\), a retained generic
wall stratum \(S\subset \Hpl_{\delta_i}\), and a retained
\((O2_{\mathrm{atlas}})\) wall chart with real normal coordinate
\(u_i\).  Suppose the chart supplies the rank-one normal Pfaffian
crossing
\[
K^{\mathrm{nor}}_{\delta_i,R,S}
\simeq [\mathbb R\xrightarrow{u_i}\mathbb R],
\qquad
\operatorname{Pf}(K^{\mathrm{nor}}_{\delta_i,R,S})
=\upsilon_{i,R,S}u_i,
\]
with
\[
s_{\delta_i}(u_i)=-u_i,\qquad
s_{\delta_i}(\upsilon_{i,R,S})=\upsilon_{i,R,S},\qquad
N_{\delta_i,R,S}^{\mathrm{Pf}}=1.
\]
Then the source-side generator sign of row \(320\) is
\[
\sigma^o_{i,R,S}
=\frac{s_{\delta_i}^{*}(\upsilon_{i,R,S}u_i)}
{\upsilon_{i,R,S}u_i}
=-1
=(-1)^{N_{\delta_i,R,S}^{\mathrm{Pf}}}.
\]
If the row also supplies stratum-independence and transition
compatibility for these signs, then
\[
\epsilon_{o,R}(s_{\delta_i})=-1,\qquad i=1,2,3,
\]
and the compatible inverse-limit character satisfies
\[
\epsilon_o(s_{\delta_i})=-1,\qquad i=1,2,3.
\]

Thus a row proving \(320\) must contain, for all three simple type-II
walls, retained wall-object rows, rank-one wall-chart rows, the normal
coordinate flip \(u_i\mapsto -u_i\), invariant tangent Pfaffian units,
normal rank \(1\), divisor order \(1\), local sign rows, the row \(319\)
orientation-character row, stratum-independence rows, and transition
compatibility rows.  The present finite tables do not prove row
\(320\): the row \(319\) packet has no constructed orientation
character or generator-sign rows, the \((O2_{\mathrm{atlas}})\) packet
has no wall-object, rank-one chart, coordinate-flip, invariant-unit, or
local-sign rows, and the finite Pfaffian packet has no type-II wall
chart or sign rows.
\end{proposition}

\begin{proof}
The local calculation is the rank-one Pfaffian crossing of
Proposition~\ref{prop:appE-local-pfaffian-sign}.  On the retained
chart the tangent Pfaffian unit is invariant and the reflection changes
only the normal coordinate:
\[
s_{\delta_i}^{*}(\upsilon_{i,R,S}u_i)
=\upsilon_{i,R,S}(-u_i)
=-\upsilon_{i,R,S}u_i.
\]
The normal complex contains one odd crossing, so
\(N_{\delta_i,R,S}^{\mathrm{Pf}}=1\) and the sign is
\((-1)^1=-1\).  Squaring removes this sign:
\[
s_{\delta_i}^{*}((\upsilon_{i,R,S}u_i)^2)
=(\upsilon_{i,R,S}u_i)^2.
\]
Hence squared determinants, scalar traces, and OP scalar branches
cannot recover the orientation character.  The Maass character value
\(\nu_{\Delta_5}(s_{\delta_i})=-1\) is an automorphic target value; it
does not supply the source wall chart, invariant unit, normal-rank row,
or generator-sign row.  Inspection of the row \(319\),
\((O2_{\mathrm{atlas}})\), and finite Pfaffian packets shows that those
finite rows are absent.
\end{proof}

The packet
\[
\texttt{certificates/orientation/rhomred\_weyl\_orientation\_simple\_sign\_computation}
\]
verifies
\(\texttt{RHOMRED\_WEYL\_ORIENTATION\_SIMPLE\_SIGN\_COMPUTATION\_OBSTRUCTION\_VERIFIED}\):
it imports the orientation-character definition, the
\((O2_{\mathrm{atlas}})\) wall-atlas ledger, the finite Pfaffian
ledger, and the automorphic Maass character as a firewall; it records
the rank-one formula for row \(320\); and it keeps the wall-object,
rank-one chart, invariant-unit, coordinate-flip, local-sign,
generator-sign, stratum-independence, and transition-compatibility
rows open.

\begin{proposition}[Determinant criterion for the orientation
character]
\label{prop:orientation-character-determinant-criterion}
\textnormal{[group-theoretic criterion proved; source generator signs
not supplied]}
Let
\[
\det_W:W^{(2)}(\La^{2,1}_{II})\longrightarrow\{\pm1\}
\]
be the determinant character of the reflection representation on
\(\La^{2,1}_{II}\otimes_{\mathbb Z}\mathbb R\).  Suppose the
source-side finite-stage character \(\epsilon_{o,R}\) of
Definition~\ref{def:orientation-character-epsilon-o} has been
constructed and the row \(320\) finite sign rows prove
\[
\epsilon_{o,R}(s_{\delta_i})=-1,\qquad i=1,2,3.
\]
Then
\[
\epsilon_{o,R}=\det_W\quad\text{on}\quad
W^{(2)}(\La^{2,1}_{II}).
\]
If the finite characters are transition-compatible, then the inverse
limit character satisfies
\[
\epsilon_o=\det_W\quad\text{on}\quad
W^{(2)}(\La^{2,1}_{II}).
\]

The present finite data do not prove row \(321\).  The row \(319\)
packet does not construct \(\epsilon_{o,R}\), and the row \(320\)
packet contains no generator-sign rows.  The automorphic packet records
\(\nu_{\Delta_5}|_{W^{(2)}}=\det_W\), but that is a target Maass
character computation, not a source orientation-character computation.
\end{proposition}

\begin{proof}
Each \(s_{\delta_i}\) is a real reflection in the retained
rank-three Lorentzian reflection representation, hence
\[
\det_W(s_{\delta_i})=-1,\qquad i=1,2,3.
\]
The type-II Coxeter presentation used in
Definition~\ref{def:orientation-character-epsilon-o} has generators
\(s_{\delta_1},s_{\delta_2},s_{\delta_3}\), relations
\(s_{\delta_i}^2=1\), and no braid relation between distinct
generators.  Therefore a \(\{\pm1\}\)-valued character is determined
by its three generator values.  If row \(320\) supplies the three
source values \(\epsilon_{o,R}(s_{\delta_i})=-1\), then
\(\epsilon_{o,R}\) and \(\det_W\) agree on the generators and hence on
all words in \(W^{(2)}(\La^{2,1}_{II})\).  The same argument passes to
the inverse limit once the transition-compatibility rows for
\(\epsilon_{o,R}\) are supplied.
\end{proof}

The packet
\[
\texttt{certificates/orientation/rhomred\_weyl\_orientation\_determinant\_comparison}
\]
verifies
\(\texttt{RHOMRED\_WEYL\_ORIENTATION\_DETERMINANT\_COMPARISON\_OBSTRUCTION\_VERIFIED}\):
it imports the row \(319\) orientation-character definition, the row
\(320\) simple-sign criterion, and the automorphic determinant
comparison as a firewall; it records the determinant criterion for row
\(321\); and it keeps the source character, source generator signs,
equality-on-generators, transition compatibility, and inverse-limit
comparison rows open.

\begin{proposition}[Maass-character criterion for the orientation
character]
\label{prop:orientation-character-maass-criterion}
\textnormal{[comparison criterion proved; source determinant comparison
not supplied]}
Assume the row \(321\) source comparison has been supplied:
\[
\epsilon_o=\det_W\quad\text{on}\quad
W^{(2)}(\La^{2,1}_{II}).
\]
Assume also the automorphic Maass character computation for the
Borcherds--Gritsenko--Nikulin square root,
\[
\nu_{\Delta_5}\big|_{W^{(2)}(\La^{2,1}_{II})}
=\det_W .
\]
Then
\[
\epsilon_o=\nu_{\Delta_5}\big|_{W^{(2)}(\La^{2,1}_{II})}.
\]

The present finite data do not prove row \(322\).  The Maass character
packet proves the automorphic equality
\(\nu_{\Delta_5}|_{W^{(2)}}=\det_W\), but the row \(321\) packet does
not prove the source equality \(\epsilon_o=\det_W\).  Therefore the
Maass character cannot be used as a substitute for the missing
source-side orientation character or the missing type-II wall sign
rows.
\end{proposition}

\begin{proof}
The assertion is the transitivity of equality of characters on
\(W^{(2)}(\La^{2,1}_{II})\).  If the source orientation character has
already been identified with \(\det_W\), and the automorphic Maass
character restricts to the same determinant character, then both
characters have identical values on every element of
\(W^{(2)}(\La^{2,1}_{II})\).  The second equality is automorphic target
data.  The first equality is the source-side obstruction in row \(321\)
and is not supplied by the current packets.
\end{proof}

The packet
\[
\texttt{certificates/orientation/rhomred\_weyl\_orientation\_maass\_comparison}
\]
verifies
\(\texttt{RHOMRED\_WEYL\_ORIENTATION\_MAASS\_COMPARISON\_OBSTRUCTION\_VERIFIED}\):
it imports the row \(321\) determinant-comparison criterion and the
automorphic Maass-character packet; it records the transitivity
criterion for row \(322\); and it keeps the source determinant
comparison, inverse-limit source character, and orientation-versus-Maass
comparison rows open.

\begin{proposition}[\(S_3\)-chamber automorphisms have trivial Maass
character]
\label{prop:s3-chamber-maass-triviality}
\textnormal{[proved; automorphic character only]}
The restriction of the Maass character of the Igusa square root to the
type-I chamber automorphism group is trivial:
\[
\nu_{\Delta_5}\big|_{\Aut(\Poly_{II})}=1,\qquad
\Aut(\Poly_{II})\simeq S_3.
\]
Equivalently,
\[
\nu_{\Delta_5}(s_{f_{-2}-f_2})
=\nu_{\Delta_5}(s_{f_2-f_3})
=\nu_{\Delta_5}(s_{f_{-2}-f_3})=+1 .
\]
This is an automorphic multiplier statement.  It does not decide
whether \(\epsilon_o\) extends from
\(W^{(2)}(\La^{2,1}_{II})\) to
\(W^{(2)}(\La^{2,1}_{II})\rtimes S_3\); that is a separate
source-side orientation-transport problem.
\end{proposition}

\begin{proof}
Lemma~\ref{lem:bkm-chamber-s3} identifies
\(\Aut(\Poly_{II})\simeq S_3\) and shows that the three type-I
reflections
\[
s_{f_{-2}-f_2},\qquad s_{f_2-f_3},\qquad s_{f_{-2}-f_3}
\]
act as the three transpositions of the type-II walls.  These
transpositions generate \(S_3\).  The Maass multiplier computation of
Lemma~\ref{lem:appE-maass-on-type-ii}, equivalently the automorphic
packet
\(\texttt{certificates/automorphic/delta5\_maass\_character}\), gives
value \(+1\) on each of the three type-I chamber reflections.  Since
\(\nu_{\Delta_5}\) is a character, it is \(+1\) on every product of
these generators, hence trivial on \(\Aut(\Poly_{II})\).
\end{proof}

The packet
\[
\texttt{certificates/automorphic/delta5\_s3\_chamber\_maass\_triviality}
\]
verifies
\(\texttt{DELTA5\_S3\_CHAMBER\_MAASS\_TRIVIALITY\_VERIFIED}\):
it imports the Maass-character packet, records the three type-I
chamber transpositions and their \(+1\) values, checks the
chamber-automorphism relation row, and excludes any conclusion about
the \(S_3\)-extension of the compact orientation character.

\begin{proposition}[Semidirect extension decision for the orientation
character]
\label{prop:orientation-character-s3-extension-decision}
\textnormal{[abstract character criterion proved; source semidirect
Hall lift not supplied]}
Let \(W=W^{(2)}(\La^{2,1}_{II})\) and
\(\Sigma=\Aut(\Poly_{II})\simeq S_3\).  The group \(\Sigma\) acts on
the simple reflections of \(W\) by permutation.  Suppose the
source-side orientation character \(\epsilon_o:W\to\{\pm1\}\) has been
constructed and satisfies
\[
\epsilon_o(s_{\delta_1})=\epsilon_o(s_{\delta_2})
=\epsilon_o(s_{\delta_3})=-1.
\]
Then \(\epsilon_o\) is \(\Sigma\)-invariant and therefore extends as
an abstract character to \(W\rtimes\Sigma\).  The extensions are
\[
\widetilde\epsilon_{\chi}(w,\sigma)=\epsilon_o(w)\chi(\sigma),
\qquad
\chi\in\Hom(\Sigma,\{\pm1\})=\{1,\operatorname{sgn}\}.
\]
The unique extension whose restriction to \(\Sigma\) equals the Maass
character of row \(323\) is the trivial-\(\Sigma\) extension
\[
\widetilde\epsilon_{\mathrm{Maass}}(w,\sigma)=\epsilon_o(w).
\]

This is only the abstract character decision.  A compact-source
extension of the orientation character requires semidirect Hall lifts
of the chamber-permuting type-I reflections, their action on the
orientation line, determinant square-root compatibility, preservation
of the quotient-orientation null-trivialisations, and the semidirect
relations.  The present finite packets do not supply those rows.
\end{proposition}

\begin{proof}
The chamber action of Lemma~\ref{lem:bkm-chamber-s3} sends
\(s_{\delta_i}\) to \(s_{\delta_{\sigma(i)}}\).  Hence, under the
displayed hypothesis,
\[
\epsilon_o(\sigma s_{\delta_i}\sigma^{-1})
=\epsilon_o(s_{\delta_{\sigma(i)}})
=-1
=\epsilon_o(s_{\delta_i})
\]
for every generator \(s_{\delta_i}\) and every
\(\sigma\in\Sigma\).  Since the type-II Coxeter presentation has only
the involutive relations \(s_{\delta_i}^2=1\) and no braid relations,
this equality on the generators gives \(\Sigma\)-invariance of
\(\epsilon_o\).  A character on \(W\) extends to the semidirect product
precisely after this invariance condition holds; any two extensions
differ by a character of \(\Sigma\).  Since
\(\Hom(S_3,\{\pm1\})=\{1,\operatorname{sgn}\}\), there are exactly two
abstract extensions.  Proposition~\ref{prop:s3-chamber-maass-triviality}
shows that \(\nu_{\Delta_5}\) is trivial on \(\Sigma\), so the
Maass-compatible abstract extension is the one with
\(\chi=1\).
\end{proof}

The packet
\[
\texttt{certificates/orientation/rhomred\_weyl\_orientation\_s3\_extension\_decision}
\]
verifies
\(\texttt{RHOMRED\_WEYL\_ORIENTATION\_S3\_EXTENSION\_DECISION\_OBSTRUCTION\_VERIFIED}\):
it imports the determinant-comparison criterion and the \(S_3\) Maass
triviality packet; it records the two abstract semidirect character
extensions and the Maass-compatible choice; and it keeps the
source-side semidirect Hall lifts, orientation-line transport,
square-root compatibility, quotient-orientation transport, semidirect
relations, and compact-source extension rows open.

\begin{proposition}[Semidirect Hall-lift construction criterion]
\label{prop:semidirect-hall-lift-construction-criterion}
\textnormal{[criterion proved; semidirect Hall lifts not supplied]}
Let \(W=W^{(2)}(\La^{2,1}_{II})\) and
\(\Sigma=\Aut(\Poly_{II})\simeq S_3\).  A compact-source lift of the
abstract extension of Proposition~\ref{prop:orientation-character-s3-extension-decision}
is a finite-stage action of \(W\rtimes\Sigma\) on the retained compact
Hall correspondence package.  It consists of:
\begin{enumerate}
\item the populated type-II partial action groupoid for \(W\), with
orientation-line arrows \(\tau_{i,R,S}\) and zero projective Coxeter
defects;
\item for the three chamber transpositions \(\sigma\in\Sigma\), Hall
correspondences
\[
\mathfrak H_{\sigma,R,S}:
\mathcal M_{R,S}\longleftarrow
\mathcal Z_{\sigma,R,S}\longrightarrow
\mathcal M_{R,\sigma S}
\]
over the retained quotient stacks, together with source and target
maps on the Hall product, Hall coproduct, unit, counit, primitive
subspace, pairing, radical, and PBW filtration rows;
\item Picard-groupoid isomorphisms
\[
\theta^o_{\sigma,R,S}:o_{R,S}\xrightarrow{\sim}
\sigma^*o_{R,\sigma S}
\]
whose squares are the determinant-line transports;
\item transport of
\(\alpha^{\mathrm{red}}\), \(\alpha^{E,\mathrm{free}}\),
\(\widetilde\beta^H\), \(\beta^H\), and \(\lambda^H=0\) through the
same chamber correspondences;
\item zero-defect rows for the \(S_3\) relations and for the
semidirect conjugation relations
\[
\theta^o_{\sigma,R,\delta_i}\,\tau_{i,R,S}\,
(\theta^o_{\sigma,R,\delta_i})^{-1}
=\tau_{\sigma(i),R,\sigma S};
\]
\item transition compatibility in \(R\) and \(R^1\varprojlim=0\) for
the semidirect Hall and Picard-groupoid data.
\end{enumerate}
If these rows are supplied, the trivial-\(S_3\) abstract extension of
row \(324\) is represented by compact-source orientation transport on
\(W\rtimes S_3\).  The present finite packets do not supply these
rows: the compact Hall source packet is empty blocked, the type-II
partial action groupoid has no object or arrow rows, the \(\tau_i\)
packet has no orientation-line transport rows, and the transition
orientation packet has no Picard-groupoid transition rows.
\end{proposition}

\begin{proof}
The listed data are exactly the data of an action of the semidirect
product in the oriented Hall correspondence 2-category.  Items \(1\)
and \(2\) give the underlying source correspondences on the retained
Hall stacks.  Item \(3\) lifts the action from line bundles to
Joyce--Upmeier orientations by enforcing the determinant square-root
diagram.  Item \(4\) says that the quotient-orientation structure is
transported with the orientation line.  Item \(5\) is the group law:
the \(S_3\) relations define the chamber part, while the conjugation
relations identify the chamber transport of the type-II Weyl lifts.
Item \(6\) is the inverse-system condition needed to pass from finite
height to the pro-object.  Without these rows there is only an abstract
character extension, not a compact-source Hall action.
\end{proof}

The packet
\[
\texttt{certificates/orientation/rhomred\_weyl\_semidirect\_hall\_lifts}
\]
verifies
\(\texttt{RHOMRED\_WEYL\_SEMIDIRECT\_HALL\_LIFTS\_OBSTRUCTION\_VERIFIED}\):
it imports the row \(324\) extension decision, the compact Hall source
obstruction ledger, the type-II partial-action-groupoid ledger, the
\(\tau_i\)-transport ledger, and the transition-orientation ledger; it
records the semidirect Hall-lift construction criterion; and it keeps
the source Hall correspondences, orientation transports, determinant
square-root rows, quotient-orientation transports, semidirect relation
rows, transition rows, and inverse-limit rows open.

\begin{proposition}[Imaginary divisor monodromy is not Weyl monodromy]
\label{prop:imaginary-divisor-not-weyl-monodromy}
\textnormal{[firewall proved; local divisor monodromy not yet defined]}
The orientation character
\[
\epsilon_o:W^{(2)}(\La^{2,1}_{II})\longrightarrow\{\pm1\}
\]
is a real-root reflection character.  Its domain is generated by the
three real simple reflections
\[
s_{\delta_1},\qquad s_{\delta_2},\qquad s_{\delta_3}.
\]
An imaginary simple root of the Borcherds--Kac--Moody denominator
algebra contributes a multiplicity in the denominator formula; it does
not define an element of \(W^{(2)}(\La^{2,1}_{II})\).  Hence monodromy
around a divisor component indexed by an imaginary root, if such
monodromy is constructed in the compact Hall theory, is not a Weyl
character value and is not a value of \(\epsilon_o\).

Thus no imaginary-divisor monodromy is being called Weyl monodromy in
the orientation-character construction.  Row \(327\) defines local
divisor monodromy as a separate datum; row \(328\) proves that this
separate datum is irrelevant to the Maass character
\(\nu_{\Delta_5}\) on \(W^{(2)}(\La^{2,1}_{II})\).
\end{proposition}

\begin{proof}
Definition~\ref{def:bkm-real-roots} defines the real-root system as
the \(W^{(2)}(\La^{2,1}_{II})\)-orbit of
\(\{\delta_1,\delta_2,\delta_3\}\), and the type-II Weyl group is
generated by the associated real simple reflections.  The imaginary
simple roots are the Borcherds correction recorded in
Table~\ref{tab:bkm-real-root-comparison}: they are absent from the
rank-three Kac--Moody real-root datum and enter \(\fg_{\Delta_5}\) by
the multiplicities \(m(a)\) and \(\tau(a)\) in the denominator
formula.  A Weyl reflection is attached to a real root of norm \(2\);
an imaginary simple root contributes a denominator factor and a root
space, not a reflection generator.  Therefore the notation
\(\epsilon_o(w)\) is meaningful only for
\(w\in W^{(2)}(\La^{2,1}_{II})\), and no loop around an imaginary
divisor component can be read as such a value without adding a new
definition.  That new definition is isolated in row \(327\).
\end{proof}

The packet
\[
\texttt{certificates/orientation/rhomred\_weyl\_imaginary\_divisor\_firewall}
\]
verifies
\(\texttt{RHOMRED\_WEYL\_IMAGINARY\_DIVISOR\_FIREWALL\_VERIFIED}\):
it imports the real-root definition, the BKM real/imaginary comparison,
the Chapter~6 imaginary-root firewall, and the row \(325\) lift
criterion; it records that real-root reflections are the only Weyl
generators; and it keeps local divisor monodromy outside the Weyl
orientation character until row \(327\) defines it.

\begin{definition}[Local divisor monodromy datum]
\label{def:local-divisor-monodromy-datum}
\textnormal{[definition specified; local divisor rows not supplied]}
Fix a finite height \(R\) and a retained compact Hall stratum
\(\mathcal M_{R,S}\).  A local divisor monodromy datum with label
\(\eta\) consists of:
\begin{enumerate}
\item a retained effective Cartier divisor or normal-crossing
component
\[
\mathcal D_{R,S,\eta}\subset \mathcal M_{R,S}
\]
together with the row identifying the source of the label \(\eta\);
if \(\eta\) is an imaginary Borcherds label, the label is only a
divisor or multiplicity label and not a Weyl-reflection generator;
\item an analytic or formal normal tube \(N_{R,S,\eta}\) of
\(\mathcal D_{R,S,\eta}\) in \(\mathcal M_{R,S}\), its punctured tube
\[
N^\times_{R,S,\eta}:=
N_{R,S,\eta}\setminus \mathcal D_{R,S,\eta},
\]
and a base point \(x_{R,S,\eta}\in N^\times_{R,S,\eta}\);
\item a positively oriented meridian class
\[
\gamma_{R,S,\eta}\in
\pi_1(N^\times_{R,S,\eta},x_{R,S,\eta});
\]
\item the oriented reduced coefficient system
\[
K_R:=\Phi^{\mathrm{red}}_R\otimes o_R
\]
restricted to \(N^\times_{R,S,\eta}\), and, when the Pfaffian line is
separately supplied, the corresponding rank-one Pfaffian-line factor;
\item the monodromy representation of this local system and the
operator
\[
T_{R,S,\eta}:=
\rho_{K_R}(\gamma_{R,S,\eta})
\in \Aut\bigl((K_R)_{x_{R,S,\eta}}\bigr).
\]
When a rank-one Pfaffian-line summand is supplied and preserved by
the meridian action, its sign is denoted
\(\mu_{R,S,\eta}\in\{\pm1\}\);
\item for every successor \(R'\le R\), transition rows identifying
the divisor components, normal tubes, meridians, coefficient systems,
and monodromy operators.
\end{enumerate}
The finite-stage local divisor monodromy is the partially defined
assignment
\[
(R,S,\eta)\longmapsto T_{R,S,\eta}
\]
on the supplied rows.  Its domain is a set of local compact Hall
divisor data, not \(W^{(2)}(\Lambda^{2,1}_{II})\).  It is not a value
of \(\epsilon_o\), not a value of the Maass character
\(\nu_{\Delta_5}\), and not determined by the Borcherds
multiplicity attached to \(\eta\).  Without the divisor component,
normal tube, meridian, oriented coefficient system, and transition
rows above, no local divisor monodromy value is defined.
\end{definition}

The packet
\[
\texttt{certificates/orientation/rhomred\_local\_divisor\_monodromy\_definition}
\]
verifies
\(\texttt{RHOMRED\_LOCAL\_DIVISOR\_MONODROMY\_DEFINITION\_OBSTRUCTION\_VERIFIED}\):
it imports the row \(326\) firewall, the compact Hall obstruction
ledger, the reduced-orientation obstruction ledger, the reduced
vanishing-cycle transition ledger, and the finite Pfaffian obstruction
ledger; it records Definition~\ref{def:local-divisor-monodromy-datum};
and it keeps the divisor-component, normal-tube, meridian,
coefficient-system, Pfaffian-line, monodromy-operator, and transition
rows open.

\begin{proposition}[Local divisor monodromy is irrelevant to the
Maass character]
\label{prop:local-divisor-monodromy-maass-irrelevance}
\textnormal{[proved as a domain-separation theorem]}
Let \(\nu_{\Delta_5}\) be the Maass character of the automorphic
Igusa square root, with values on the type-II Weyl generators and on
the \(S_3\) chamber automorphisms as recorded by the automorphic
Maass-character packet.  Let
\((R,S,\eta)\mapsto T_{R,S,\eta}\) be any supplied local divisor
monodromy datum in the sense of
Definition~\ref{def:local-divisor-monodromy-datum}.  Then the values
of \(\nu_{\Delta_5}\) on
\[
W^{(2)}(\Lambda^{2,1}_{II})\quad\text{and}\quad
\Aut(\Poly_{II})\simeq S_3
\]
are independent of the operators \(T_{R,S,\eta}\).

More precisely, there is no expression
\[
\nu_{\Delta_5}(\gamma_{R,S,\eta})
\]
in the present data.  Such an expression would require an additional
comparison map from the local meridian groupoid of the compact Hall
divisor rows to the automorphic or Weyl group on which
\(\nu_{\Delta_5}\) is defined.  No such comparison map is part of
Definition~\ref{def:local-divisor-monodromy-datum}, and no such row is
supplied by the current packets.
\end{proposition}

\begin{proof}
The Maass character is automorphic target data.  Its recorded inputs
are the type-II real-root reflections \(s_{\delta_i}\) and the
chamber automorphisms in \(\Aut(\Poly_{II})\); the packet
\(\texttt{certificates/automorphic/delta5\_maass\_character}\) records
these values and contains no compact Hall source rows.  By
Definition~\ref{def:local-divisor-monodromy-datum}, local divisor
monodromy is instead a meridian action on the restricted oriented
coefficient system \(K_R=\Phi_R^{\mathrm{red}}\otimes o_R\) over a
punctured normal tube of a retained compact Hall divisor component.
Its input is the meridian class \(\gamma_{R,S,\eta}\), not an element
of \(W^{(2)}(\Lambda^{2,1}_{II})\) and not a chamber automorphism.

Therefore the Maass-character computation has no argument on which a
local divisor monodromy value can be evaluated.  Changing or supplying
operators \(T_{R,S,\eta}\) while keeping the automorphic Maass packet
fixed changes only source-side local-system data; it does not change
the automorphic multiplier rows defining \(\nu_{\Delta_5}\).  A
comparison would require a new meridian-to-automorphic map and a proof
that the Maass character pulls back along it.  Those rows are absent.
\end{proof}

The packet
\[
\texttt{certificates/orientation/rhomred\_local\_divisor\_monodromy\_maass\_irrelevance}
\]
verifies
\(\texttt{RHOMRED\_LOCAL\_DIVISOR\_MONODROMY\_MAASS\_IRRELEVANCE\_VERIFIED}\):
it imports Definition~\ref{def:local-divisor-monodromy-datum}, the row
\(322\) Maass-comparison criterion, and the automorphic
Maass-character packet; it records the domain-separation theorem; and
it keeps any meridian-to-automorphic comparison map, local monodromy
operator rows, and row \(329\) simple-wall orientation computations
separate.

\begin{proposition}[Orientation on the three simple wall charts]
\label{prop:three-simple-wall-chart-orientation}
\textnormal{[chart-level computation criterion proved; finite chart
rows not supplied]}
For \(i=1,2,3\), let a retained \((O2_{\mathrm{atlas}})\) chart over
the simple type-II wall \(\Hpl_{\delta_i}\) supply the data
\[
(U_{i,R,S},T_{i,R,S},u_i,\mathscr L^{\mathrm{tan}}_{i,R,S},
\upsilon_{i,R,S},o_{i,R,S})
\]
with
\[
U_{i,R,S}=T_{i,R,S}\times(-\varepsilon,\varepsilon),\qquad
s_{\delta_i}(t,u_i)=(t,-u_i),
\]
rank-one normal Pfaffian block
\[
K^{\mathrm{cross}}_{i,R,S}=[\mathbb R\xrightarrow{u_i}\mathbb R],
\qquad
\operatorname{pf}_{i,R,S}^{\mathrm{loc}}=\upsilon_{i,R,S}u_i,
\]
divisor order \(1\), and invariant tangent Pfaffian unit
\[
s_{\delta_i}(\upsilon_{i,R,S})=\upsilon_{i,R,S}.
\]
Then the source orientation sign on each of the three supplied simple
wall charts is
\[
\omega^{\mathrm{chart}}_{i,R,S}
:=\frac{s_{\delta_i}^*\operatorname{pf}_{i,R,S}^{\mathrm{loc}}}
        {\operatorname{pf}_{i,R,S}^{\mathrm{loc}}}
=-1,\qquad i=1,2,3.
\]
Equivalently, the three-chart orientation vector is
\[
(\omega^{\mathrm{chart}}_{1,R,S},
 \omega^{\mathrm{chart}}_{2,R,S},
 \omega^{\mathrm{chart}}_{3,R,S})
=(-1,-1,-1)
\]
on the supplied rows.

The present finite packets do not supply the three rows to which this
calculation applies.  The \((O2_{\mathrm{atlas}})\) packet has no
retained wall objects and no wall-chart rows; the finite Pfaffian
packet has no type-II wall chart, Pfaffian line, invariant unit, or
local sign row; and the row \(320\) packet contains only the abstract
rank-one sign criterion, not populated chart-orientation rows.
\end{proposition}

\begin{proof}
The computation is componentwise.  On the \(i\)-th supplied rank-one
chart the reflection fixes the tangent factor and sends the normal
coordinate to its negative.  Hence
\[
s_{\delta_i}^*(\upsilon_{i,R,S}u_i)
=\upsilon_{i,R,S}(-u_i)
=-\upsilon_{i,R,S}u_i.
\]
The quotient of the transported local Pfaffian section by the original
section is therefore \(-1\).  The same calculation holds for
\(i=1,2,3\).  The assertion is a source chart-orientation computation:
it uses the normal coordinate, invariant tangent unit, divisor order,
and retained quotient orientation on the chart.  By
Proposition~\ref{prop:local-divisor-monodromy-maass-irrelevance}, no
Maass-character value or local divisor monodromy datum supplies these
rows.
\end{proof}

The packet
\[
\texttt{certificates/orientation/rhomred\_three\_simple\_wall\_chart\_orientation}
\]
verifies
\(\texttt{RHOMRED\_THREE\_SIMPLE\_WALL\_CHART\_ORIENTATION\_OBSTRUCTION\_VERIFIED}\):
it imports the row \(320\) simple-wall sign criterion, the
\((O2_{\mathrm{atlas}})\) wall-atlas ledger, the finite Pfaffian
ledger, and the row \(328\) Maass-irrelevance firewall; it records the
three-chart orientation computation criterion for row \(329\); and it
keeps the three retained chart rows, invariant-unit rows, quotient
orientation rows, local sign rows, and transition rows open.

\begin{proposition}[Orientation on the isotropic \(a_{ij}\)-charts]
\label{prop:isotropic-aij-chart-orientation}
\textnormal{[chart-level parity criterion proved; source charts not
supplied]}
For \(1\le i<j\le3\), let \(a_{ij}:=\delta_i+\delta_j\) be the
primitive isotropic boundary ray of the type-II chamber.  Suppose a
retained compact Hall chart \(\mathcal U_{ij,R,S}\) in degree
\(a_{ij}\) supplies:
\begin{enumerate}
\item a source decomposition of the oriented reduced coefficient block
into \(10\) even and \(0\) odd one-dimensional summands,
compatible with the target \(a_{ij}:10|0\) row;
\item the quotient-orientation line and its determinant-square row on
\(\mathcal U_{ij,R,S}\);
\item identification of the single affine \(A_1^{(1)}\) null-root
direction and the nine Gritsenko--Nikulin isotropic directions
\(u_{ij,1},\ldots,u_{ij,9}\);
\item transition compatibility for \(R'\le R\).
\end{enumerate}
Then the local isotropic chart orientation parity is even and the
chart sign is
\[
\omega^{\mathrm{iso}}_{ij,R,S}
=(-1)^{\dim K^{\mathrm{iso}}_{\overline1,ij,R,S}}
=(-1)^0=+1 .
\]
Thus the three supplied isotropic chart rows would give
\[
(\omega^{\mathrm{iso}}_{12,R,S},
 \omega^{\mathrm{iso}}_{13,R,S},
 \omega^{\mathrm{iso}}_{23,R,S})
=(+1,+1,+1).
\]

The present packets do not supply these source chart rows.  The
\(\texttt{type\_ii\_chamber\_isotropic}\) and
\(\texttt{wle3\_target\_parity}\) packets certify the target arithmetic
\(a_{ij}:10|0\), the decomposition \(10=1+9\), and the three
isotropic rays \(a_{12},a_{13},a_{23}\); they do not construct compact
Hall source charts, determinant-square orientation rows, or transition
rows.  Consequently row \(330\) is presently a chart-orientation
criterion, not an unconditional source orientation computation.
\end{proposition}

\begin{proof}
The target \(W_{\le3}\) parity table records
\[
a_{ij}:10|0,\qquad i<j,
\]
and the chamber-isotropic packet records the decomposition of each
primitive ray into the affine null-root contribution \(1\) and the
Gritsenko--Nikulin isotropic contribution \(9\).  If a source compact
Hall chart realizes these ten directions as even oriented summands and
supplies no odd summand, the parity exponent of the local determinant
orientation is the odd rank.  This rank is \(0\).  Hence the sign is
\((-1)^0=+1\) for each of the three isotropic charts.  The computation
uses source chart and orientation rows; the target parity table alone
does not provide them.
\end{proof}

The packet
\[
\texttt{certificates/orientation/rhomred\_isotropic\_aij\_chart\_orientation}
\]
verifies
\(\texttt{RHOMRED\_ISOTROPIC\_AIJ\_CHART\_ORIENTATION\_OBSTRUCTION\_VERIFIED}\):
it imports the type-II chamber-isotropic packet, the \(W_{\le3}\)
target-parity packet, the row \(329\) simple-wall chart-orientation
packet, and the compact Hall source obstruction ledger; it records the
isotropic chart-orientation criterion for row \(330\); and it keeps
the compact source chart, quotient-orientation, determinant-square,
isotropic-source parity, and transition rows open.

\begin{proposition}[Orientation on the \(\delta_{123}\)-chart]
\label{prop:delta123-chart-orientation}
\textnormal{[chart-level parity criterion proved; source chart not
supplied]}
Let
\[
\delta_{123}:=\delta_1+\delta_2+\delta_3.
\]
Suppose a retained compact Hall chart \(\mathcal U_{123,R,S}\) in
degree \(\delta_{123}\) supplies:
\begin{enumerate}
\item a source decomposition of the oriented reduced coefficient block
into \(29\) even and \(93\) odd one-dimensional summands, compatible
with the target \(\delta_{123}:29|93\) row;
\item source rows realizing the two real--real--real directions, the
\(27\) mixed real--isotropic directions, and the \(93\) odd
imaginary-simple directions of
Proposition~\ref{prop:bkm-delta123-presentation-count};
\item the quotient-orientation line and its determinant-square row on
\(\mathcal U_{123,R,S}\);
\item transition compatibility for \(R'\le R\).
\end{enumerate}
Then the local first-timelike chart orientation parity is odd and the
chart sign is
\[
\omega^{\mathrm{time}}_{123,R,S}
=(-1)^{\dim K^{\mathrm{time}}_{\overline1,123,R,S}}
=(-1)^{93}=-1 .
\]

The present packets do not supply the source chart to which this
calculation applies.  The
\(\texttt{delta123\_presentation\_split}\) and
\(\texttt{wle3\_target\_parity}\) packets certify the target
presentation count \(29|93\), with \(29=2+27\) and \(29-93=-64\);
they do not construct a compact Hall source chart, a source parity
decomposition, a quotient orientation, a determinant-square row, or
transition rows.  Consequently row \(331\) is presently a
chart-orientation criterion, not an unconditional source orientation
computation.
\end{proposition}

\begin{proof}
Proposition~\ref{prop:bkm-delta123-presentation-count} and the finite
target packet
\(\texttt{certificates/targets/delta5\_gn\_kac/delta123\_presentation\_split}\)
compute the first timelike target presentation:
\[
\rootmult_{\overline0}(\delta_{123})=29,\qquad
\rootmult_{\overline1}(\delta_{123})=93,\qquad
29-93=-64.
\]
The even part is the direct presentation count \(2+27\): two
real--real--real directions after the Jacobi relation and twenty-seven
mixed real--isotropic brackets.  The odd part consists of the
ninety-three negative-norm imaginary-simple directions.  If a compact
Hall chart realizes this decomposition as the source reduced
coefficient block and supplies the determinant-square orientation row,
the parity exponent of the local determinant orientation is the odd
rank.  That rank is \(93\).  Hence the chart sign is
\((-1)^{93}=-1\).  The target presentation count alone is not a source
orientation row.
\end{proof}

The packet
\[
\texttt{certificates/orientation/rhomred\_delta123\_chart\_orientation}
\]
verifies
\(\texttt{RHOMRED\_DELTA123\_CHART\_ORIENTATION\_OBSTRUCTION\_VERIFIED}\):
it imports the target \(\delta_{123}\) presentation-split packet, the
\(W_{\le3}\) target-parity packet, the compact Hall source obstruction
ledger, the finite Pfaffian obstruction ledger, and the row \(330\)
isotropic chart-orientation firewall; it records the
\(\delta_{123}\) chart-orientation criterion for row \(331\); and it
keeps the compact source chart, source parity, source presentation
decomposition, quotient-orientation, determinant-square, and transition
rows open.

\begin{proposition}[Orientation-line compatibility with the \(29|93\)
parity split]
\label{prop:orientation-line-delta123-parity-compatibility}
\textnormal{[line-level compatibility criterion proved; source parity
and orientation line not supplied]}
Let \(\mathcal U_{123,R,S}\) be a retained compact Hall chart in degree
\(\delta_{123}\), and let
\[
K^{\mathrm{red}}_{123,R,S}
\]
be the corresponding reduced coefficient block.  Suppose the following
source data are supplied:
\begin{enumerate}
\item a parity involution \(J_{123,R,S}\) on
\(K^{\mathrm{red}}_{123,R,S}\), with eigenspace decomposition
\[
K^{\mathrm{red}}_{123,R,S}
=K_{\overline0,123,R,S}\oplus K_{\overline1,123,R,S},
\qquad
\dim K_{\overline0,123,R,S}=29,\quad
\dim K_{\overline1,123,R,S}=93;
\]
\item source comparison rows identifying the even summand with the two
real--real--real directions and the \(27\) mixed real--isotropic
directions, and identifying the odd summand with the \(93\) odd
imaginary-simple directions of the target
\(\delta_{123}\)-presentation;
\item a Joyce--Upmeier orientation line \(o_{123,R,S}\) and a commuting
determinant-square diagram
\[
o_{123,R,S}^{\otimes2}
\xrightarrow{\ \sim\ }
\det K^{\mathrm{red}}_{123,R,S}
\xrightarrow{\ \sim\ }
\det K_{\overline0,123,R,S}\otimes
\bigl(\det K_{\overline1,123,R,S}\bigr)^{-1};
\]
\item quotient-orientation null-trivialisations on the chart, compatible
with the parity splitting;
\item transition compatibility: for every \(R'\le R\), the transition
map commutes with \(J_{123}\), preserves the determinant-square diagram
above, and has zero Picard-groupoid, parity-commutator, off-diagonal,
composition, and \(R^1\varprojlim\) defects.
\end{enumerate}
Then the orientation line is compatible with the \(29|93\) parity split.
In this case the determinant parity exponent is the source odd rank
\(93\), so the local orientation sign on the first timelike chart is
\[
(-1)^{\dim K_{\overline1,123,R,S}}=(-1)^{93}=-1.
\]

The present packets do not supply these hypotheses.  The target
\(\delta_{123}:29|93\) presentation is verified, and
Proposition~\ref{prop:delta123-chart-orientation} records the local
chart criterion, but the compact Hall source packet has no parity block
rows, the reduced-orientation packet has no orientation-line
square-root rows, the transition-parity packet has no parity involution
or block-diagonal transition rows, and the transition-orientation packet
has no Picard-groupoid transport rows.  Hence row \(332\) is presently
a compatibility criterion and obstruction ledger, not a proof that the
orientation line is compatible with the \(29|93\) split.
\end{proposition}

\begin{proof}
The determinant functor on a \(\mathbb Z/2\)-graded finite block gives
the Berezinian factorisation
\[
\det K^{\mathrm{red}}_{123,R,S}
\simeq
\det K_{\overline0,123,R,S}\otimes
\bigl(\det K_{\overline1,123,R,S}\bigr)^{-1}.
\]
A Joyce--Upmeier orientation on the chart is a square root of this
determinant line in the Picard groupoid, not a scalar dimension.  If the
displayed square diagram commutes, the chosen orientation line is an
orientation of the parity-split determinant line itself.  The target
presentation count \(29|93\) identifies the required ranks; the source
comparison rows in item \(2\) assert that these target ranks are the
actual source parity ranks on \(\mathcal U_{123,R,S}\).  The local sign
is then the parity of the odd determinant factor, namely
\((-1)^{93}=-1\).  The transition rows in item \(5\) make the statement
independent of the finite HN height and prevent a parity split at one
stage from being read as a pro-object orientation.
\end{proof}

The packet
\[
\texttt{certificates/orientation/rhomred\_delta123\_parity\_compatibility}
\]
verifies
\(\texttt{RHOMRED\_DELTA123\_PARITY\_COMPATIBILITY\_OBSTRUCTION\_VERIFIED}\):
it imports the row \(331\) \(\delta_{123}\)-chart criterion, the
reduced determinant-line packet, the square-root obstruction ledger, the
compact Hall source obstruction ledger, the transition-parity
obstruction ledger, and the transition-orientation obstruction ledger;
it records the line-level compatibility criterion for row \(332\); and
it keeps the source parity block, parity involution,
source-to-target comparison, orientation square-root,
determinant-square, quotient-orientation, parity-transition, and
Picard-groupoid transition rows open.

\begin{proposition}[Orientation-line compatibility with negative roots]
\label{prop:orientation-line-negative-root-compatibility}
\textnormal{[negative-root compatibility criterion proved; source
pairing and orientation transport not supplied]}
Let \(\gamma\) be a retained positive degree whose target row carries a
formal negative-dual block, and let \(-\gamma\) denote the corresponding
negative root degree.  Suppose the compact Hall source supplies:
\begin{enumerate}
\item retained source charts \(\mathcal U_{\gamma,R,S}\) and
\(\mathcal U_{-\gamma,R,S}\) and a compact-source duality morphism
\[
\mathbb D_{\gamma,R,S}:\mathcal U_{\gamma,R,S}\longrightarrow
\mathcal U_{-\gamma,R,S}
\]
lifting the retained finite-window dual closure;
\item parity decompositions of
\[
K^{\mathrm{red}}_{\gamma,R,S},\qquad
K^{\mathrm{red}}_{-\gamma,R,S}
\]
with equal even ranks and equal odd ranks, identified by the
positive--negative target pairing block for \(\gamma\);
\item a nondegenerate source Hopf-pairing matrix
\[
G_{\gamma,R,S}:
K^{\mathrm{red}}_{\gamma,R,S}\otimes
K^{\mathrm{red}}_{-\gamma,R,S}\longrightarrow \mathbb C
\]
whose off-parity blocks vanish and whose radical is the recorded
source radical;
\item orientation square roots \(o_{\gamma,R,S}\) and
\(o_{-\gamma,R,S}\), with determinant-square diagrams commuting with
the duality isomorphism
\[
\det K^{\mathrm{red}}_{-\gamma,R,S}
\simeq
\bigl(\det K^{\mathrm{red}}_{\gamma,R,S}\bigr)^{-1}
\otimes \varepsilon^{\mathrm{CY}_3}_{\gamma,R,S},
\]
where \(\varepsilon^{\mathrm{CY}_3}_{\gamma,R,S}\) is the supplied
Calabi--Yau threefold Serre-pairing sign line;
\item quotient-orientation null-trivialisations and transition rows
preserving the duality morphism, the pairing matrix, the parity
decompositions, and the two orientation square roots.
\end{enumerate}
Then the reduced orientation line is compatible with the negative root
\(-\gamma\): the negative-root orientation is the Serre-dual
Picard-groupoid transport of the positive-root orientation, and the
positive--negative pairing is even and nondegenerate modulo the recorded
radical.

The present packets do not supply these hypotheses.  The retained
dual-closure packet proves only that the finite HN window is closed
under duals.  The target presentation packet records formal
positive--negative blocks, but every such row has
\(\texttt{source\_pairing=false}\).  The compact Hall source packet has
no \(G\)-entries, Hopf-pairing identities, radical rows, source
negative-root parity blocks, or Chevalley anti-involution rows; the
orientation packets have no square-root row, no quotient-orientation
transport row, and no Picard-groupoid transition row.  Hence row \(333\)
is presently a negative-root compatibility criterion and obstruction
ledger, not a proof of negative-root orientation compatibility.
\end{proposition}

\begin{proof}
The proof is formal once the listed source data are supplied.  The
pairing matrix \(G_{\gamma,R,S}\) identifies the negative source block
with the dual of the positive source block after quotienting by the
recorded radical.  Since its off-parity blocks vanish, the
identification respects the even and odd determinant factors.  The
displayed determinant-duality isomorphism then identifies the
determinant line of the negative block with the dual determinant line of
the positive block, up to the supplied Calabi--Yau Serre sign line.
Commutativity of the square-root diagrams lifts this determinant
identity to the Joyce--Upmeier orientation lines.  The quotient
null-trivialisations make the statement one on the reduced quotient
chart, and the transition rows make the positive--negative
identification a statement in the pro-object rather than at a single
finite height.  No target formal pairing block supplies the source
pairing matrix, radical, or orientation-line transport.
\end{proof}

The packet
\[
\texttt{certificates/orientation/rhomred\_negative\_root\_orientation\_compatibility}
\]
verifies
\(\texttt{RHOMRED\_NEGATIVE\_ROOT\_ORIENTATION\_COMPATIBILITY\_OBSTRUCTION\_VERIFIED}\):
it imports the retained dual-closure packet, the target presentation
packet, the formal Hall-pairing pushforward packet, the compact Hall
source obstruction ledger, the row \(332\) parity-compatibility packet,
and the transition-orientation obstruction ledger; it records the
negative-root orientation-compatibility criterion for row \(333\); and
it keeps the compact source negative-root charts, source
positive--negative pairing matrices, Hopf-pairing identities, radical
quotient rows, source parity blocks, orientation square roots,
Serre-duality sign line, quotient-orientation transport, and
Picard-groupoid transition rows open.

\begin{proposition}[Chevalley anti-involution on orientation data]
\label{prop:orientation-chevalley-antiinvolution}
\textnormal{[anti-involution criterion proved; source maps not supplied]}
Let \(P^{\mathrm{red}}_{R,S}\) be the reduced compact Hall primitive
quotient at finite stage \((R,S)\), with retained positive and negative
degree pieces and Joyce--Upmeier orientation lines
\(o_{\gamma,R,S}\).  Suppose the compact Hall source supplies:
\begin{enumerate}
\item source simple representatives \(e_i\), \(f_i\), and \(h_i\),
including the imaginary simple generators and the recorded parity
blocks, with \(e_i\) and \(f_i\) lying in opposite retained degrees and
\(h_i\) in the Cartan part;
\item an even involutive linear map
\[
\vartheta_{R,S}:P^{\mathrm{red}}_{R,S}\longrightarrow
P^{\mathrm{red}}_{R,S}
\]
which sends \(P^{\mathrm{red}}_{\gamma,R,S}\) to
\(P^{\mathrm{red}}_{-\gamma,R,S}\), fixes the Cartan part, sends
\(e_i\leftrightarrow f_i\), and preserves the source parity
decomposition;
\item source bracket and relation rows showing
\[
\vartheta_{R,S}([x,y])=[\vartheta_{R,S}(y),\vartheta_{R,S}(x)]
\]
for homogeneous primitive classes, together with zero-defect Cartan,
Chevalley, real-Serre, Borcherds-orthogonality, and super-sign
relation matrices after quotienting by the recorded radical;
\item a source positive--negative Hopf pairing for which
\(\vartheta_{R,S}\) is adjoint, the radical is stable under
\(\vartheta_{R,S}\), and the quotient pairing is nondegenerate and
parity preserving;
\item orientation-line isomorphisms in the Picard groupoid
\[
\vartheta^o_{\gamma,R,S}:
o_{\gamma,R,S}\xrightarrow{\sim}
o_{-\gamma,R,S}^{\vee}\otimes
\varepsilon^{\mathrm{CY}_3}_{\gamma,R,S}
\]
whose determinant-square diagrams commute with the source duality map,
the Hopf pairing, and the Calabi--Yau threefold Serre sign line;
\item quotient-orientation null-trivialisations and transition rows
preserving \(\vartheta_{R,S}\), \(\vartheta^o_{\gamma,R,S}\), the
radical quotient, the parity decomposition, and the determinant-square
diagrams.
\end{enumerate}
Then \((\vartheta_{R,S},\vartheta^o_{R,S})\) is a Chevalley
anti-involution on the oriented reduced compact Hall primitive package.
It exchanges the positive and negative simple root data, fixes the
Cartan orientation, is anti-multiplicative for the Lie superbracket, is
adjoint for the Hopf pairing on the radical quotient, and transports
Joyce--Upmeier orientations by the Serre-dual Picard-groupoid
isomorphism.

The present packets do not supply these hypotheses.  The target
presentation contains Chevalley and Serre language only at the formal
target level.  The compact Hall source packet has no simple
representative rows, no bracket \(B\)-entries, no Chevalley relation
rows, no \(G\)-entries, no radical rows, and no parity blocks.  The
negative-root compatibility packet records the dual-orientation
criterion but supplies no source anti-involution map.  The
orientation-transition packet has no Picard-groupoid transition map.
Hence row \(334\) is presently an anti-involution criterion and
obstruction ledger, not a construction of a Chevalley anti-involution on
orientation data.
\end{proposition}

\begin{proof}
All assertions follow from the listed source identities.  The linear map
\(\vartheta_{R,S}\) is involutive and exchanges the positive and
negative primitive source pieces by item \(2\).  Item \(3\) is exactly
the anti-automorphism identity for the even Lie-superalgebra
anti-involution, with the Cartan, Chevalley, Serre, Borcherds, and
super-sign relations checked after passage to the radical quotient.
Item \(4\) makes this anti-involution adjoint for the source Hopf
pairing and prevents a target formal pairing block from replacing the
compact source pairing.  Item \(5\) lifts the determinant duality to the
Joyce--Upmeier orientation lines and records the Calabi--Yau Serre sign
line.  Item \(6\) makes the construction compatible with the quotient
null-trivialisations and with finite-stage transitions.  These are
precisely the data of a Chevalley anti-involution on orientation data.
\end{proof}

The packet
\[
\texttt{certificates/orientation/rhomred\_orientation\_chevalley\_antiinvolution}
\]
verifies
\(\texttt{RHOMRED\_ORIENTATION\_CHEVALLEY\_ANTIINVOLUTION\_OBSTRUCTION\_VERIFIED}\):
it imports the row \(333\) negative-root compatibility packet, the
compact Hall source obstruction ledger, the target presentation packet,
the retained dual-closure packet, the formal Hall-pairing pushforward
packet, and the transition-orientation obstruction ledger; it records
the source Chevalley anti-involution criterion for row \(334\); and it
keeps the source simple representatives, positive--negative source
charts, bracket matrices, Chevalley relation rows, Cartan orientation
rows, Hopf-pairing matrices, radical quotient, parity preservation,
orientation square roots, determinant-duality square, Serre sign line,
quotient-orientation transport, Picard-groupoid transition, involutivity
identity, and anti-multiplicativity identity open.

\begin{proposition}[Frobenius property of the orientation pairing]
\label{prop:orientation-frobenius-pairing}
\textnormal{[Frobenius criterion proved; source cyclicity not supplied]}
Let \(P^{\mathrm{red}}_{R,S}\) be the oriented reduced compact Hall
primitive quotient at finite stage \((R,S)\).  Suppose the compact Hall
source supplies:
\begin{enumerate}
\item Hall product and coproduct matrices \(M_{R,S}\) and \(D_{R,S}\),
unit and counit rows, and zero-defect associativity,
coassociativity, bialgebra-compatibility, and primitive-closure rows;
\item a source Hopf pairing
\[
\langle-,-\rangle^o_{R,S}:
P^{\mathrm{red}}_{R,S}\otimes P^{\mathrm{red}}_{R,S}
\longrightarrow \ell_{R,S}
\]
with orientation-line factors included in the target line
\(\ell_{R,S}\), together with a source trace functional whose degree is
recorded in the next row;
\item zero-defect Hopf-adjointness rows
\[
\langle xy,z\rangle^o_{R,S}
=
\langle x\otimes y,\Delta z\rangle^o_{R,S}
\]
with the Koszul signs and Thom--Sebastiani orientation signs recorded
at the source;
\item zero-defect primitive Frobenius-cyclicity rows
\[
\langle [x,y],z\rangle^o_{R,S}
=
\langle x,[y,z]\rangle^o_{R,S}
\]
for homogeneous primitive classes, with the parity signs and the
Chevalley anti-involution of
Proposition~\ref{prop:orientation-chevalley-antiinvolution} acting by
adjunction;
\item radical-kernel, quotient-splitting, radical ideal/coideal, and
quotient nondegeneracy rows, so that the pairing descends to a
nondegenerate form on the reduced radical quotient;
\item Picard-groupoid orientation pairing squares, Calabi--Yau Serre
sign-line rows, radical-transition rows, pairing-kernel
Mittag--Leffler rows, and orientation-transition rows making the
cyclicity identities independent of the finite HN height.
\end{enumerate}
Then \(\langle-,-\rangle^o_{R,S}\) is a Frobenius pairing on the
oriented reduced compact Hall primitive quotient: it is nondegenerate
after quotienting by the recorded radical, it is adjoint for the Hall
product and coproduct, and its restriction to primitives is cyclic for
the Lie superbracket with the recorded orientation signs.

The present packets do not supply these hypotheses.  The formal
Hall-pairing pushforward packet proves only degree compatibility for a
supplied pairing; it explicitly does not prove Hopf adjointness,
Frobenius cyclicity, or quotient nondegeneracy.  The compact Hall source
packet has no \(M\)-, \(D\)-, \(B\)-, \(G\)-, \(K\)-, or \(Q\)-entries,
no unit or counit rows, no Hall-bialgebra identities, no
Hopf-pairing identity rows, and no radical ideal/coideal rows.  The row
\(334\) packet records the Chevalley anti-involution only as a
criterion.  Hence row \(335\) is presently a Frobenius-pairing
criterion and obstruction ledger, not a proof that the orientation
pairing is Frobenius.
\end{proposition}

\begin{proof}
The argument is the standard finite-dimensional Frobenius argument once
the listed source data are present.  The product, coproduct, unit,
counit, and bialgebra identities make the finite-stage Hall package a
bialgebra on the retained source quotient.  Hopf adjointness identifies
left multiplication with the adjoint of the coproduct.  Restricting to
primitive classes turns this adjointness into the displayed cyclic
identity for the Lie superbracket; the Koszul signs are exactly the
source parity signs, and the orientation signs are the
Thom--Sebastiani and Serre signs recorded in the Picard groupoid.
Radical ideal/coideal stability and quotient nondegeneracy descend the
pairing to the reduced radical quotient.  The transition and
Mittag--Leffler rows make the construction independent of the finite
height.  No scalar trace, target pairing block, or formal degree
compatibility row supplies these source identities.
\end{proof}

The packet
\[
\texttt{certificates/orientation/rhomred\_orientation\_frobenius\_pairing}
\]
verifies
\(\texttt{RHOMRED\_ORIENTATION\_FROBENIUS\_PAIRING\_OBSTRUCTION\_VERIFIED}\):
it imports the row \(334\) Chevalley anti-involution criterion, the
compact Hall source obstruction ledger, the formal Hall-pairing
pushforward packet, the transition-radical obstruction ledger, the
pairing-kernel \(R^1\varprojlim\) obstruction ledger, and the
transition-orientation obstruction ledger; it records the Frobenius
pairing criterion for row \(335\); and it keeps the Hall product and
coproduct matrices, bialgebra identities, source Hopf-pairing matrix,
Hopf-adjointness row, Frobenius-cyclicity row, quotient
nondegeneracy row, radical kernel, quotient splitting, radical
ideal/coideal, orientation pairing line, orientation cyclic signs, Serre
sign line, radical transition rows, pairing-kernel Mittag--Leffler rows,
orientation-transition rows, and row \(336\) trace-degree computation
open.

\begin{proposition}[Degree of the Frobenius trace]
\label{prop:orientation-frobenius-trace-degree}
\textnormal{[trace-degree criterion proved; source trace not supplied]}
Let \(P^{\mathrm{red}}_{R,S}\) be the oriented reduced compact Hall
primitive quotient, and suppose the Frobenius pairing of
Proposition~\ref{prop:orientation-frobenius-pairing} has been supplied
on the radical quotient.  Suppose in addition that the compact Hall
source supplies:
\begin{enumerate}
\item a source trace functional
\[
\operatorname{tr}^o_{R,S}:
P^{\mathrm{red}}_{R,S,\mathrm{rad}}\longrightarrow
\ell^{\mathrm{tr}}_{R,S}[-3]
\]
on the reduced radical quotient, with the target trace line
\(\ell^{\mathrm{tr}}_{R,S}\) recorded in the orientation Picard
groupoid;
\item a source grading convention row for cohomological degree,
parity, orientation-line degree, and Calabi--Yau threefold Serre shift,
with the trace-degree identity
\[
\deg(\operatorname{tr}^o_{R,S})=-3
\]
having defect rank zero;
\item zero-defect trace--pairing rows
\[
\langle x,y\rangle^o_{R,S}
=
\operatorname{tr}^o_{R,S}(xy)
\]
for homogeneous classes \(x,y\), including the Koszul and
Thom--Sebastiani orientation signs;
\item zero-defect trace cyclicity rows
\[
\operatorname{tr}^o_{R,S}(xy)
=(-1)^{|x||y|}\operatorname{tr}^o_{R,S}(yx)
\]
and compatibility with the Chevalley anti-involution and the
Calabi--Yau Serre sign line;
\item radical-quotient, quotient-orientation, transition, and
Mittag--Leffler rows preserving the trace functional, the trace degree,
and the trace--pairing identity in the pro-object.
\end{enumerate}
Then the Frobenius trace has degree \(-3\).  Equivalently, the
orientation pairing has cohomological degree \(-3\):
\[
\deg\langle x,y\rangle^o_{R,S}
=\deg x+\deg y-3
\]
for homogeneous classes on the reduced radical quotient.

The present packets do not supply these hypotheses.  The row \(335\)
packet records the Frobenius pairing only as a criterion and explicitly
keeps the source trace functional open.  The compact Hall source packet
has no source trace rows, no product or pairing matrices, no
Hopf-pairing identities, and no radical quotient rows.  The formal
Hall-pairing pushforward packet proves only homogeneous degree-zero
compatibility for supplied pairing rows.  The protected-integration
Gram-degree packet records the scalar target monomial \(s\)-degree
\(m_R\) of \(I_R^{\mathrm{prot}}\), not the cohomological degree of a
source Frobenius trace.  The level-\(Z\) protected trace packet is an
empty obstruction ledger.  Hence row \(336\) is presently a trace-degree
criterion and obstruction ledger, not a proof that the Frobenius trace
has degree \(-3\).
\end{proposition}

\begin{proof}
The statement follows immediately from the supplied source trace data.
The trace-degree row fixes the target shift of
\(\operatorname{tr}^o_{R,S}\) to \([-3]\).  The trace--pairing identity
then identifies the Frobenius pairing with the composite
\((x,y)\mapsto \operatorname{tr}^o_{R,S}(xy)\), so a homogeneous product
of degree \(\deg x+\deg y\) is sent to the trace line in degree
\(\deg x+\deg y-3\).  Trace cyclicity and the Chevalley/Serre
compatibility rows show that this degree is the Frobenius trace degree,
not only a degree of one chosen product expression.  Radical and
transition rows descend the identity to the reduced pro-object.  A
scalar protected-integration monomial degree does not provide the
source trace functional or the cohomological shift.
\end{proof}

The packet
\[
\texttt{certificates/orientation/rhomred\_orientation\_frobenius\_trace\_degree}
\]
verifies
\(\texttt{RHOMRED\_ORIENTATION\_FROBENIUS\_TRACE\_DEGREE\_OBSTRUCTION\_VERIFIED}\):
it imports the row \(335\) Frobenius-pairing criterion, the compact Hall
source obstruction ledger, the formal Hall-pairing pushforward packet,
the protected-integration Gram-degree packet, the level-\(Z\) protected
trace obstruction ledger, and the transition-orientation obstruction
ledger; it records the trace-degree criterion for row \(336\); and it
keeps the source trace functional, trace-degree row, trace--pairing
identity, trace cyclicity, orientation trace line, Calabi--Yau Serre
shift, radical quotient, quotient-orientation transport, transition
compatibility, and row \(337\) Serre-sign computation open.

\begin{proposition}[Calabi--Yau threefold Serre signs]
\label{prop:orientation-cy3-serre-signs}
\textnormal{[Serre-sign criterion proved; source sign rows not supplied]}
Let \(K^{\mathrm{red}}_{\gamma,R,S}\) and
\(K^{\mathrm{red}}_{-\gamma,R,S}\) be the reduced oriented compact
Hall complexes in opposite retained degrees.  Suppose the compact Hall
source supplies:
\begin{enumerate}
\item a reduced Serre-duality isomorphism
\[
S_{\gamma,R,S}:
K^{\mathrm{red}}_{\gamma,R,S}
\xrightarrow{\sim}
\bigl(K^{\mathrm{red}}_{-\gamma,R,S}\bigr)^{\vee}[-3]
\]
on the radical quotient, compatible with the source Ext-pairing and
with the quotient null-trivialisations;
\item source rows fixing cohomological degree, parity, determinant-line
ordering, and the Koszul sign convention for the shift \([-3]\), with
a sign-exponent row
\(\sigma^{\mathrm{CY}_3}_{\gamma,R,S}\) of defect rank zero;
\item a Picard-groupoid sign line
\(\varepsilon^{\mathrm{CY}_3}_{\gamma,R,S}\) and a commuting
determinant-duality square
\[
\det K^{\mathrm{red}}_{-\gamma,R,S}
\simeq
\bigl(\det K^{\mathrm{red}}_{\gamma,R,S}\bigr)^{-1}
\otimes \varepsilon^{\mathrm{CY}_3}_{\gamma,R,S};
\]
\item zero-defect comparison rows identifying the
Thom--Sebastiani sign, the Chevalley anti-involution sign, the
Frobenius-cyclicity sign, and the trace-cyclicity sign with
\(\varepsilon^{\mathrm{CY}_3}_{\gamma,R,S}\);
\item radical-quotient, quotient-orientation, finite-stage transition,
and Mittag--Leffler rows preserving \(S_{\gamma,R,S}\), the sign
exponent, and the determinant-duality square.
\end{enumerate}
Then the orientation signs agree with Calabi--Yau threefold Serre
duality: every sign used in negative-root transport, Chevalley
adjunction, Frobenius cyclicity, and trace cyclicity is the sign
obtained from the displayed Serre-duality isomorphism and the recorded
source grading convention.

The present packets do not supply these hypotheses.  The PTVV packet
imports the ambient \((-1)\)-shifted symplectic form on retained finite
derived substacks, but it explicitly does not construct the Joyce
d-critical truncation, reduced \(\operatorname{RHom}\), determinant
line, orientation square root, or quotient orientation.  The row
\(333\), \(334\), \(335\), and \(336\) packets all keep the
Calabi--Yau Serre sign line or shift open.  The compact Hall source
packet has no source pairing matrices, Hopf-pairing identities, parity
blocks, radical rows, determinant-duality squares, or Serre sign rows.
The formal Hall-pairing pushforward and target negative-dual blocks are
not source sign computations.  Hence row \(337\) is presently a
Serre-sign criterion and obstruction ledger, not a proof that the
orientation signs agree with Calabi--Yau threefold Serre duality.
\end{proposition}

\begin{proof}
Once the listed source rows are present, the assertion is a determinant
calculation.  The reduced Serre-duality isomorphism identifies the
negative complex with the shifted dual of the positive complex.  The
source degree and parity rows determine the Koszul exponent attached to
the shift \([-3]\); the determinant-ordering row transports this
exponent to the determinant line.  The determinant-duality square is
therefore the Picard-groupoid form of the Calabi--Yau threefold Serre
sign line.  The comparison rows force the signs used in
Thom--Sebastiani gluing, Chevalley adjunction, Frobenius cyclicity, and
trace cyclicity to be this same line.  Radical and transition rows make
the equality one on the reduced pro-object.  An ambient shifted
symplectic form or a target pairing block gives no source determinant
ordering and therefore no Serre-sign computation.
\end{proof}

The packet
\[
\texttt{certificates/orientation/rhomred\_orientation\_cy3\_serre\_signs}
\]
verifies
\(\texttt{RHOMRED\_ORIENTATION\_CY3\_SERRE\_SIGNS\_OBSTRUCTION\_VERIFIED}\):
it imports the row \(333\) negative-root compatibility criterion, the
row \(334\) Chevalley anti-involution criterion, the row \(335\)
Frobenius-pairing criterion, the row \(336\) trace-degree criterion,
the PTVV finite-substack shifted-symplectic packet, the compact Hall
source obstruction ledger, and the transition-orientation obstruction
ledger; it records the Calabi--Yau Serre-sign criterion for row \(337\);
and it keeps the reduced Serre-duality isomorphism, source
sign-exponent table, orientation sign line, determinant-duality square,
Thom--Sebastiani sign comparison, Chevalley sign comparison,
Frobenius-cyclicity sign comparison, trace-cyclicity sign comparison,
radical quotient, quotient orientation, and transition compatibility
open.

\begin{proposition}[Reduced Calabi--Yau pairing modulo the radical]
\label{prop:orientation-cy3-pairing-nondegeneracy}
\textnormal{[nondegeneracy criterion proved; quotient matrices not supplied]}
Let \(P^{\mathrm{red}}_{\gamma,R,S}\) and
\(P^{\mathrm{red}}_{-\gamma,R,S}\) be the retained opposite-degree
primitive source spaces at finite stage \((R,S)\).  Suppose the compact
Hall source supplies:
\begin{enumerate}
\item a Calabi--Yau threefold Hopf-pairing matrix
\[
G^{\mathrm{CY}_3}_{\gamma,R,S}:
P^{\mathrm{red}}_{\gamma,R,S}\otimes
P^{\mathrm{red}}_{-\gamma,R,S}\longrightarrow
\ell^{\mathrm{CY}_3}_{\gamma,R,S}
\]
with the source Serre sign line of
Proposition~\ref{prop:orientation-cy3-serre-signs} included in the
target line;
\item left- and right-kernel rows
\[
K^L_{\gamma,R,S}=\ker G^{\mathrm{CY}_3}_{\gamma,R,S},\qquad
K^R_{\gamma,R,S}=\ker (G^{\mathrm{CY}_3}_{\gamma,R,S})^{t}
\]
and zero-defect rows identifying these kernels with the recorded
Hopf-pairing radical in the positive and negative source spaces;
\item quotient splitting matrices
\[
Q_{\gamma,R,S}:P^{\mathrm{red}}_{\gamma,R,S}\twoheadrightarrow
\overline P_{\gamma,R,S},\qquad
Q_{-\gamma,R,S}:P^{\mathrm{red}}_{-\gamma,R,S}\twoheadrightarrow
\overline P_{-\gamma,R,S}
\]
whose kernels are exactly the recorded radicals;
\item a quotient-pairing matrix
\[
\overline G_{\gamma,R,S}
   = Q_{\gamma,R,S}\,G^{\mathrm{CY}_3}_{\gamma,R,S}\,
     Q_{-\gamma,R,S}^{t}
\]
with full left rank and full right rank, equivalently determinant or
Smith defect rank zero after choosing finite quotient bases;
\item radical ideal/coideal rows, transition-radical rows, and
pairing-kernel Mittag--Leffler rows proving that the quotient pairing
and its nondegeneracy are independent of the finite HN height.
\end{enumerate}
Then the reduced Calabi--Yau threefold pairing is nondegenerate modulo
the radical:
\[
\overline G_{\gamma,R,S}:
\overline P_{\gamma,R,S}\otimes
\overline P_{-\gamma,R,S}
\longrightarrow \ell^{\mathrm{CY}_3}_{\gamma,R,S}
\]
has zero left and right kernels, with the orientation sign line fixed
by Calabi--Yau threefold Serre duality.

The present packets do not supply these hypotheses.  The compact Hall
source packet has no \(G\)-entries, kernel matrices, quotient
splittings, quotient-nondegeneracy identities, radical ideal/coideal
rows, or quotient transition rows.  The formal Hall-pairing pushforward
packet proves only homogeneous degree-zero compatibility for three
supplied formal pairings and explicitly does not prove quotient
nondegeneracy.  The row \(333\), \(335\), and \(337\) packets record
negative-root compatibility, Frobenius pairing, and Serre signs only as
criteria.  The transition-radical and pairing-kernel Mittag--Leffler
packets are obstruction ledgers with empty transition and kernel
tables.  Hence row \(338\) is presently a nondegeneracy-modulo-radical
criterion and obstruction ledger, not a proof that the reduced
Calabi--Yau pairing is nondegenerate modulo the radical.
\end{proposition}

\begin{proof}
The assertion is the elementary quotient-pairing argument once the
listed finite matrices exist.  By the kernel rows, the left and right
radicals of \(G^{\mathrm{CY}_3}_{\gamma,R,S}\) are exactly the kernels
of \(Q_{\gamma,R,S}\) and \(Q_{-\gamma,R,S}\).  Therefore
\(G^{\mathrm{CY}_3}_{\gamma,R,S}\) factors through
\(\overline P_{\gamma,R,S}\otimes\overline P_{-\gamma,R,S}\).  The
full-rank row for
\(Q_{\gamma,R,S}G^{\mathrm{CY}_3}_{\gamma,R,S}Q_{-\gamma,R,S}^{t}\)
states precisely that the induced form has no left or right kernel.
The radical ideal/coideal rows make the quotient compatible with the
Hall algebra operations; the transition and Mittag--Leffler rows make
the finite-stage statement stable in the pro-object.  The Serre-sign
line fixes the orientation of the target line but does not by itself
imply full rank.
\end{proof}

The packet
\[
\texttt{certificates/orientation/rhomred\_orientation\_cy3\_pairing\_nondegeneracy}
\]
verifies
\(\texttt{RHOMRED\_ORIENTATION\_CY3\_PAIRING\_NONDEGENERACY\_OBSTRUCTION\_VERIFIED}\):
it imports the row \(333\) negative-root compatibility criterion, the
row \(335\) Frobenius-pairing criterion, the row \(337\) Serre-sign
criterion, the compact Hall source obstruction ledger, the formal
Hall-pairing pushforward packet, the transition-radical obstruction
ledger, and the pairing-kernel Mittag--Leffler obstruction ledger; it
records the reduced Calabi--Yau pairing nondegeneracy criterion for
row \(338\); and it keeps the source pairing matrix, left and right
kernel rows, radical identification, quotient splitting matrices,
quotient-pairing full-rank row, radical ideal/coideal rows, radical
transition rows, pairing-kernel Mittag--Leffler rows, and inverse-limit
nondegeneracy open.

\begin{definition}[Orientation-preserving transition datum]
\label{def:transition-orientation-preservation}
Let \(R'\le R\), and suppose the finite-moduli transition geometry of
Definition~\ref{def:finite-moduli-transition-geometry} is supplied.
An orientation-preserving transition datum consists of:
\begin{enumerate}
\item a Picard-groupoid isomorphism over the retained transition stack
\[
\theta^o_{RR'}:(t_{RR'})^*o_{R'}\xrightarrow{\sim}o_R;
\]
\item commutativity of the determinant square-root diagram
\[
\bigl((t_{RR'})^*o_{R'}\bigr)^{\otimes2}
\longrightarrow (t_{RR'})^*\det K^{\mathrm{red}}_{R'}
\longrightarrow \det K^{\mathrm{red}}_R,
\]
with the square \(o_R^{\otimes2}\simeq\det K^{\mathrm{red}}_R\);
\item transport of the null-trivialisations of
\(\alpha^{\mathrm{red}}\), \(\alpha^{E,\mathrm{free}}\),
\(\widetilde\beta^H\), \(\beta^H\), and
\(\lambda^H=0\) for every retained finite stabilizer stratum;
\item compatibility of \(\theta^o_{RR'}\) with the
Thom--Sebastiani orientation isomorphisms on extension and two-step
flag stacks;
\item compatibility with the Weyl lifts \(\tau_i\), the cochain
killing the projective Coxeter cocycle, and the quotient-cocycle
transport rows of \((O1)^+\);
\item zero Picard-groupoid, square-root, null-trivialization,
Thom--Sebastiani, Weyl, and composition defect ranks, together with
\(R^1\varprojlim=0\) for the orientation Picard data.
\end{enumerate}
Proper or closed stack transitions alone do not define this datum.
Scalar traces, squared determinants, OP scalar branches, Maass-character
values, Pfaffian products, and target root windows are not orientation
transport data.
\end{definition}

\begin{proposition}[Transition preservation of Joyce--Upmeier
orientations; \textit{proved}]
\label{prop:transition-orientation-preservation-criterion}
Given an orientation-preserving transition datum of
Definition~\ref{def:transition-orientation-preservation}, the
transition \(R\to R'\) preserves the reduced Joyce--Upmeier
orientation line, the quotient-orientation null-trivialisations, the
Thom--Sebastiani multiplicative structure, and the \((O1)^+\)
Weyl-equivariant determinant-line transport.
\end{proposition}

\begin{proof}
The isomorphism \(\theta^o_{RR'}\) identifies the pulled-back target
orientation line with the source orientation line in the Picard
groupoid.  The square-root diagram in item \(2\) shows that this
identification is an identification of Joyce--Upmeier orientations, not
only of underlying line bundles.  Item \(3\) transports the quotient
Cech--Borel null-trivialisations and the zero finite-stabilizer
linearization characters.  Item \(4\) makes the Hall
Thom--Sebastiani orientation isomorphisms commute with transition, and
item \(5\) makes the type-II Weyl lifts and Coxeter cochain commute
with transition.  The zero-defect and \(R^1\varprojlim\) rows in item
\(6\) give strict composition and inverse-limit exactness for the
orientation Picard data.
\end{proof}

The packet
\(\texttt{certificates/orientation/transition\_orientation\_preservation}\)
records the present state of this datum.  Its verifier status is
\texttt{TRANSITION\_ORIENTATION\_PRESERVATION\_OBSTRUCTION\_VERIFIED}:
the stack-transition input, orientation-line pullback isomorphisms,
determinant square compatibility, quotient Cech--Borel
null-trivialization transport rows, finite-stabilizer and
linearization transport rows, Thom--Sebastiani transition rows,
Weyl-lift transition rows, Coxeter-cochain transition rows, strict
Picard-groupoid transition rows, and orientation Mittag--Leffler rows
are not yet supplied.  Hence transition preservation of orientations
remains open.

\begin{proposition}[Survival of orientation classes in the inverse limit]
\label{prop:orientation-class-inverse-limit-survival}
\textnormal{[inverse-limit survival criterion proved; orientation tower not supplied]}
Let \(\mathcal R\) be the retained cofinal HN index category, and for
each \(R\in\mathcal R\) let \(o_R\) be a Joyce--Upmeier orientation
line on the reduced compact Hall frame at stage \(R\).  Suppose the
compact orientation source supplies:
\begin{enumerate}
\item finite-stage orientation rows
\[
o_R^{\otimes 2}\simeq \det K^{\mathrm{red}}_R,\qquad
\operatorname{cl}(o_R)=0
\]
with zero square-root and orientation-class defect ranks for every
retained stratum;
\item for every successor \(R'\le R\), Picard-groupoid transition
isomorphisms
\[
\theta^o_{RR'}:(t_{RR'})^*o_{R'}\xrightarrow{\sim}o_R
\]
whose determinant-square, quotient-orientation, Thom--Sebastiani, Weyl,
and Coxeter transition defects are zero;
\item strict identity and composition rows for the transitions
\(\theta^o_{RR'}\), so that the \(o_R\) form an inverse system in the
orientation Picard groupoid;
\item an orientation Mittag--Leffler row with
\[
R^1\varprojlim_R\,\pi_0\operatorname{Pic}^{or}_R=0
\]
and compatible image-stabilization witnesses for the orientation
classes;
\item a limit row identifying the pro-orientation class
\[
o_\infty=\varprojlim_R(o_R,\theta^o_{RR'})
\]
with the class used by the Dirac--Igusa pro-object.
\end{enumerate}
Then the finite-stage orientation classes survive the inverse limit:
the compatible system \((o_R,\theta^o_{RR'})\) determines an
orientation class \(o_\infty\) on the pro-object, and every finite
projection of \(o_\infty\) is the prescribed \(o_R\).

The present packets do not supply these hypotheses.  The reduced
orientation ledger has no orientation-line rows, no square-root rows,
no orientation-class rows, and no transition rows.  The square-root
packet records three missing square-root rows and has no orientation
class zero row.  The transition-orientation packet has no
Picard-groupoid pullback maps, no determinant-square transition rows,
no strict composition rows, and no orientation Mittag--Leffler row.
The row \(338\) nondegeneracy packet keeps inverse-limit
nondegeneracy open and does not construct orientation classes.  Hence
row \(339\) is presently an inverse-limit survival criterion and
obstruction ledger, not a proof that orientation classes survive the
inverse limit.
\end{proposition}

\begin{proof}
The claim is the standard exactness statement for an inverse system in
Picard groupoids once the listed rows are present.  The finite-stage
rows give objects \(o_R\) with the prescribed determinant square and
vanishing orientation obstruction.  The transition rows make these
objects compatible under pullback.  Strict identity and composition
rows replace weak pairwise compatibility by an inverse system.  The
Mittag--Leffler row kills the \(R^1\varprojlim\) obstruction to lifting
the compatible finite classes to a pro-class.  The limit row then
identifies the resulting object with the orientation datum used in the
Dirac--Igusa pro-object.  This proves survival of the prescribed
classes.  It does not prove that every orientation class at the limit
comes from a finite-stage class; that is the separate row \(340\)
obligation.
\end{proof}

The packet
\[
\texttt{certificates/orientation/rhomred\_orientation\_class\_limit\_survival}
\]
verifies
\(\texttt{RHOMRED\_ORIENTATION\_CLASS\_LIMIT\_SURVIVAL\_OBSTRUCTION\_VERIFIED}\):
it imports the reduced-orientation obstruction ledger, the square-root
obstruction packet, the transition-orientation obstruction ledger, and
the row \(338\) nondegeneracy criterion; it records the inverse-limit
survival criterion for row \(339\); and it keeps the finite orientation
line rows, square-root rows, orientation-class zero rows, Picard
transition maps, strict composition rows, orientation
Mittag--Leffler rows, limit class row, and row \(340\) no-new-class
statement open.

\begin{proposition}[No new orientation class at the limit]
\label{prop:orientation-no-new-limit-class}
\textnormal{[no-phantom criterion proved; pro-Picard comparison not supplied]}
Let \(\operatorname{Pic}^{or}_\infty\) be the orientation Picard
groupoid of the Dirac--Igusa pro-object, and let
\(\operatorname{Pic}^{or}_R\) be the finite-stage orientation Picard
groupoid at HN height \(R\).  Suppose the compact orientation source
supplies:
\begin{enumerate}
\item a pro-Picard comparison functor
\[
\Xi:\operatorname{Pic}^{or}_\infty
\longrightarrow
\varprojlim_R\operatorname{Pic}^{or}_R
\]
whose object and automorphism components are defined by the finite
stack transitions;
\item finite-stage effectivity rows showing that every orientation line
on the pro-object restricts to the recorded finite-stage orientation
line, square-root, quotient-orientation, Thom--Sebastiani, Weyl, and
Coxeter data;
\item separatedness rows proving that a pro-orientation class whose
finite restrictions are all trivial is itself trivial;
\item automorphism Mittag--Leffler rows with
\[
R^1\varprojlim_R\,\pi_1\operatorname{Pic}^{or}_R=0
\]
and object-class Mittag--Leffler rows with zero \(R^1\varprojlim\)
defect for the relevant orientation-class tower;
\item a zero-defect comparison row
\[
\pi_0\operatorname{Pic}^{or}_\infty
\xrightarrow{\sim}
\varprojlim_R\pi_0\operatorname{Pic}^{or}_R
\]
on the subfunctor cut out by the determinant square, quotient
orientation, Thom--Sebastiani, Weyl, and Coxeter conditions.
\end{enumerate}
Then no new orientation class appears at the limit: every
Dirac--Igusa pro-orientation class is detected by its compatible
finite-stage restrictions, and the class is the image of the finite
tower under Proposition~\ref{prop:orientation-class-inverse-limit-survival}.

The present packets do not supply these hypotheses.  The row \(339\)
packet records only the survival criterion and has no limit-survival
rows.  The reduced-orientation ledger has no finite orientation-line or
transition rows.  The transition-orientation packet has no
Picard-groupoid transition maps, no strict composition rows, no
orientation Mittag--Leffler rows, and no pro-Picard comparison rows.
The square-root packet records missing finite square roots and no
orientation-class zero rows.  Hence row \(340\) is presently a
no-phantom-class criterion and obstruction ledger, not a proof that no
new orientation class appears at the limit.
\end{proposition}

\begin{proof}
Assume the listed comparison and exactness rows.  The functor \(\Xi\)
restricts a pro-orientation class to its finite-stage orientation
tower.  The effectivity rows identify this tower with the same
finite-stage orientation data used in the Dirac--Igusa construction.
Separatedness says that the kernel of \(\pi_0(\Xi)\) is zero: a class
with trivial finite restrictions is trivial at the pro-level.
Vanishing of \(R^1\varprojlim\) on automorphisms and object classes
removes the usual derived-limit ambiguity in passing between Picard
groupoids and their towers.  The displayed zero-defect comparison row
therefore identifies pro-orientation classes with compatible finite
orientation classes.  By
Proposition~\ref{prop:orientation-class-inverse-limit-survival}, each
compatible finite tower gives the prescribed limit class.  No additional
class can occur at the limit.
\end{proof}

The packet
\[
\texttt{certificates/orientation/rhomred\_orientation\_no\_new\_limit\_class}
\]
verifies
\(\texttt{RHOMRED\_ORIENTATION\_NO\_NEW\_LIMIT\_CLASS\_OBSTRUCTION\_VERIFIED}\):
it imports the row \(339\) inverse-limit survival criterion, the
reduced-orientation obstruction ledger, the square-root obstruction
packet, and the transition-orientation obstruction ledger; it records
the no-new-limit-class criterion for row \(340\); and it keeps the
pro-Picard comparison functor, finite-stage effectivity rows,
separatedness rows, automorphism Mittag--Leffler rows, object-class
Mittag--Leffler rows, zero-defect comparison row, and pro-orientation
class table open.

\begin{definition}[Vanishing-cycle-preserving transition datum]
\label{def:transition-vanishing-cycle-preservation}
Let \(R'\le R\).  Suppose the finite-moduli transition geometry of
Definition~\ref{def:finite-moduli-transition-geometry} and the
orientation-preserving transition datum of
Definition~\ref{def:transition-orientation-preservation} are supplied.
A vanishing-cycle-preserving transition datum consists of:
\begin{enumerate}
\item transition maps of oriented d-critical charts, with transition
graph \(\mathfrak G_{RR'}\) and projections
\[
p_R:\mathfrak G_{RR'}\to\mathfrak Z_R,\qquad
p_{R'}:\mathfrak G_{RR'}\to\mathfrak Z_{R'};
\]
\item potential compatibility on charts:
\[
f_R=p_{R'}^*f_{R'}+q_{RR'}
\]
after pullback to the transition graph, where \(q_{RR'}\) is the
recorded non-degenerate quadratic stabilization and its parity shift is
recorded;
\item pullback compatibility of the K3 semiregularity cosections and
a morphism of reduced obstruction complexes with zero cosection and
reduced-complex defect ranks;
\item a constructible-complex isomorphism
\[
\theta^\Phi_{RR'}:
p_R^*\Phi^{\mathrm{red}}_R
\xrightarrow{\sim}
p_{R'}^*\Phi^{\mathrm{red}}_{R'}
\]
on \(\mathfrak G_{RR'}\), with BBDJS perverse normalization and
support defects zero;
\item compatibility of \(\theta^\Phi_{RR'}\) with the
Thom--Sebastiani vanishing-cycle isomorphisms on extension and
two-step flag stacks and with proper base change on retained support;
\item zero d-critical, potential, cosection, vanishing-cycle
transport, Thom--Sebastiani, base-change, and composition defect
ranks, together with \(R^1\varprojlim=0\) for the reduced
vanishing-cycle complexes.
\end{enumerate}
Scalar traces, Euler characteristics, Behrend-weighted numbers,
Pfaffian products, orientation lines alone, target root windows, D0
scalar tests, and signed multiplicities are not vanishing-cycle
transition data.
\end{definition}

\begin{proposition}[Transition preservation of reduced vanishing
cycles; \textit{proved}]
\label{prop:transition-vanishing-cycle-preservation-criterion}
Given a vanishing-cycle-preserving transition datum of
Definition~\ref{def:transition-vanishing-cycle-preservation}, the
transition \(R\to R'\) preserves the reduced BBDJS perverse sheaves of
vanishing cycles, their cosection reduction, their
Thom--Sebastiani products, and their compact-support base-change
domains.
\end{proposition}

\begin{proof}
The d-critical chart transition and potential compatibility identify
the local vanishing-cycle functors after the recorded quadratic
stabilization.  The cosection compatibility identifies the reduced
obstruction theories, so the reduced vanishing-cycle complexes lie on
the same reduced critical locus after pullback to
\(\mathfrak G_{RR'}\).  The isomorphism
\(\theta^\Phi_{RR'}\) identifies the two BBDJS complexes in
\(D^b_c(\mathfrak G_{RR'})\), and the perverse-normalization row
records the shift introduced by stabilization.  The
Thom--Sebastiani and proper-base-change rows make this identification
compatible with extension and two-step flag correspondences.  The
zero-defect and \(R^1\varprojlim\) rows give strict composition and
inverse-limit exactness for the reduced vanishing-cycle system.
\end{proof}

The packet
\(\texttt{certificates/vanishing\_cycles/transition\_vanishing\_cycle\_preservation}\)
records the present state of this datum.  Its verifier status is
\texttt{TRANSITION\_VANISHING\_CYCLE\_PRESERVATION\_OBSTRUCTION\_VERIFIED}:
the stack-transition input, orientation-transition input, d-critical
chart transition rows, potential-compatibility rows, quadratic
stabilization rows, cosection-pullback rows, reduced obstruction
complex maps, graph-pullback vanishing-cycle isomorphisms, perverse
normalization rows, support and base-change rows,
Thom--Sebastiani vanishing-cycle compatibility rows, strict
composition rows, and vanishing-cycle Mittag--Leffler rows are not yet
supplied.  Hence transition preservation of reduced vanishing cycles
remains open.

\subsubsection*{Cosection localization and the cosection-reduced
heart}

After the surjectivity, kernel, and perfectness rows have been
supplied, the K3 semiregularity cosection
\(\operatorname{CS}_{\mathcal C}\) selects the reduced obstruction
theory.  Only then does the corresponding Behrend constructible
function \(\nu^{\mathrm{red}}\) and the reduced vanishing-cycle complex
\(\Phi^{\mathrm{red}}\) enter the Hall product
\cite{BEHREND,JoyceSong}.  The filtered Hall additivity
\(q^*\operatorname{CS}_C=p_1^*\operatorname{CS}_A+p_2^*\operatorname{CS}_B\)
on extension stacks is the later cosection-localization compatibility
required by the Hall correspondence; the Joyce--Upmeier orientation
lines are required to be Thom--Sebastiani-coherent on the resulting
two-step and four-input flag stacks.

\begin{proposition}[Finite d-critical Hall product;
\textit{conditional on the (O1) and atlas data}]
\label{prop:finite-dcritical-hall-product}
Fix $R$ and closed-configuration substacks
$\mathfrak M_{c,R,I}^{\mathrm{loc,cl}}$ in a noetherian abelian Hall
heart, and assume:
\begin{enumerate}
\item finite-type quasi-smooth derived Artin enhancements of the
selected object, extension, and two-step flag stacks, with perfect
obstruction theories;
\item finite residual inertia after scalar, $E$-translation, and
wrapped rigidifications;
\item retained closed flag-Quot/filtration extension and flag stacks,
or an explicitly specified compact-support six-functor model, with
admissible $p^*$ and $q_!$;
\item PTVV $(-1)$-shifted symplectic structures restricting to these
substacks, classical truncations carrying Joyce d-critical
structures \cite{PTVV2013,BBDJS2015};
\item surjective K3 semiregularity cosections with filtered Hall
additivity;
\item Joyce--Upmeier orientation lines and orientation transports
\cite{JoyceUpmeier} with all finite-stabilizer characters
trivialized;
\item Thom--Sebastiani transport for the reduced vanishing-cycle
sheaves satisfying two-step flag pentagons and higher-coloured
configuration coherences.
\end{enumerate}
Then the finite reduced local-local product
\[
m_R^{\mathrm{red}}=q_!\circ\operatorname{TS}^{\mathrm{red}}_R\circ p^*
\]
is an honest morphism on
$H_c^{\bullet}(\mathfrak M_{c,R,I}^{\mathrm{loc,cl}},
\Phi_{c,R,I}^{\mathrm{red}}\otimes o_{c,R,I}^{\mathrm{red}})$
and is associative on the finite stage.
\end{proposition}

The first obstruction to removing these hypotheses is finite inertia:
proper finite-type correspondences do not by themselves eliminate
residual $\mathbb G_m$, direct-sum, or parabolic automorphisms, and a
residual $\mathbb G_m$ already contributes
$H^\bullet(B\mathbb G_m)=\C[u]$ to protected cohomology.  The phrase
\emph{finite stage} here means finite charge support, finite-type
stacks, and finite-dimensional protected cohomology after the stated
quotients and admissibility hypotheses; it does not mean that the
Hall correspondence morphisms are themselves finite morphisms.
The raw exact-triangle stack is not a proper Hall correspondence; the
proper object is the retained compactified filtration stack, or else
the compact-support functor \(q_!\) is an extra datum.

\subsubsection*{Quotient-after-correspondence ordering}

The $E$-quotient is taken \emph{after} the extension correspondences,
not before.  A pseudofunctor
$\mathsf{Corr}^{E,\mathrm{hyb}}_R\to
\mathsf{Corr}^{\mathrm{red},\mathrm{hyb}}_R$ supplies the descent;
applying the vanishing-cycle Borel--Moore functor with orientation
lines included gives $Q_{E,R}$ between the corresponding finite
correspondence module categories.  For each operation
$\mu_R^\bullet=q_!\operatorname{TS}p^*$ a chosen comparison
isomorphism
\[
\theta^Q_{\mu,R}:Q_{E,R}\,q_!\operatorname{TS}p^*
\xrightarrow{\sim}
\overline q_!\overline{\operatorname{TS}}\,\overline p^*
(Q_{E,R}^{\boxtimes k})
\]
is part of the datum.  Its compatibility with the eight-word flag
$2$-morphisms and HN transitions is the (O1) descent calculus.  A
quotient-first construction (object-stack quotient followed by a new
Hall product) is not the same datum.

\subsubsection*{The orientation obstruction class}

The orientation lane carries the residual
\[
\mathfrak O_{\mathrm{or},R}
=\bigl(
\mathfrak o^{\mathrm{red}}_{\mathcal C},
\mathfrak o^{E,\mathrm{free}}_{\mathcal C},
\{\mathfrak o^{H}_{\mathcal C,S}\}_{H,S},
\{\mathfrak o^{\lambda^H}_{\mathcal C,S}\}_{H,S},
[c_o]_R,\,\xi_{c_o,R},\,
\{\mathfrak o^{\tau,2}_{i,R}\}_i,\,
\{\omega_{i,\mathcal C,R}\}_i\bigr)
\]
required to vanish at every finite HN height, with chosen
null-trivializations.  Theorems~\ref{thm:G-O1}
and~\ref{thm:G-O1-plus} of
Appendix~\ref{app:obstruction-discharges} are the quotient-orientation
and Weyl-transport criteria relative to that retained datum.  The
packet
\[
\texttt{certificates/orientation/k3e\_reduced\_orientation/blocked\_obligations.csv}
\]
records the missing finite equations: determinant square-root and
orientation-class rows, Thom--Sebastiani and pentagon rows, the four
quotient Cech--Borel class rows, finite-stabilizer edge reduction,
odd-transfer, Klein-four and two-primary coefficient rows, Weyl lift
involutivity, torsor-defect, quotient-cocycle transport,
projective-cocycle, cochain, Coxeter, transition, Mittag--Leffler, and
scalar-firewall rows.  Its verifier returns
\(\texttt{ORIENTATION\_OBSTRUCTION\_LEDGER\_VERIFIED}\), with
\(\texttt{orientation\_certification=false}\) and
\(\texttt{mathematical\_certification=false}\).  It is the finite list
of equations still required for \((O1)\) and \((O1)^+\), not a proof of
those equations.
The
Klein-four and two-primary finite-stabilizer lemmas there give the
edge coordinates of (O1).  Proposition~\ref{prop:finite-stabilizer-orientation-detection}
is the finite-site test: the Borel spectral sequence must first reduce
the equivariant class to the classifying-space edge, and for
two-primary rank two stabilizers the mixed coefficient must be computed
directly.  The wrapped-stratum Weyl-equivariant transport along
\(W^{(2)}(\La^{2,1}_{II})\) is reduced to the no-braid type-II Coxeter
presentation and the Maass-character match.

\subsubsection*{Comparison with the Calabi--Yau quantum-group stratum}

Vol III's Calabi--Yau quantum-groups stratification observes that
$K3\times E$ is not presented by a global toric/quiver-variety NCCR
atlas of the Schiffmann--Vasserot type, so the orientation lane
cannot be supplied by such an NCCR-induced heart; the
cosection-twisted reduced d-critical structure is forced on the
chosen stability heart instead.  The displayed Borel
calculation at $E[2]$ is the smallest finite stabilizer detected by
the cosection-reduced determinant; even $N\ge4$ stabilizers carry the
mixed coefficient $A_{12}^{(N)}$, which the displayed cyclic
restrictions miss, and which (O1) requires to be null-trivialized as
part of the full degree-two gerbe.

The orientation $o_R$ trivializes the Joyce--Upmeier obstruction to
a compatible Hall product on the cosection-twisted heart.  Once
$o_R$ is in place, the Hall product, the collision coproduct, and
the protected primitive extraction are unambiguous.  Entry
$(\mathrm L 3)$ is the Hall apparatus of
\S\ref{subsec:hall-correspondences}.
